% ==========================================================
%  MASTER SUMMARY TEMPLATE — AI / ML / MATHEMATICS
%  Design Philosophy:
%    1. Readable running text is the default. Prose over boxes.
%    2. Boxes are used sparingly as visual landmarks for skimming.
%    3. Visual hierarchy: chapter > section > subsection > prose.
%    4. Every box is breakable — no orphaned whitespace at page ends.
%    5. A single consistent accent color ties all structure together.
% ==========================================================

\documentclass[11pt, oneside]{book}


% ══════════════════════════════════════════════════════════
%  PAGE GEOMETRY
% ══════════════════════════════════════════════════════════
\usepackage[
  a4paper,
  top    = 2.6cm,
  bottom = 2.1cm,
  left   = 2.5cm,
  right  = 2.5cm,
  headheight = 15pt
]{geometry}


% ══════════════════════════════════════════════════════════
%  TYPOGRAPHY
% ══════════════════════════════════════════════════════════
\usepackage[T1]{fontenc}
\usepackage[utf8]{inputenc}
\usepackage{lmodern}         % Latin Modern: improved Computer Modern
\usepackage{microtype}       % Microtypographic refinements (better justification,
                             %   fewer overfull hboxes — invaluable for math text)

% If inconsolata is not installed, simply remove this line.
% It provides a crisper monospace font for code listings.
\usepackage{inconsolata}

\setlength{\parindent}{0pt}
\setlength{\parskip}{0.62em}


% ══════════════════════════════════════════════════════════
%  MATHEMATICS
% ══════════════════════════════════════════════════════════
\usepackage{amsmath, amssymb, amsthm, mathtools}
\usepackage{bm}        % Bold math symbols: \bm{\theta}
\usepackage{bbm}       % Blackboard bold for indicator: \mathbbm{1}
\usepackage{mathrsfs}  % Script letters: \mathscr{F}, \mathscr{H}
\usepackage{nicefrac}  % Compact inline fractions: \nicefrac{1}{2}
\usepackage{cancel}    % Strike-through in math: \cancel{x}


% ══════════════════════════════════════════════════════════
%  SCIENTIFIC NOTATION
% ══════════════════════════════════════════════════════════
% siunitx — consistent typesetting of numbers and SI units.
%   Usage: \SI{1.5}{\micro\meter}  \SI{50}{\kilo\hertz}  \num{6.022e23}
\usepackage{siunitx}
\sisetup{per-mode = symbol}   % render "per unit" as e.g. m/s rather than m s^{-1}

% mhchem — chemical formulas and reaction equations.
%   Usage: \ce{H2O}   \ce{ATP -> ADP + P_i}   \ce{^{14}C}
\usepackage[version=4]{mhchem}


% ══════════════════════════════════════════════════════════
%  FIGURES, TABLES, LISTS
% ══════════════════════════════════════════════════════════
\usepackage{graphicx}
\usepackage{booktabs}         % \toprule, \midrule, \bottomrule
\usepackage{array}            % Extended column types in tabular
\usepackage{float}            % [H] float placement
\usepackage{enumitem}

\setlist[itemize]{topsep=3pt, itemsep=2pt, leftmargin=2.2em}
\setlist[enumerate]{topsep=3pt, itemsep=2pt, leftmargin=2.4em}

\usepackage{needspace}        % keep headings together with their first lines

% ── Course addition: self-drawn figures (TikZ / pgfplots) ──
\usepackage{tikz}
\usetikzlibrary{arrows.meta, positioning, shapes.geometric, calc, fit,
                backgrounds, decorations.pathreplacing}
\usepackage{pgfplots}
\pgfplotsset{compat=1.18}
\usepgfplotslibrary{groupplots, fillbetween}

% Float pages: stack floats from the top instead of spreading them out.
\makeatletter
\setlength{\@fptop}{0pt}
\setlength{\@fpsep}{14pt plus 4pt}
\setlength{\@fpbot}{0pt plus 1fil}
\makeatother


% ══════════════════════════════════════════════════════════
%  ALGORITHMS / PSEUDOCODE
% ══════════════════════════════════════════════════════════
% Remove this block if you do not need pseudocode.
\usepackage[ruled, vlined, linesnumbered]{algorithm2e}
\SetAlFnt{\small}
\SetAlCapFnt{\small\bfseries}
\SetCommentSty{textit}


% ══════════════════════════════════════════════════════════
%  COLORS  (a refined, academic palette)
% ══════════════════════════════════════════════════════════
\usepackage{xcolor}

% Primary structural accent — deep navy used for headings, rules, ToC
\definecolor{accent}{HTML}{1B3A5C}

% Box accent colors: frame and background, designed to be legible but
% unobtrusive so running text remains the visual focus.
\definecolor{def@frame}{HTML}{1A5276}   \definecolor{def@bg}{HTML}{EBF5FB}
\definecolor{thm@frame}{HTML}{922B21}   \definecolor{thm@bg}{HTML}{FDEDEC}
\definecolor{ex@frame} {HTML}{1E8449}   \definecolor{ex@bg} {HTML}{EAFAF1}
\definecolor{int@frame}{HTML}{9A6A00}   \definecolor{int@bg}{HTML}{FEF9E7}
\definecolor{note@frame}{HTML}{5B2C6F}  \definecolor{note@bg}{HTML}{F5EEF8}


% ══════════════════════════════════════════════════════════
%  CODE LISTINGS
% ══════════════════════════════════════════════════════════
\usepackage{listings}
\lstset{
  basicstyle       = \ttfamily\small,
  breaklines       = true,
  numbers          = left,
  numberstyle      = \tiny\color{gray!70},
  frame            = tb,
  framerule        = 0.4pt,
  rulecolor        = \color{gray!35},
  tabsize          = 2,
  showstringspaces = false,
  keywordstyle     = \color{def@frame}\bfseries,
  commentstyle     = \color{gray!55!black}\itshape,
  stringstyle      = \color{ex@frame!90!black},
  backgroundcolor  = \color{gray!5},
  xleftmargin      = 2em,
  xrightmargin     = 0.4em,
  aboveskip        = 0.8em,
  belowskip        = 0.5em
}


% ══════════════════════════════════════════════════════════
%  CUSTOM BOXES  (tcolorbox)
%
%  KEY DESIGN DECISIONS:
%  1. enhanced   — enables advanced tcolorbox features (borderline etc.)
%  2. breakable  — box may split across a page break; eliminates the ugly
%                  white gap that appears when a long box would overflow.
%  3. Left-accent bar style — a thick colored stripe on the left margin.
%                  This is legible, modern, and unobtrusive relative to
%                  a full colored frame. The box background is a very
%                  light tint, keeping the reading experience calm.
%  4. sharp corners=west — left side is flat (squared), right is rounded.
%                  This mirrors the visual logic of the left accent bar.
%  5. parbox=false — allows display math to stretch properly inside boxes.
% ══════════════════════════════════════════════════════════
\usepackage[most]{tcolorbox}

\tcbset{
  enhanced,
  breakable,
  arc             = 2.5mm,
  boxrule         = 0.5pt,
  left            = 4.5mm,
  right           = 4mm,
  top             = 2.5mm,
  bottom          = 2.5mm,
  fonttitle       = \bfseries\small,
  sharp corners   = west,
  parbox          = false
}

\newtcolorbox{defbox}[1][]{
  title  = {Definition},
  colback = def@bg, colframe = def@frame,
  borderline west = {3.5pt}{0pt}{def@frame},
  #1
}

\newtcolorbox{thmbox}[1][]{
  title  = {Theorem / Key Formula},
  colback = thm@bg, colframe = thm@frame,
  borderline west = {3.5pt}{0pt}{thm@frame},
  #1
}

\newtcolorbox{exbox}[1][]{
  title  = {Example},
  colback = ex@bg, colframe = ex@frame,
  borderline west = {3.5pt}{0pt}{ex@frame},
  #1
}

\newtcolorbox{intbox}[1][]{
  title  = {Intuition \& Mental Model},
  colback = int@bg, colframe = int@frame,
  borderline west = {3.5pt}{0pt}{int@frame},
  #1
}

\newtcolorbox{notebox}[1][]{
  title  = {Note},
  colback = note@bg, colframe = note@frame,
  borderline west = {3.5pt}{0pt}{note@frame},
  #1
}


% ══════════════════════════════════════════════════════════
%  TABLE OF CONTENTS  (tocloft)
%
%  The [titles] option prevents tocloft from conflicting
%  with titlesec's redefinition of chapter/section titles.
% ══════════════════════════════════════════════════════════
\usepackage[titles]{tocloft}

% Chapter entries: accented, bold, with dots leading to page number
\renewcommand{\cftchapfont}     {\bfseries\color{accent}}
\renewcommand{\cftchappagefont} {\bfseries\color{accent}}
\renewcommand{\cftchapleader}   {\bfseries\color{accent!40}\cftdotfill{\cftdotsep}}
\setlength{\cftbeforechapskip} {7pt}

% Section entries: standard weight, small size
\renewcommand{\cftsecfont}      {\small}
\renewcommand{\cftsecpagefont}  {\small}
\setlength{\cftsecindent}      {1.8em}
\setlength{\cftbeforesecskip}  {1.5pt}

% Subsection entries: muted gray to visually recede
\renewcommand{\cftsubsecfont}      {\small\color{gray!65!black}}
\renewcommand{\cftsubsecpagefont}  {\small\color{gray!65!black}}
\setlength{\cftsubsecindent}      {3.4em}
\setlength{\cftbeforesubsecskip}  {0.5pt}


% ══════════════════════════════════════════════════════════
%  CHAPTER & SECTION HEADINGS  (titlesec)
%
%  Chapter: large faded chapter number floated right, title
%           left-aligned beneath it, then a double rule.
%           The oversized ghost number is a quiet nod to
%           modern textbook design without being gimmicky.
%
%  Section: accent-colored, with a thin rule beneath.
%           The rule draws the eye without shouting.
%
%  Subsection: slightly smaller accent color, no rule — the
%           rule at section level is already enough structure.
% ══════════════════════════════════════════════════════════
\usepackage{titlesec}

\titleformat{\chapter}[display]
  {\normalfont\bfseries\color{accent}\huge}
  {\hfill\fontsize{68}{68}\selectfont\color{accent!13}\thechapter}
  {-10pt}
  {}
  [{\vspace{4pt}%
    {\color{accent}\rule{\linewidth}{1.6pt}}\\[-2pt]%
    {\color{accent!35}\rule{\linewidth}{0.5pt}}}]
\titlespacing*{\chapter}{0pt}{-28pt}{22pt}

\titleformat{\section}
  {\normalfont\large\bfseries\color{accent}}
  {\thesection}
  {0.8em}
  {}
  [{\vspace{1pt}\color{accent!45}\titlerule[0.45pt]}]
\titlespacing*{\section}{0pt}{3.8ex plus 1ex minus .2ex}{2.2ex plus .2ex}

\titleformat{\subsection}
  {\normalfont\normalsize\bfseries\color{accent!80!black}}
  {\thesubsection}
  {0.7em}
  {}
\titlespacing*{\subsection}{0pt}{2.6ex plus 1ex minus .2ex}{1.4ex plus .2ex}

% Run-in paragraph label — useful for named sub-items in a derivation.
\titleformat{\paragraph}[runin]
  {\normalfont\small\bfseries}
  {}
  {0em}
  {}[.\enspace]
\titlespacing{\paragraph}{0pt}{1.2ex plus .5ex minus .2ex}{0.5em}


% ══════════════════════════════════════════════════════════
%  HEADER & FOOTER
% ══════════════════════════════════════════════════════════
\usepackage{fancyhdr}
\pagestyle{fancy}
\fancyhf{}
\fancyhead[L]{\small\color{gray!70}\nouppercase{\leftmark}}
\fancyhead[R]{\small\color{gray!70}\thepage}
\renewcommand{\headrulewidth}{0.4pt}

% Chapter-opening pages use LaTeX's 'plain' style — redefine it here
% to match the rest of the document (page number at foot, no header line).
\fancypagestyle{plain}{%
  \fancyhf{}
  \fancyfoot[C]{\small\color{gray!70}\thepage}
  \renewcommand{\headrulewidth}{0pt}%
}


% ══════════════════════════════════════════════════════════
%  LINKS & REFERENCES
% ══════════════════════════════════════════════════════════
\usepackage[hidelinks]{hyperref}
\usepackage[nameinlink, capitalize]{cleveref}


% ══════════════════════════════════════════════════════════
%  DOCUMENT METADATA  (fill in for each new subject)
% ══════════════════════════════════════════════════════════
\newcommand{\coursetitle}{Biological Sequence Analysis}
\newcommand{\documenttype}{Study Notes -- Working Draft}
\newcommand{\authorname}{Stefan Eder}
\newcommand{\basedon}{Based on the lecture \emph{Biological Sequence Analysis} (365.060; P.~Seidl, G.~Klambauer, A.~Regl; JKU Linz), WS 2026/27}


% ══════════════════════════════════════════════════════════
%  MATHEMATICAL MACROS
%  Organized by domain. Add your own at the end of each group.
% ══════════════════════════════════════════════════════════

% ── Number Sets & Spaces ──────────────────────────────────
\newcommand{\R}{\mathbb{R}}
\newcommand{\N}{\mathbb{N}}
\newcommand{\Z}{\mathbb{Z}}
\newcommand{\C}{\mathbb{C}}
\newcommand{\Q}{\mathbb{Q}}

% ── Probability & Statistics ──────────────────────────────
\newcommand{\E}{\mathbb{E}}              % Expectation
\newcommand{\Prob}{\mathbb{P}}           % Probability measure
\newcommand{\Var}{\mathrm{Var}}
\newcommand{\Cov}{\mathrm{Cov}}
\newcommand{\Corr}{\mathrm{Corr}}
\newcommand{\iid}{\overset{\mathrm{iid}}{\sim}}
% Indicator function. Usage: \ind_{\{x > 0\}}
\newcommand{\ind}{\mathbbm{1}}

% ── Information Theory ────────────────────────────────────
% KL divergence. Usage: \KL{q}{p}  →  D_KL(q ‖ p)
\newcommand{\KL}[2]{D_{\mathrm{KL}}\!\left(#1 \,\|\, #2\right)}
\newcommand{\Ent}{\mathrm{H}}            % Entropy:         \Ent(X)
\newcommand{\MI}{\mathrm{I}}             % Mutual information: \MI(X;Y)

% ── Calculus & Optimization ───────────────────────────────
\newcommand{\argmax}{\operatorname*{arg\,max}}
\newcommand{\argmin}{\operatorname*{arg\,min}}
\newcommand{\grad}{\nabla}               % Gradient
\newcommand{\hess}{\nabla^2}             % Hessian
\newcommand{\lapl}{\Delta}               % Laplacian
% Differential (upright d). Usage: \int f(x) \dd x
\newcommand{\dd}{\,\mathrm{d}}
% Partial derivatives. Usage: \pd{f}{x}  or  \pdd{f}{x}
\newcommand{\pd}[2]{\frac{\partial #1}{\partial #2}}
\newcommand{\pdd}[2]{\frac{\partial^2 #1}{\partial #2^2}}

% ── Norms, Abs, Inner Products ────────────────────────────
\newcommand{\norm}[1]{\left\lVert #1 \right\rVert}
\newcommand{\normsq}[1]{\left\lVert #1 \right\rVert^2}
\newcommand{\abs}[1]{\left\lvert #1 \right\rvert}
\newcommand{\inner}[2]{\left\langle #1,\, #2 \right\rangle}

% ── Linear Algebra ────────────────────────────────────────
\newcommand{\T}{^{\mathsf{T}}}           % Transpose:        \mathbf{A}\T
\newcommand{\inv}{^{-1}}                 % Inverse:          A\inv
\newcommand{\pinv}{^{\dagger}}           % Pseudo-inverse:   A\pinv
\newcommand{\tr}{\operatorname{tr}}
\newcommand{\rank}{\operatorname{rank}}
\newcommand{\diag}{\operatorname{diag}}
\newcommand{\spanop}{\operatorname{span}}
\newcommand{\eye}{\mathbf{I}}            % Identity matrix
\newcommand{\zeros}{\mathbf{0}}
\newcommand{\ones}{\mathbf{1}}

% ── ML / AI Core Notation ────────────────────────────────
\newcommand{\loss}{\mathcal{L}}          % Loss:          \loss(\params)
\newcommand{\reg}{\mathcal{R}}           % Regularizer:   \reg(\params)
\newcommand{\dataset}{\mathcal{D}}       % Dataset:       \dataset
\newcommand{\params}{\boldsymbol{\theta}}   % Parameters
\newcommand{\weights}{\mathbf{w}}
\newcommand{\xvec}{\mathbf{x}}
\newcommand{\yvec}{\mathbf{y}}
\newcommand{\zvec}{\mathbf{z}}
\newcommand{\hvec}{\mathbf{h}}           % Hidden state
\newcommand{\feature}{\boldsymbol{\phi}} % Feature map

% Activation functions
\DeclareMathOperator{\softmax}{softmax}
\DeclareMathOperator{\relu}{ReLU}
\DeclareMathOperator{\sign}{sign}
% sigmoid: just use \sigma (standard Greek letter)

% ── Complexity ────────────────────────────────────────────
\newcommand{\bigO}[1]{\mathcal{O}\!\left(#1\right)}
\newcommand{\bigOm}[1]{\Omega\!\left(#1\right)}
\newcommand{\bigTh}[1]{\Theta\!\left(#1\right)}

% ── Common Distributions ──────────────────────────────────
\newcommand{\Normal}{\mathcal{N}}
\newcommand{\Uniform}{\mathcal{U}}
\newcommand{\Bernoulli}{\mathrm{Ber}}
\newcommand{\Categorical}{\mathrm{Cat}}
\newcommand{\Dirichlet}{\mathrm{Dir}}
\newcommand{\Binomial}{\mathrm{Bin}}

% ── Miscellaneous ─────────────────────────────────────────
\newcommand{\given}{\,\vert\,}           % Conditional: \Prob(A \given B)
\newcommand{\st}{\;\text{s.t.}\;}        % "subject to"
\newcommand{\eqdef}{\coloneqq}           % Definitional equality  :=
\newcommand{\eqq}{\overset{?}{=}}        % "Does this equal?"

% ══════════════════════════════════════════════════════════
%  COURSE-SPECIFIC MACROS  (Biological Sequence Analysis)
% ══════════════════════════════════════════════════════════
\usepackage{multirow}
\usepackage{bigdelim}
\usepackage{chemfig}                     % small structural formulas
\setchemfig{atom sep=1.6em, bond offset=1pt, double bond sep=2.5pt}

% Sequences and alignments in monospace
\newcommand{\seq}[1]{\texttt{#1}}
% An alignment block: lines separated by \\ ; use ~ for a space
\newcommand{\alnblock}[1]{%
  \par\smallskip
  {\centering\ttfamily\small
   \begin{tabular}{@{}l@{}}#1\end{tabular}\par}%
  \smallskip}
\newcommand{\ma}[1]{\textcolor{ex@frame}{\textbf{#1}}}   % highlighted match
\newcommand{\mm}[1]{\textcolor{thm@frame}{\textbf{#1}}}  % highlighted mismatch

% Log-odds notation
\newcommand{\obs}{\mathrm{observed}}
\newcommand{\expd}{\mathrm{expected}}
\newcommand{\PAM}[1]{\mathrm{PAM}#1}
\newcommand{\BLOSUM}[1]{\mathrm{BLOSUM}#1}

% Lecture reference in chapter openings
\newcommand{\lectureref}[1]{%
  \par{\small\color{gray!70!black}\textit{Lecture:} #1}\par}

% ── Figure colours (aliases of the template palette) ──────
\colorlet{figblue}{def@frame}
\colorlet{figred}{thm@frame}
\colorlet{figgreen}{ex@frame}
\colorlet{figamber}{int@frame}
\colorlet{figpurple}{note@frame}

% ── TikZ styles ───────────────────────────────────────────
\tikzset{
  box/.style    = {draw=accent, thick, rounded corners=2pt, fill=accent!6,
                   inner sep=4pt, font=\small, align=center},
  flow/.style   = {-{Stealth[length=2.4mm]}, thick, accent!85!black},
  lbl/.style    = {font=\scriptsize, inner sep=1pt},
  base/.style   = {draw, thick, minimum width=6mm, minimum height=5mm,
                   inner sep=0pt, font=\small\ttfamily\bfseries},
}
\colorlet{baseA}{ex@frame!25}  \colorlet{baseT}{thm@frame!22}
\colorlet{baseG}{int@frame!30} \colorlet{baseC}{def@frame!22}

% ── pgfplots look shared by all plots ─────────────────────
\pgfplotsset{
  bsa/.style = {
    axis lines = left,
    axis line style = {gray!70},
    tick label style = {font=\scriptsize, gray!70!black},
    label style = {font=\small},
    grid = major, grid style = {gray!15},
  },
}

% ── Placeholders ──────────────────────────────────────────
% Figure placeholder (same signature as in other course notes).
% #1 height, #2 concept, #3 what the picture should show + search hint,
% #4 caption, #5 label.  Replace with \includegraphics{figures/...}
% once a suitable image is in figures/.
\newcommand{\figureplaceholder}[5]{%
\begin{figure}[htbp]
\centering
\begin{tcolorbox}[unbreakable,width=\linewidth,height=#1,colback=accent!3,
  colframe=accent!30,boxrule=0.5pt,arc=2.5mm,sharp corners=all,
  valign=center,halign=center,fontupper=\small]
{\bfseries\color{accent}Figure to insert: #2}\par\medskip
{\color{gray!75!black}#3}
\end{tcolorbox}
\caption{#4}\label{#5}
\end{figure}}

% Outline text for chapters that are not written yet.
\newcommand{\planned}[1]{%
  \par{\small\color{gray!70!black}\textit{Planned content:} #1}\par}

% Keep section headings together with the start of their text.
\let\originalsection\section
\renewcommand{\section}{\Needspace{7\baselineskip}\originalsection}
\let\originalsubsection\subsection
\renewcommand{\subsection}{\Needspace{6\baselineskip}\originalsubsection}


% ==========================================================
%  BEGIN DOCUMENT
% ==========================================================
\begin{document}

\frontmatter  % Roman-numeral page numbers for front matter


% ──────────────────────────────────────────────────────────
%  TITLE PAGE
% ──────────────────────────────────────────────────────────
\thispagestyle{empty}
\vspace*{1.8cm}

\begin{center}
  {\color{accent}\rule{\linewidth}{1.8pt}}
  \vspace{0.65cm}

  {\fontsize{28}{34}\selectfont\bfseries\color{accent}\coursetitle\par}
  \vspace{0.25cm}
  {\large\color{gray!80!black}\documenttype\par}

  \vspace{0.55cm}
  {\color{accent!35}\rule{\linewidth}{0.5pt}}
  \vspace{0.55cm}

  {\large\authorname\par}
  \vspace{0.2cm}
  {\small\color{gray!70}\basedon\par}
\end{center}

\vspace{1.6cm}

\begin{center}
\begin{minipage}{0.84\textwidth}
  \small\color{gray!80!black}
  These notes are written for people who, like the author, come to this
  course from computer science and not from biology. Every biological idea is
  explained from scratch, with analogies to things a programmer already
  knows, and every scoring idea comes with a small example with actual
  numbers. The notes combine the lecture slides with what was \emph{said} in
  the recorded lectures --- explanations, side remarks, and hints about what
  matters for the exam. All worked numbers were re-computed by script.
  Use the notes as a study aid alongside the slides, not as a replacement for
  them.

  \medskip
  \textbf{Working draft:} Chapters 1--9 (lectures 1--4: biology intro,
  scoring, PAM, BLOSUM, tools) are written out. The remaining chapters
  (pairwise and multiple alignment, phylogenetics, HMMs, AI-based sequence
  models) are outlines and will be filled in as the lecture progresses.

  \medskip
  If you spot an error or a gap, please open an issue or pull request at
  \textbf{github.com/steveTgit/AI\_Master} --- corrections are always welcome.

  \medskip
  \color{gray!60!black}
  Good study materials should be shared, not gatekept.
  If this helped you, take it as a reason to help the next colleague you see struggling with something.
\end{minipage}
\end{center}

\vfill
{\centering\small\color{gray!65} \today\par}
\clearpage


% ──────────────────────────────────────────────────────────
%  TABLE OF CONTENTS
% ──────────────────────────────────────────────────────────
\setcounter{tocdepth}{1}  % chapters and sections

\begingroup
  \let\cleardoublepage\relax
  \let\clearpage\relax
  \tableofcontents
\endgroup

\mainmatter  % Switch to Arabic numerals


% ==========================================================
%  CHAPTER 1 — COURSE OVERVIEW AND MOTIVATION
%  Lecture VO01 (5.10., in person), first ~20 minutes
%  Slides: slides/VO01_bio-intro-1.pdf, p. 1--22
% ==========================================================
\chapter{Course Overview and Motivation}
\label{ch:intro}

\lectureref{VO01 (Bio-Intro I), slides 1--22. Recording: first 20 minutes.}

This course is about \emph{biological sequences} --- long strings over small
alphabets that encode how living things are built --- and about the
algorithms we use to compare them. Before any of that makes sense we need
two things: a rough idea of what the course looks like, and a reason to care.
This chapter provides both.

\section{What Is Sequence Analysis?}

The slides open with a definition:

\begin{defbox}[title={Sequence analysis}]
Sequence analysis is the ``process of applying analytical methods to DNA,
RNA, or protein sequences to retrieve meaningful information. Among the
sequence analysis methods, sequence alignment and similarity analysis are
widely used \ldots{} The main goal of sequence alignment and similarity
analysis is to show similarities and differences among biological
sequences.''
\end{defbox}

The field was born when DNA was discovered and people suddenly had a new kind
of data --- strings of the letters \seq{A}, \seq{C}, \seq{G}, \seq{T} ---
and needed methods to analyse it. It is now a sub-field of bioinformatics.
The main focus is DNA and everything derived from it: RNA is copied from DNA,
and proteins are built according to RNA (\cref{ch:bio1}). But most methods in
this course work on \emph{any} kind of sequence; if you understand them for
DNA, you understand them for text, time series or source code.

\begin{intbox}
Think of sequence analysis as \emph{diff for biology}. When you run
\texttt{diff} on two versions of a file, you want to know which lines were
kept, changed, inserted or deleted. Aligning two DNA or protein sequences
asks exactly the same question --- except that the ``versions'' are the
same gene in two species, separated by millions of years of evolution, and
we want to know how closely related they are.
\end{intbox}

\section{Organisation of the Course}

\begin{itemize}
  \item \textbf{People.} Philipp Seidl leads the course in WS 2026/27,
    together with Günter Klambauer (on sabbatical this year, so everything
    goes through Philipp). Earlier editions were taught by Florian Sestak and
    by Alois Regl, who still gives the scoring lectures and whose slides much
    of the course is based on.
  \item \textbf{Format.} A full distance-learning lecture: 13 lectures, two of
    them in person (Bio-Intro~I and the AI-based sequence models), the rest as
    video releases, with a Zoom Q\&A roughly every third week. All material
    is on Moodle; general questions go to the Moodle forum, personal ones to
    \texttt{seidl@ml.jku.at}.
  \item \textbf{Exam.} Written, pen and paper, \textbf{in person} (Linz, or
    the distance-learning centres in Bregenz and Vienna). First date: Monday
    25.01.2027; retry exams about a month later; a third one only on request.
    No registration means no grade.
  \item \textbf{Exercise.} Separate course (UE), weekly, mostly Python;
    students present their solutions. The exercise may run a bit ahead of
    the lecture.
\end{itemize}

\Cref{tab:roadmap} shows the schedule and how the chapters of these notes
correspond to the lectures and to the files in the course folder
(\texttt{slides/}, \texttt{recordings/}, \texttt{transcripts/} share the same
\texttt{VOxx\_\ldots} names).

\begin{table}[htbp]
\centering
\small
\begin{tabular}{@{}llll@{}}
\toprule
\textbf{Date} & \textbf{Topic} & \textbf{Lecture / files} & \textbf{Chapter} \\
\midrule
5.10.  & Bio-Intro I                         & VO01 \texttt{bio-intro-1}  & \ref{ch:intro}, \ref{ch:bio1} \\
12.10. & Bio-Intro II                        & VO02 \texttt{bio-intro-2}  & \ref{ch:bio2} \\
19.10. & Scoring: PAM                        & VO03 \texttt{scoring-1}    & \ref{ch:alignment}, \ref{ch:scoring}, \ref{ch:logodds} \\
9.11.  & Scoring: BLOSUM                     & VO04 \texttt{scoring-2}    & \ref{ch:logodds}--\ref{ch:tools} \\
16.11. & Pairwise alignment: global          & --                          & \ref{ch:global} (outline) \\
23.11. & Pairwise alignment: local           & --                          & \ref{ch:local} (outline) \\
30.11., 7.12. & Multiple alignment           & --                          & \ref{ch:msa} (outline) \\
14.12. & Phylogenetics                       & --                          & \ref{ch:phylo} (outline) \\
21.12. & HMM basics                          & --                          & \ref{ch:hmm} (outline) \\
11.01., 18.01. & AI-based biological sequence models & --              & \ref{ch:ai} (outline) \\
25.01. & Exam                                & & \\
\bottomrule
\end{tabular}
\caption{Course schedule (WS 2026/27) and where to find each lecture in these
notes. VO04 also contains the ``add-ons'' (tools and a case study,
\cref{ch:tools}).}
\label{tab:roadmap}
\end{table}

\section{Why Should We Care? Motivation}

\subsection{Sequences Are the Languages of Life}

The lecture starts with the 2024 Nobel Prize in Chemistry, awarded to David
Baker, Demis Hassabis and John Jumper --- the latter two for
\textbf{AlphaFold}. Given the sequence of amino acids of a protein,
AlphaFold predicts the three-dimensional shape the protein will fold into.
This was a famously hard problem: people spent entire PhDs determining the 3D
structure of a single protein experimentally. Knowing the shape matters,
because the shape determines what a protein does --- and therefore how
diseases work and how drugs could act on them. As we will see in
\cref{sec:alphafold-msa}, AlphaFold's success rests heavily on the topic of
this course: it searches for \emph{similar sequences} and aligns them.

Life on Earth is old: single-celled organisms (prokaryotes) appeared
billions of years ago, then cells with a nucleus (eukaryotes), then fungi,
plants and animals. A striking event in this history is the \emph{Cambrian
explosion}, a geologically short period in which a huge diversity of animal
species appeared; why it happened so suddenly is still not fully understood.
Compared with all of this, our own existence is a tiny sliver of time. And a
remarkable amount of what an organism is --- how it looks, which diseases it
is prone to --- is encoded in its DNA.

\begin{figure}[htbp]
\centering
\begin{tikzpicture}[node distance=4mm]
  \node[box, text width=3.6cm, minimum height=2.5cm] (dna) {%
    \textbf{DNA}\\[2pt]
    {\scriptsize alphabet: 4 bases}\\
    \texttt{\scriptsize GGACGGAGGGGCTCC\ldots}\\[2pt]
    {\scriptsize human genome: $\approx 3$\,GB}};
  \node[box, text width=3.6cm, minimum height=2.5cm, right=of dna] (prot) {%
    \textbf{Protein}\\[2pt]
    {\scriptsize alphabet: 20 amino acids}\\
    \texttt{\scriptsize LCLYTHIGRNIYYGS\ldots}\\[2pt]
    {\scriptsize human: $\approx 20{,}000$ proteins}};
  \node[box, text width=3.6cm, minimum height=2.5cm, right=of prot] (mol) {%
    \textbf{Small molecules}\\[2pt]
    {\scriptsize written as SMILES strings}\\
    \texttt{\scriptsize CC(=O)Oc1ccccc1C(=O)O}\\[2pt]
    {\scriptsize (here: aspirin)}};
\end{tikzpicture}
\caption{Three ``languages of life'' that can be written as strings: DNA,
proteins and small molecules. This course deals with the first two.}
\label{fig:languages}
\end{figure}

If you sequenced your own genome, you would get roughly 3\,GB of text over
the alphabet $\{\seq{A}, \seq{C}, \seq{G}, \seq{T}\}$; compressed, about
1\,GB, because much of it is repetitive. Imagine being the first person to
look at that string: it is like an alien code, and we had to learn how to read
it. Proteins are a second language, written in amino acids, and molecules
can be written as strings too (SMILES), see \cref{fig:languages}.

\subsection{Genotype and Phenotype}

\begin{defbox}[title={Genotype vs.\ phenotype}]
The \textbf{genotype} is the genetic information an organism carries in its
DNA. The \textbf{phenotype} is what actually shows up: its appearance,
properties and behaviour. The two are not the same --- a trait can be
present in the DNA without being expressed (for example a recessive gene),
and the environment influences the phenotype as well.
\end{defbox}

\subsection{Questions Sequence Analysis Can Answer}

Sequence analysis gives us the tools to decode, compare and infer function
from this huge amount of data; phylogenetics (\cref{ch:phylo}) traces the
evolutionary relationships between organisms. Typical questions:
\begin{itemize}
  \item How do we identify genes (the parts of DNA that code for proteins)
    or other functional elements in a genome?
  \item What are the evolutionary relationships between species?
  \item Which mutations and genetic variants are associated with diseases?
  \item How do proteins evolve, and what are the functional consequences of
    their sequence variations?
\end{itemize}

And some concrete applications mentioned in the lecture:
\begin{itemize}
  \item \textbf{Epidemiology.} Sequencing the DNA/RNA found in sewage water
    reveals which coronavirus variants circulate in a city --- often before
    they show up in hospital data.
  \item \textbf{Forensics.} Companies such as 23andMe do not sequence the whole
    genome; they use microarrays that read only the positions that vary a
    lot between people. That is cheap, and it is enough to find
    \emph{relatives} of a person in their databases --- which has been used
    to solve criminal cases from DNA found at a crime scene.
  \item \textbf{Human origins.} DNA is chemically very stable (it can survive
    in bones for hundreds of thousands of years). Comparing ancient and
    modern DNA shows that modern humans originated in Africa and spread from
    there across the world.
  \item \textbf{Neanderthals.} Are they our ancestors? No: they are
    considered a separate species that split from our lineage roughly
    550,000--750,000 years ago. But there was interbreeding (about
    40,000--54,000 years ago), so many people today carry some Neanderthal
    genes.
  \item \textbf{The origin of SARS-CoV-2.} A study compared the virus's
    sequence with earlier coronavirus sequences in databases and asked how
    likely the observed changes are to arise naturally; it concluded that a
    natural origin is likely.
  \item And more: origins of life, evolution of diseases, and the chain
    \emph{amino-acid sequence $\to$ 3D structure $\to$ function $\to$
    intervention} (drug design).
\end{itemize}

\begin{notebox}[title={What is a species?}]
A common working definition: two organisms belong to the same species if
they can have \emph{fertile} offspring. \emph{Hybrids} show that the line is
blurry: a horse and a donkey can produce a mule, but the mule is (normally)
sterile, so horses and donkeys are different species.
\end{notebox}

\section{Chapter Checklist}

\begin{itemize}
  \item what sequence analysis is, and that alignment and similarity analysis
    are its core tools;
  \item course organisation: distance learning, in-person written exam,
    separate Python exercise;
  \item why the topic matters: AlphaFold, sequences as languages of life,
    and applications from epidemiology to forensics;
  \item genotype vs.\ phenotype; the fertile-offspring definition of a species.
\end{itemize}

% ==========================================================
%  CHAPTER 2 — BIOLOGY BASICS I: CELLS, CENTRAL DOGMA, DNA, RNA
%  Lecture VO01 (5.10.), slides 23--59, recording 17:00--61:00
% ==========================================================
\chapter{Biology Basics I: Cells, the Central Dogma, DNA and RNA}
\label{ch:bio1}

\lectureref{VO01 (Bio-Intro I), slides 23--59. Recording from about 17:00.}

You do not need to become a biologist for this course --- the lecturers
say so explicitly. But every algorithm we will meet makes biological
assumptions (``mutations are rare'', ``some substitutions are more likely
than others'', ``proteins are read in triplets''), and they only make sense
if you know what the strings represent. This chapter covers the first half
of the biology crash course: where the information lives (the cell), how it
flows (the central dogma), and how it is stored (DNA and RNA).

\begin{intbox}
A mental model that carries surprisingly far: a cell is a tiny
\textbf{factory with a computer inside}. The \textbf{DNA} is the source code
on a protected hard drive. Copies of individual ``files'' (genes) are made as
\textbf{RNA} and sent to the factory floor, where \textbf{ribosomes} read them
and assemble \textbf{proteins} --- the machines that do practically all the
work. Most of the biology in this course is about this pipeline.
\end{intbox}

\section{The Cell}

\begin{defbox}[title={Cell}]
The cell is the smallest fundamental, self-replicating unit of life --- and
at the same time the largest unit that \emph{all} life forms have in common.
Every organism is either a single cell or built from many cells.
\end{defbox}

\subsection{Prokaryotes and Eukaryotes}

The first and most important distinction is a single yes/no question: does
the cell have a \textbf{nucleus}, a membrane-enclosed compartment that holds
the DNA?

\begin{table}[htbp]
\centering
\small
\begin{tabular}{@{}p{3.0cm}p{5.4cm}p{5.4cm}@{}}
\toprule
 & \textbf{Prokaryotes} (no nucleus) & \textbf{Eukaryotes} (nucleus) \\
\midrule
Examples & bacteria, archaea & animals, plants, fungi \\
DNA & floats freely in the cell & locked in the nucleus (``vault'') \\
Size (typical) & 1--5\,\textmu m (a fine dust particle) & 10--100\,\textmu m (thickness of a hair) \\
Internal structure & few compartments & many organelles \\
From DNA to protein & directly: transcription and translation in the same space & extra step: the information has to be copied out of the nucleus first \\
\bottomrule
\end{tabular}
\caption{The two big kinds of cells.}
\label{tab:pro-eu}
\end{table}

\textbf{Archaea} deserve a special mention: they are prokaryotes that live in
places where you would not expect life at all, for example near deep-sea
hydrothermal vents at temperatures above 100\,\textdegree C. Among the
eukaryotes, the slides contrast three kinds of cells:
\begin{itemize}
  \item \textbf{animal cells}: flexible boundary, no cell wall --- built for
    mobility and consumption;
  \item \textbf{plant cells}: rigid cell wall (cellulose) and chloroplasts ---
    built for growing upwards and harvesting solar energy;
  \item \textbf{fungal cells}: rigid cell wall (chitin) --- built to penetrate
    their environment and absorb nutrients.
\end{itemize}
A question from the lecture: are animal cells more closely related to plant
cells or to fungal cells? Surprisingly, to \textbf{fungi}: in the tree of life,
animals and fungi share a more recent common ancestor than either does with
plants.

\subsection{Inside a Eukaryotic Cell}

A eukaryotic cell contains many specialised compartments, the
\emph{organelles}. You do not need to memorise their chemistry, but you
should know what each one is for (\cref{tab:organelles}, \cref{fig:cell}).

\begin{table}[htbp]
\centering
\small
\renewcommand{\arraystretch}{1.15}
\begin{tabular}{@{}p{3.3cm}p{6.6cm}p{3.9cm}@{}}
\toprule
\textbf{Organelle} & \textbf{Function} & \textbf{Analogy} \\
\midrule
nucleus & stores the DNA, packed into chromosomes & vault / hard drive \\
nucleolus & inside the nucleus; assembles ribosome subunits & \\
ribosome & complex of RNA and protein that builds proteins from mRNA (80S in the cytoplasm of eukaryotes, 70S in prokaryotes and mitochondria) & factory / 3D printer \\
mitochondria & power generators (oxidative phosphorylation); also involved in programmed cell death; have their own DNA, inherited from the mother only & power plant \\
chloroplasts (plants) & photosynthesis; also have their own genome & solar panels \\
endoplasmic reticulum (ER) & transport network; the rough ER carries ribosomes on its surface and produces proteins that are shipped onwards & road network \\
Golgi apparatus & processes and packages macromolecules into vesicles & post office / parcel service \\
lysosomes & enzymes that digest old organelles, food particles and engulfed viruses or bacteria & recycling centre \\
peroxisomes & get rid of toxic peroxides & waste disposal \\
cytoskeleton, centrosome & scaffolding; microtubules direct transport and pull chromosomes apart during cell division & scaffolding, cranes \\
vacuoles & storage (in animals: shipping containers; in plants and fungi: big water-filled ``hydraulic cylinders'') & warehouse \\
cytoplasm & the fluid filling the cell --- crowded with proteins and other molecules & factory floor \\
plasma membrane & boundary that protects the inside & walls \\
\bottomrule
\end{tabular}
\caption{The organelles of a eukaryotic cell (after slide 27 and the
lecture's explanations).}
\label{tab:organelles}
\end{table}

\figureplaceholder{5.5cm}{labelled eukaryotic (animal) cell}%
{A cross-section of an animal cell with the organelles from
\cref{tab:organelles} labelled: nucleus, nucleolus, ribosomes, rough/smooth
ER, Golgi apparatus, mitochondria, lysosomes, cytoskeleton, plasma membrane.
\textit{Search:} ``eukaryotic animal cell diagram labeled organelles''
(e.g.\ Britannica or Wikimedia Commons ``Animal cell structure en.svg'').}%
{A eukaryotic cell and its organelles.}{fig:cell}

The most important point for this course is the separation of storage and
production. The DNA sits in the nucleus, but the factories (ribosomes) are
outside, in the cytoplasm. DNA is far too large to leave the nucleus, so the
cell \emph{copies the relevant part into RNA}, and the RNA travels out
through pores in the nuclear envelope. Prokaryotes do not need this step:
their DNA lies around in the same space as the ribosomes.

\subsection{From Organelles to Organisms}

Biology is organised in levels (\cref{fig:hierarchy}): organelles are parts
of cells, many cells form a tissue, tissues form organs, organs work together
in organ systems, and all of them together make up the organism.

\begin{figure}[htbp]
\centering
\begin{tikzpicture}[node distance=3.5mm, every node/.style={box, minimum height=7mm}]
  \node (a) {organelle};
  \node[right=of a] (b) {cell};
  \node[right=of b] (c) {tissue};
  \node[right=of c] (d) {organ};
  \node[right=of d] (e) {organ system};
  \node[right=of e] (f) {organism};
  \foreach \x/\y in {a/b,b/c,c/d,d/e,e/f} \draw[flow] (\x) -- (\y);
\end{tikzpicture}
\caption{Levels of organisation, from micro to macro (e.g.\ mitochondrion
$\to$ intestinal cell $\to$ intestinal lining $\to$ intestine $\to$ digestive
system $\to$ human).}
\label{fig:hierarchy}
\end{figure}

Cells are also extremely diverse in shape and size: a motor neuron, a red
blood cell, a sperm cell and an egg cell look nothing alike. The human egg is
one of the largest cells and just about visible with the naked eye; the
longest is the sciatic motor neuron, which runs from the spine to the toe
(up to about 1.4\,m). Mature human red blood cells are special: they have no
nucleus and no DNA. And a cell looks very different depending on its stage
in the cell cycle --- a dividing cell, with its condensed chromosomes, looks
nothing like a resting one.

\subsection{Some Numbers}

\begin{itemize}
  \item A human consists of roughly $30 \times 10^{12}$ cells (older estimates
    said 100 trillion).
  \item We carry about as many bacteria as human cells (ratio about 1.3 : 1,
    not the often-quoted 10 : 1), about ten times as many viruses, and a small
    amount of fungi (0.1\,\%). ``We'' are quite an ecosystem.
  \item The DNA in a single nucleus is about 2\,m long when stretched out.
    Over all cells, that adds up to about 200 round trips from the Earth to
    the Sun.\footnote{$3 \times 10^{13}$ cells $\times$ 2\,m $= 6 \times
    10^{13}$\,m, and one round trip is $2 \times 1.5 \times 10^{11}$\,m, so
    about 200. Another slide says 500; the 200 is consistent with the cell
    count.}
  \item Every day, 50--70 billion of our cells die in a controlled way
    (\emph{apoptosis}, programmed cell death) --- about 10\,g of tissue.
  \item The first cells appeared about 3.5 billion years ago; cells were
    discovered by Robert Hooke in 1665.
\end{itemize}

\section{The Central Dogma of Molecular Biology}
\label{sec:central-dogma}

Who tells a cell what to do, and when? Inside the cell, the work is done by
\textbf{proteins}: they act as sensors, as data cables and as engines. But a
cell cannot simply build everything all the time --- that would waste
energy. It needs strict control over \emph{which} protein to build,
\emph{when}, and \emph{where} to send it. That requires an organised
``source code'': the DNA, plus a pipeline that turns code into machines.

\begin{defbox}[title={Central dogma of molecular biology}]
Genetic information flows
\[
  \text{DNA} \;\xrightarrow{\ \text{transcription}\ }\; \text{RNA}
  \;\xrightarrow{\ \text{translation}\ }\; \text{protein},
\]
and DNA is copied by \textbf{replication}. The DNA is the library of
instructions (``genes'') for how to build proteins.
\begin{itemize}
  \item \textbf{Replication} (DNA $\to$ DNA): the genome is copied before a cell
    divides (enzyme: DNA polymerase).
  \item \textbf{Transcription} (DNA $\to$ RNA): a gene is copied into a
    messenger RNA (enzyme: RNA polymerase). The ``language'' stays the same ---
    nucleotides in, nucleotides out --- which is why it is called a
    \emph{transcription}.
  \item \textbf{Translation} (RNA $\to$ protein): the ribosome reads the mRNA
    and builds a chain of amino acids. Now the language changes --- from
    nucleotides to amino acids --- hence \emph{translation}.
\end{itemize}
\end{defbox}

\begin{figure}[htbp]
\centering
\begin{tikzpicture}[node distance=3.3cm]
  \node[box, minimum width=2.2cm, minimum height=9mm] (dna) {\textbf{DNA}};
  \node[box, minimum width=2.2cm, minimum height=9mm, right=of dna] (rna) {\textbf{RNA}};
  \node[box, minimum width=2.2cm, minimum height=9mm, right=of rna] (prot) {\textbf{Protein}};
  \draw[flow] (dna) -- node[above, lbl] {transcription} node[below, lbl] {RNA polymerase} (rna);
  \draw[flow] (rna) -- node[above, lbl] {translation} node[below, lbl] {ribosome} (prot);
  \draw[flow] (dna.north west) to[out=120, in=60, looseness=4] node[above, lbl] {replication (DNA polymerase)} (dna.north east);
  % exceptions
  \draw[-{Stealth[length=2.2mm]}, thick, dashed, figgreen] (rna.south west) to[out=-150, in=-30] node[below, lbl, figgreen!70!black] {reverse transcription (retroviruses)} (dna.south east);
  \draw[-{Stealth[length=2.2mm]}, thick, dashed, figgreen] (rna.north west) to[out=120, in=60, looseness=4] node[above, lbl, figgreen!70!black] {RNA replication (RNA viruses)} (rna.north east);
\end{tikzpicture}
\caption{The central dogma (solid) and the most important ``violations''
found later (dashed): reverse transcription and RNA replication. Prions are
a third, special case: misfolded proteins that pass on their \emph{shape}
to other proteins.}
\label{fig:central-dogma}
\end{figure}

\paragraph{Why ``dogma''}
Francis Crick named it so in the late 1950s. As the lecture tells it, he
misunderstood the word: in science one talks of hypotheses and theories, and
``dogma'' sounds almost religious. The name stuck anyway --- biology is full
of such historical names (``chromosome'' means ``coloured body'', because
these structures took up dye under the microscope).

\paragraph{Exceptions}
For about 20 years after the discovery of DNA, the dogma was believed to hold
strictly. Since then, several exceptions have been found
(\cref{fig:central-dogma}):
\begin{itemize}
  \item \textbf{reverse transcription} (RNA $\to$ DNA): retroviruses such as
    HIV write their RNA genome back into the DNA of the host cell;
  \item \textbf{RNA replication} (RNA $\to$ RNA): many RNA viruses copy their
    RNA directly, with an RNA-dependent RNA polymerase;
  \item \textbf{prions}: proteins that can change the shape of other
    proteins of the same kind, causing diseases such as mad cow disease (BSE)
    and Creutzfeldt--Jakob disease.
\end{itemize}

\paragraph{Prokaryotes vs.\ eukaryotes}
In prokaryotes the pipeline is short: DNA is transcribed into mRNA, and the
mRNA is translated right away. In eukaryotes there are extra steps
(\cref{ch:bio2}): transcription happens in the nucleus; the primary RNA
transcript is processed (a protective ``cap'' and a poly-A ``tail'' are
added, and the non-coding parts are \emph{spliced} out); only then is the
mature mRNA exported to the cytoplasm and translated. The extra complexity
pays off: splicing allows the cell to reuse parts of a gene in different
combinations (\cref{sec:splicing}).

\section{DNA}

\subsection{A Short History}

The structure of DNA was identified in 1953 by James Watson and Francis
Crick --- with decisive help from \textbf{Rosalind Franklin}, whose X-ray
image ``Photo~51'' revealed the helical structure. Whether Franklin got the
credit she deserved is a well-known controversy. The key insight: DNA is
\emph{double-stranded}, and the two strands wind around each other into the
famous \textbf{double helix}.

\subsection{Building Blocks: Nucleotides}

\begin{defbox}[title={DNA (deoxyribonucleic acid)}]
DNA is a long chain of \textbf{nucleotides}. Each nucleotide consists of
\begin{itemize}
  \item a \textbf{sugar} (deoxyribose, a ring of five carbon atoms),
  \item a \textbf{phosphate} group, which links the sugar to the next sugar
    (sugar and phosphate form the \emph{backbone}),
  \item one of four \textbf{bases}: adenine (\seq{A}), cytosine (\seq{C}),
    guanine (\seq{G}) and thymine (\seq{T}).
\end{itemize}
The information is in the order of the bases --- so for us, DNA is simply a
string over the alphabet $\{\seq{A}, \seq{C}, \seq{G}, \seq{T}\}$.
\end{defbox}

\figureplaceholder{4.2cm}{one nucleotide, chemically}%
{A single nucleotide with the phosphate group, the 5-carbon sugar (carbons
numbered 1$'$ to 5$'$) and the nitrogenous base; ideally two nucleotides of the
backbone to show how phosphate links the 5$'$ carbon of one sugar to the 3$'$
carbon of the next. \textit{Search:} ``DNA nucleotide structure phosphate
deoxyribose base 5 prime 3 prime''.}%
{The chemical structure of a nucleotide.}{fig:nucleotide}

\subsection{Base Pairing}

The two strands are held together by \textbf{hydrogen bonds} between the
bases, and only two pairings fit:
\[
  \seq{A} \text{---} \seq{T} \quad (2 \text{ hydrogen bonds}),
  \qquad
  \seq{G} \text{---} \seq{C} \quad (3 \text{ hydrogen bonds}).
\]
So each strand completely determines the other: the strands are
\textbf{complementary}. A G--C pair, with its three hydrogen bonds, is harder
to break than an A--T pair. Organisms that live in harsh conditions, such as
heat-loving archaea, therefore tend to have a higher \textbf{GC content}.

\subsection{Direction: 5\texorpdfstring{$'$}{'} and 3\texorpdfstring{$'$}{'}}

The carbon atoms of the sugar are numbered $1'$ to $5'$ (``one prime'' to
``five prime''). In the backbone, the phosphate connects the $5'$ carbon of
one sugar with the $3'$ carbon of the next. A strand therefore has a
direction, from its \textbf{$5'$ end} to its \textbf{$3'$ end}, and the two
strands of the double helix run in \emph{opposite} directions
(\emph{antiparallel}, \cref{fig:dna-strands}). By convention, sequences are
always written from $5'$ to $3'$, and polymerases build new strands in the
$5' \to 3'$ direction.

\begin{figure}[htbp]
\centering
\begin{tikzpicture}[x=8mm, y=1cm]
  \def\top{A,T,G,C,G,T,T,A}
  \def\bot{T,A,C,G,C,A,A,T}
  \foreach \b [count=\i] in \top {
    \node[base, fill=base\b] (t\i) at (\i,1.2) {\b};
  }
  \foreach \b [count=\i] in \bot {
    \node[base, fill=base\b] (b\i) at (\i,0) {\b};
  }
  % hydrogen bonds: 2 for AT, 3 for GC
  \foreach \b [count=\i] in \top {
    \ifnum\pdfstrcmp{\b}{A}=0 \def\nb{2}\else\ifnum\pdfstrcmp{\b}{T}=0 \def\nb{2}\else\def\nb{3}\fi\fi
    \ifnum\nb=2
      \draw[gray, densely dashed] ([xshift=-1.2mm]t\i.south) -- ([xshift=-1.2mm]b\i.north);
      \draw[gray, densely dashed] ([xshift=1.2mm]t\i.south) -- ([xshift=1.2mm]b\i.north);
    \else
      \draw[gray, densely dashed] ([xshift=-1.8mm]t\i.south) -- ([xshift=-1.8mm]b\i.north);
      \draw[gray, densely dashed] (t\i.south) -- (b\i.north);
      \draw[gray, densely dashed] ([xshift=1.8mm]t\i.south) -- ([xshift=1.8mm]b\i.north);
    \fi
  }
  \draw[accent, very thick] (0.55,1.55) -- (8.45,1.55);
  \draw[accent, very thick] (0.55,-0.35) -- (8.45,-0.35);
  \node[lbl, anchor=east, font=\small] at (0.4,1.2) {$5'$};
  \node[lbl, anchor=west, font=\small] at (8.6,1.2) {$3'$};
  \node[lbl, anchor=east, font=\small] at (0.4,0) {$3'$};
  \node[lbl, anchor=west, font=\small] at (8.6,0) {$5'$};
  \draw[flow] (9.5,1.2) -- (11,1.2) node[right, lbl, font=\small] {top strand};
  \draw[flow] (11,0) -- (9.5,0) node[pos=0, right, lbl, font=\small] {bottom strand};
  \node[lbl, gray!60!black, align=left, anchor=west] at (0.4,-0.85) {backbone (sugar + phosphate) \quad dashed: hydrogen bonds (A--T: 2, G--C: 3)};
\end{tikzpicture}
\caption{The two strands of DNA are complementary (A pairs with T, G with C)
and antiparallel. The top strand reads \seq{ATGCGTTA} from $5'$ to $3'$; the
bottom strand, read in its own $5' \to 3'$ direction, is \seq{TAACGCAT}.}
\label{fig:dna-strands}
\end{figure}

\begin{exbox}[title={The other strand}]
Given the strand $5'$-\seq{ATGCGTTA}-$3'$:
\begin{itemize}
  \item Its \textbf{complement}, base by base, is \seq{TACGCAAT} --- but that
    is written $3' \to 5'$.
  \item Written in the conventional $5' \to 3'$ direction, the other strand is
    the \textbf{reverse complement}: \seq{TAACGCAT}.
  \item Hydrogen bonds holding this stretch together: 5 A/T pairs $\times 2$ +
    3 G/C pairs $\times 3$ = 19. A stretch of 8 G/C pairs would have 24, a
    stretch of 8 A/T pairs only 16 --- that is the GC-stability argument in
    numbers.
\end{itemize}
Whenever a tool asks whether to search ``both strands'', this is the reason:
a gene can sit on either strand, and on the other strand you see its reverse
complement.
\end{exbox}

\subsection{Packing Two Metres into a Nucleus}
\label{sec:packing}

Two metres of DNA have to fit into a nucleus a few micrometres across. The
cell solves this the way you would stow a long cable: it winds it up, several
times over (\cref{fig:packing}). The DNA is wrapped around ball-shaped
proteins called \textbf{histones}; the resulting ``beads on a string''
(chromatin) are coiled again and again, until --- during cell division ---
they form the compact \textbf{chromosomes}.

\figureplaceholder{4.2cm}{DNA packing}%
{The levels of DNA packing: double helix $\to$ DNA wrapped around histones
(nucleosomes, ``beads on a string'') $\to$ chromatin fibre $\to$ looped
domains $\to$ condensed chromosome. \textit{Search:} ``DNA packaging
histones nucleosome chromatin chromosome diagram''.}%
{How DNA is packed into chromosomes.}{fig:packing}

The histones do more than store the DNA. Their protruding ``tails'' carry
chemical marks that decide whether a region can be unwound and read. These
\textbf{epigenetic} marks let a cell switch genes on or off without changing
the DNA sequence itself --- changing the sequence is hard, changing how
accessible it is, much easier.

\begin{intbox}
If the DNA is the hard drive, epigenetic marks are like \emph{file
permissions}: every cell in your body has the same files, but a liver cell
and a neuron have very different files ``readable''. That is why cells with
identical DNA can be so different (\cref{sec:expression}).
\end{intbox}

\subsection{Chromosomes}

Humans have \textbf{46 chromosomes, in 23 pairs}: one chromosome of each
pair comes from the mother, one from the father. 22 pairs are the same in
both sexes; the 23rd pair determines biological sex (XX: female, XY: male).
Some facts mentioned in the lecture:
\begin{itemize}
  \item The familiar X shape only appears after the DNA has been copied, just
    before cell division (two identical copies joined in the middle). Most of
    the time, chromosomes look more like sausages.
  \item The Y chromosome is much smaller than the X. Since men have only one
    X, they have no ``backup copy'' of the many genes on it, which is why
    some X-linked diseases affect men much more often.
  \item Mitochondria have their own small DNA (mtDNA), which is inherited from
    the mother only --- useful for tracing maternal lineages.
\end{itemize}

\begin{figure}[htbp]
\centering
\begin{tikzpicture}[x=1cm, y=1cm]
  % chromosome as two rounded bars with a waist
  \fill[accent!15, draw=accent, thick, rounded corners=4mm] (0,0) rectangle (3.0,0.9);
  \fill[accent!15, draw=accent, thick, rounded corners=4mm] (3.4,0) rectangle (9.0,0.9);
  \fill[white] (2.9,-0.05) rectangle (3.5,0.95);
  \draw[accent, thick] (2.9,0.15) to[out=0,in=180] (3.2,0.35) to[out=0,in=180] (3.5,0.15);
  \draw[accent, thick] (2.9,0.75) to[out=0,in=180] (3.2,0.55) to[out=0,in=180] (3.5,0.75);
  % bands
  \foreach \x in {0.8,1.4,2.2,4.2,5.0,5.6,6.6,7.4,8.1}
    \fill[accent!45] (\x,0.05) rectangle ++(0.18,0.8);
  % telomeres
  \fill[figgreen!60] (0.02,0.12) rectangle (0.35,0.78);
  \fill[figgreen!60] (8.65,0.12) rectangle (8.98,0.78);
  \node[lbl, font=\small, above] at (1.5,0.95) {p arm (short)};
  \node[lbl, font=\small, above] at (6.2,0.95) {q arm (long)};
  \draw[figred, thick, -{Stealth}] (3.2,-0.6) node[below, font=\small] {centromere} -- (3.2,0.25);
  \draw[figgreen!70!black, thick, -{Stealth}] (0.2,-0.6) node[below, font=\small, align=center] {telomere} -- (0.2,0.15);
  \draw[figgreen!70!black, thick, -{Stealth}] (8.8,-0.6) node[below, font=\small] {telomere} -- (8.8,0.15);
  \node[lbl, gray!60!black, align=center] at (6.0,-0.75) {bands after staining\\(``coordinates'' such as q21.1)};
\end{tikzpicture}
\caption{Anatomy of a chromosome. The \textbf{centromere} is the attachment
point where the cell's machinery pulls the copies apart during division; it
divides the chromosome into the short \textbf{p arm} and the long \textbf{q
arm}. The ends are protected by \textbf{telomeres} (thousands of repeats of
\seq{TTAGGG}). The staining bands give a rough, historical coordinate system.}
\label{fig:chromosome}
\end{figure}

\paragraph{Telomeres and ageing}
When DNA is copied, the copying machinery cannot quite reach the very end of
the chromosome, so each round of replication shortens it a little. The
telomeres are expendable repeats at the ends that absorb this loss. Their
length therefore acts like a clock for how often a cell has divided, which is
why longevity research is very interested in them. Telomeres and centromeres
are mostly structural and regulatory --- and, being repetitive, notoriously
difficult to sequence.

\section{Single Nucleotide Polymorphisms}

No two genomes are identical (except, almost, those of identical twins), yet
about 99.9\,\% of the DNA is the same between any two people. The
differences are concentrated in specific positions.

\begin{defbox}[title={Single Nucleotide Polymorphism (SNP)}]
A \textbf{SNP} (``snip'') is a position in the genome where a single base
varies within a population, and where the less common variant is present in
more than 1\,\% of the population. In humans, SNPs occur roughly every 1000
bases, so each of us carries about 4--5 million of them.
\end{defbox}

SNPs are exactly what consumer DNA tests read. Some have clear effects:
\begin{itemize}
  \item \textbf{Lactose tolerance.} Babies produce the enzyme \emph{lactase},
    which breaks down lactose (milk sugar). Keeping it switched on after
    weaning costs energy, so originally it was switched off. A SNP in an
    intron of the gene \emph{MCM6} keeps it on --- which is why most Central
    Europeans can digest milk as adults.
  \item \textbf{Diseases.} Current research associates many diseases with
    SNPs (e.g.\ schizophrenia, alcohol dependence, Crohn's disease, ALL).
  \item \textbf{Trade-offs.} Variants are rarely purely good or bad. Lighter
    skin (less melanin) makes more vitamin D but sunburns more easily;
    ``double-muscled'' cattle breeds with a disabled muscle-growth brake carry
    a heavy burden on their joints.
\end{itemize}

\section{RNA}

RNA is DNA's short-lived, single-stranded cousin.

\begin{table}[htbp]
\centering
\small
\begin{tabular}{@{}lll@{}}
\toprule
 & \textbf{DNA} & \textbf{RNA} \\
\midrule
full name & deoxyribonucleic acid & ribonucleic acid \\
sugar & deoxyribose (one oxygen less) & ribose \\
bases & A, C, G, \textbf{T} (thymine) & A, C, G, \textbf{U} (uracil) \\
strands & double-stranded & single-stranded \\
stability & very stable (survives for ages) & unstable (degrades quickly) \\
length & long (millions of bases) & short (order of hundreds to thousands) \\
role & \emph{stores} information & \emph{transports} and uses it \\
\bottomrule
\end{tabular}
\caption{DNA versus RNA.}
\label{tab:dna-rna}
\end{table}

The instability of RNA is a feature, not a bug: a message that should only
be read for a while must not last forever. (It was also a practical problem
for mRNA vaccines, whose RNA had to be chemically modified to survive long
enough.) There are several types of RNA:
\begin{itemize}
  \item \textbf{messenger RNA (mRNA)}: the ``classical'' working copy of a
    gene that carries its information to the ribosome;
  \item \textbf{transfer RNA (tRNA)}: small adapters that bring the matching
    amino acid for each triplet of the mRNA (\cref{sec:translation});
  \item \textbf{ribosomal RNA (rRNA)}: part of the ribosome itself, which is a
    mixture of RNA and proteins;
  \item \textbf{regulatory RNA} (e.g.\ miRNA, siRNA): many varieties that
    influence how much of which protein is made, and when.
\end{itemize}

\section{``Junk DNA''? Coding and Non-Coding DNA}

Only a small part of our genome codes for proteins. Early papers called the
rest ``junk DNA''. That turned out to be wrong: much of it is used, just not
for building proteins. \Cref{fig:genome} shows the rough breakdown from the
slides.

\begin{figure}[htbp]
\centering
\begin{tikzpicture}
\begin{axis}[bsa, xbar, width=0.82\linewidth, height=5.2cm,
    xmin=0, xmax=58, xlabel={share of the human genome (\%)},
    ytick={1,...,6}, ymin=0.4, ymax=6.6,
    yticklabels={unknown, regulatory, repeats, introns, legacy code (viruses \ldots), coding (proteins)},
    bar width=9pt, nodes near coords, nodes near coords style={font=\scriptsize},
    grid=none, xmajorgrids]
  \addplot[fill=accent!60, draw=accent] coordinates {
    (2,6) (50,5) (26,4) (10,3) (10,2) (2,1)};
\end{axis}
\end{tikzpicture}
\caption{What our DNA consists of (rough estimates from the lecture). About
50\,\% is ``legacy code'' --- mostly remnants of viruses, largely silenced,
some repurposed for regulation; 26\,\% introns, which make genes modular
(\cref{sec:splicing}); 10\,\% repeats (e.g.\ \seq{TTAGGG} in telomeres); 10\,\%
regulatory elements (promoters as ``start flags'', enhancers as ``volume
dials'', silencers as ``brakes'') and templates for non-coding RNA.}
\label{fig:genome}
\end{figure}

\begin{intbox}
The genome looks a lot like an old, grown codebase: only about 2\,\% is the
actual ``business logic'' (the $\sim$20,000 protein-coding genes), surrounded
by configuration and control flow (regulatory regions), modularisation
(introns), and a lot of legacy code that is commented out but never deleted.
One reason given in the lecture: cutting things out of the genome is risky
and costly, so it is often ``cheaper'' to keep them.
\end{intbox}

DNA is such a dense and stable storage medium that it can even be used as a
hard drive: in 2017, researchers encoded a short film clip in DNA, stored it
in living bacteria, and read it back by sequencing.

\paragraph{A note on scale and speed}
Pictures of cells show lots of empty space between the organelles. In
reality the cell is packed with molecules, all moving extremely fast because
of thermal motion; a molecule visits practically every place in the cell
within seconds. That is what makes the whole machinery work: a protein does
not ``find'' its binding site by navigating, it bumps into it by chance, and
because the bound state is energetically favourable, it stays there.

\section{Chapter Checklist}

The lecturer's tip: write the whole course onto one sheet of paper, by hand
--- it helps to see where each piece fits. For this chapter you should know:
\begin{itemize}
  \item prokaryotes vs.\ eukaryotes (nucleus!), and what the main organelles
    do --- especially nucleus, ribosomes and mitochondria;
  \item the central dogma: replication, transcription, translation, and the
    known exceptions;
  \item DNA: nucleotides (sugar, phosphate, base), the four bases, base
    pairing A--T (2 H-bonds) and G--C (3 H-bonds), GC content, $5' \to 3'$
    direction, antiparallel strands, reverse complement;
  \item packing (histones, chromosomes), 23 chromosome pairs, centromere,
    p/q arms, telomeres, epigenetics in one sentence;
  \item SNPs (definition, frequency, examples);
  \item DNA vs.\ RNA, and the types of RNA;
  \item that only $\sim$2\,\% of the genome codes for proteins, and what
    the rest roughly is.
\end{itemize}

% ==========================================================
%  CHAPTER 3 — BIOLOGY BASICS II: FROM GENES TO PROTEINS
%  Lecture VO02 (12.10., video), slides VO02 p. 1--37
% ==========================================================
\chapter{Biology Basics II: From Genes to Proteins}
\label{ch:bio2}

\lectureref{VO02 (Bio-Intro II), slides 50--86 of the deck. The recording
starts with a few minutes of technical problems; the lecture proper begins
at about 3:45.}

The first biology chapter followed the information from the nucleus to the
RNA. This chapter follows it all the way to the machine: what proteins are
made of, how they fold, and how the cell turns a gene into a protein
(transcription and translation). Along the way we meet the concepts that
the scoring chapters rely on directly: the 20 amino acids and their
properties, codons, and reading frames.

\section{Proteins: The Machines of the Cell}

\begin{defbox}[title={Protein}]
Proteins are large biomolecules made of one or more \textbf{chains of amino
acids} (20 different kinds), linked by \textbf{peptide bonds}. Short chains
are called \emph{peptides} or \emph{polypeptides}. Proteins are the molecular
machines of the cell: they are involved in practically every process.
\end{defbox}

Examples from the lecture show how different their jobs are: the spike
protein of SARS-CoV-2 (the target of the vaccines), the serotonin receptor
(which modulates mood and digestion), \textbf{hemoglobin} (transports oxygen
in the blood), cytokines (modulate immune defences), collagen (gives
structure to tissues), firefly luciferase (emits light), antibodies (identify
invaders) and insulin (regulates blood sugar). Others catalyse chemical
reactions (enzymes such as ligases), give structure (hair), or produce
movement (muscle). A small bacterium contains 50,000 to 1 million protein
molecules, a human cell 1--3 billion; individual proteins live from minutes
to years before they are recycled.

\paragraph{Example: hemoglobin}
Each red blood cell contains several hundred million hemoglobin molecules.
Oxygen binds to the \emph{heme} groups in the hemoglobin, and the red blood
cells carry it through the body. Hemoglobin will be our running example for
sequence comparison in \cref{ch:alignment}.

\section{Amino Acids}
\label{sec:aminoacids}

Every amino acid has the same backbone --- an \emph{amino group}
($\mathrm{NH_2}$), a central carbon atom ($\mathrm{C_\alpha}$), and a
\emph{carboxyl group} ($\mathrm{COOH}$) --- plus a \textbf{side chain}
$R$ that differs from one amino acid to the next (\cref{fig:aminoacid}).
Only the side chain changes along a protein; it gives each amino acid its
chemical character (size, charge, whether it likes water).

\begin{figure}[htbp]
\centering
{\setchemfig{atom sep=2.4em}\Large\chemfig{H_2N-C(-[2]\textcolor{thm@frame}{R})(-[6]H)-C(=[1]O)-[7]OH}}
\qquad\qquad
\begin{tikzpicture}[baseline=-1.4cm]
  \node[font=\small, align=left, anchor=west] at (0,0) {%
    \textcolor{figblue}{amino group}: $\mathrm{NH_2}$\\[2pt]
    \textcolor{figblue}{central carbon}: $\mathrm{C_\alpha}$ (with one H)\\[2pt]
    \textcolor{figblue}{carboxyl group}: $\mathrm{COOH}$\\[2pt]
    \textcolor{figred}{side chain}: $R$ --- the only part that differs};
\end{tikzpicture}
\caption{General structure of an amino acid. Glycine is the smallest one
($R$ is just a hydrogen atom); tryptophan, with its double ring, is one of
the largest.}
\label{fig:aminoacid}
\end{figure}

There are 20 standard amino acids (\cref{tab:aminoacids}), plus two rare ones
(selenocysteine and pyrrolysine) that play a minor role. Learn the
\textbf{one-letter codes}: protein sequences in databases and in every
alignment are written with them. Several are not the first letter of the name
(because that letter was already taken): F = Phe, W = Trp, Y = Tyr, K = Lys,
R = Arg, D = Asp, E = Glu, N = Asn, Q = Gln. \emph{Essential} amino acids are
those the human body cannot make itself and has to get from food.

\begin{table}[htbp]
\centering
\small
\renewcommand{\arraystretch}{1.12}
\begin{tabular}{@{}ll@{\hspace{2.2em}}ll@{\hspace{2.2em}}ll@{\hspace{2.2em}}ll@{}}
\toprule
\multicolumn{2}{@{}l}{\textbf{non-polar}} & \multicolumn{2}{l}{\textbf{polar (uncharged)}} & \multicolumn{2}{l}{\textbf{positive charge}} & \multicolumn{2}{l@{}}{\textbf{negative charge}} \\
\midrule
glycine & Gly G     & serine & Ser S      & lysine & Lys K    & aspartate & Asp D \\
alanine & Ala A     & threonine & Thr T   & arginine & Arg R  & glutamate & Glu E \\
valine & Val V      & tyrosine & Tyr Y    & histidine & His H & & \\
cysteine & Cys C    & asparagine & Asn N  & & & & \\
proline & Pro P     & glutamine & Gln Q   & & & & \\
leucine & Leu L     & & & & & & \\
isoleucine & Ile I  & & & & & & \\
methionine & Met M  & & & & & & \\
tryptophan & Trp W  & & & & & & \\
phenylalanine & Phe F & & & & & & \\
\bottomrule
\end{tabular}
\caption{The 20 standard amino acids with three- and one-letter codes,
grouped by the chemistry of their side chain as on slide 57. Non-polar side
chains are \emph{hydrophobic} (they avoid water); polar and charged ones are
\emph{hydrophilic}.}
\label{tab:aminoacids}
\end{table}

\section{The Peptide Bond}

Amino acids are linked into a chain by a \textbf{peptide bond}: the carboxyl
group of one amino acid reacts with the amino group of the next, and a
molecule of water ($\mathrm{H_2O}$) is released (\cref{fig:peptide}). The
bond is strong (covalent).

\begin{figure}[htbp]
\centering
\begin{tikzpicture}[scale=0.9, transform shape, node distance=3mm,
    aa/.style={circle, draw=accent, thick, fill=accent!8, minimum size=9mm, font=\small}]
  \node[aa] (a1) {Met};
  \node[aa, right=of a1] (a2) {Gly};
  \node[aa, right=of a2] (a3) {Leu};
  \draw[thick, accent] (a1) -- (a2) -- (a3);
  \node[left=1mm of a1, font=\small] {$\mathrm{H_2N}$--};
  \node[right=1mm of a3, font=\small] {--$\mathrm{COOH}$};
  \node[font=\Large] at (5.0,0) {$+$};
  \node[aa] (a4) at (6.0,0) {Ser};
  \draw[flow] (7.0,0) -- (8.2,0);
  \begin{scope}[xshift=9.4cm]
    \node[aa] (b1) {Met};
    \node[aa, right=of b1] (b2) {Gly};
    \node[aa, right=of b2] (b3) {Leu};
    \node[aa, right=of b3] (b4) {Ser};
    \draw[thick, accent] (b1) -- (b2) -- (b3);
    \draw[very thick, figred] (b3) -- (b4);
    \node[font=\scriptsize, figred, below=6mm] at ($(b3)!0.5!(b4)$) {new peptide bond};
  \end{scope}
  \node[font=\small] at (14.6,0) {$+\ \mathrm{H_2O}$};
  \node[font=\scriptsize, gray!60!black] at (1.3,-0.9) {N-terminus \hspace{1.2cm} C-terminus};
\end{tikzpicture}
\caption{Building a protein chain. Each new amino acid is attached to the
free carboxyl end (C-terminus) of the growing chain; each bond releases one
water molecule.}
\label{fig:peptide}
\end{figure}

Two consequences matter for us:
\begin{itemize}
  \item \textbf{Proteins are linear chains.} Each amino acid has exactly one
    amino and one carboxyl group in its backbone, so it can connect to at most
    one predecessor and one successor --- no branching. That is why a protein
    can be written as a \emph{string}.
  \item \textbf{Chains have a direction.} One end has a free amino group, the
    \textbf{N-terminus}; the other a free carboxyl group, the
    \textbf{C-terminus}. Protein sequences are written from N to C, which is
    also the order in which the ribosome builds them.
\end{itemize}
The reverse reaction also exists: water can attack the peptide bond, so a
protein left in water long enough falls apart into its amino acids.

\section{Protein Structure}
\label{sec:protein-structure}

A freshly made amino-acid chain is a floppy string. In the cell it folds
into a specific three-dimensional shape, and \emph{this shape determines what
the protein does}. Biologists distinguish four levels of structure.

\begin{defbox}[title={The four levels of protein structure}]
\begin{itemize}
  \item \textbf{Primary structure}: the amino-acid \emph{sequence} --- the string.
  \item \textbf{Secondary structure}: local, regular folding patterns, held together by
    (weak) hydrogen bonds between backbone atoms: the \textbf{$\alpha$-helix}
    (a coil), the \textbf{$\beta$-sheet} (strands lying side by side), and the
    irregular \textbf{loops} connecting them.
  \item \textbf{Tertiary structure}: how these elements are arranged in 3D to form the
    overall shape of \emph{one} chain.
  \item \textbf{Quaternary structure}: how \emph{several} chains (subunits) assemble
    into a complex. Hemoglobin, for example, consists of four chains.
\end{itemize}
\end{defbox}

\begin{figure}[htbp]
\centering
\begin{tikzpicture}[x=1cm, y=1cm]
  % primary
  \foreach \i in {0,...,7} \draw[accent, fill=accent!20] (0.32*\i,0) circle (0.13);
  \node[font=\small\bfseries, above] at (1.1,0.6) {primary};
  \node[font=\scriptsize, below] at (1.1,-0.3) {\seq{MVHLTPEEK\ldots}};
  % secondary
  \begin{scope}[xshift=3.6cm]
    \draw[accent, thick, domain=0:4*pi, samples=120, variable=\t] plot ({0.08*\t},{0.28*sin(\t r)});
    \foreach \y in {-0.35,-0.05,0.25} \draw[-{Stealth[length=2mm,width=3mm]}, line width=3pt, figgreen!60] (1.25,\y) -- (2.1,\y);
    \node[font=\small\bfseries, above] at (1.0,0.6) {secondary};
    \node[font=\scriptsize, below] at (1.0,-0.5) {$\alpha$-helix \quad $\beta$-sheet};
  \end{scope}
  % tertiary
  \begin{scope}[xshift=7.4cm]
    \draw[accent, thick, fill=accent!10] plot[smooth cycle, tension=0.8] coordinates {(0,0) (0.4,0.5) (1.1,0.45) (1.5,0) (1.2,-0.45) (0.5,-0.4)};
    \draw[accent!70, thick, domain=0:3*pi, samples=80, variable=\t] plot ({0.25+0.07*\t},{0.12*sin(\t r)});
    \node[font=\small\bfseries, above] at (0.75,0.6) {tertiary};
    \node[font=\scriptsize, below] at (0.75,-0.5) {one folded chain};
  \end{scope}
  % quaternary
  \begin{scope}[xshift=10.6cm]
    \draw[accent, thick, fill=accent!10] plot[smooth cycle, tension=0.8] coordinates {(0,0) (0.3,0.45) (0.9,0.4) (1.1,0) (0.8,-0.4) (0.3,-0.35)};
    \draw[figred, thick, fill=figred!10] plot[smooth cycle, tension=0.8] coordinates {(1.0,0.05) (1.3,0.5) (1.9,0.45) (2.1,0.05) (1.8,-0.4) (1.3,-0.35)};
    \node[font=\small\bfseries, above] at (1.05,0.6) {quaternary};
    \node[font=\scriptsize, below] at (1.05,-0.5) {several chains};
  \end{scope}
  % difficulty arrows
  \draw[-{Stealth}, thick, figamber] (1.1,-1.0) to[bend right=18] node[below, font=\scriptsize] {easy to hard} (4.6,-1.0);
  \draw[-{Stealth}, thick, figblue] (1.1,-1.0) to[bend right=24] node[below, font=\scriptsize] {hard to impossible (before AlphaFold)} (8.15,-1.0);
  \draw[-{Stealth}, thick, figred] (8.15,-1.0) to[bend right=30] node[below, font=\scriptsize] {hard} (11.65,-1.0);
\end{tikzpicture}
\caption{The four structure levels (schematic). The arrows show how hard it
is to \emph{predict} each level from the sequence: secondary structure is
comparatively easy; the full 3D (tertiary) structure was ``hard to
impossible'' until AlphaFold; complexes of several proteins are still
hard.}
\label{fig:structure-levels}
\end{figure}

\figureplaceholder{4.2cm}{real protein structure levels}%
{One protein shown at all four levels, e.g.\ hemoglobin: the sequence, an
$\alpha$-helix, one folded subunit, and the assembled four-subunit complex,
in cartoon representation. \textit{Search:} ``four levels of protein
structure primary secondary tertiary quaternary hemoglobin''.}%
{Protein structure levels on a real example.}{fig:structure-real}

\begin{notebox}[title={Exam hint from the lecture}]
When showing the amino acids and the structure levels, the lecturer said:
``please remember this picture \ldots{} it is also important for the exam,
and that's really what you should take with you from the course.'' Be able to
name and explain the four levels, give examples of secondary-structure
elements, and explain that the side chains determine the chemistry.
\end{notebox}

\paragraph{Representations}
A protein can be drawn as a graph of all its atoms --- correct, but almost
impossible to read. Bioinformaticians therefore invented the \emph{cartoon
representation}: helices as spirals or ribbons, sheets as arrows, loops as
thin tubes. Experts read function from such pictures. Other views show the
\emph{solvent-accessible surface}, i.e.\ where a water molecule can touch
the protein. None of these is ``the truth''; even the atom graph is a model
(electrons are really clouds of probability). We pick whichever
representation helps us see what we are looking for. The same is true for
sequences: a nucleotide is a molecule, but for analysing sequences the
abstraction ``a letter'' is all we need.

Experimentally determined structures are collected in the \textbf{Protein Data
Bank (PDB)}, \texttt{rcsb.org}, where each entry has a four-character code
(e.g.\ \seq{1A3N}, human hemoglobin). AlphaFold was trained on this data; it
now predicts the coordinates of all heavy atoms directly (hydrogens are easy
to add afterwards).

\paragraph{Prosthetic groups and ligands}
Myoglobin (which stores oxygen in muscle cells) consists of several
$\alpha$-helices plus a \emph{prosthetic group}: a heme group carrying the
oxygen, which is \emph{not} made of amino acids. Small molecules that bind
to proteins are called \textbf{ligands}. Drug discovery is largely about
finding small molecules that bind a harmful protein and switch it off --- or
that boost a protein the body lacks.

\section{Transcription: From DNA to RNA}
\label{sec:transcription}

\begin{defbox}[title={Transcription}]
An enzyme, \textbf{RNA polymerase}, copies one gene into a messenger RNA,
in three phases:
\begin{enumerate}
  \item \textbf{Initiation}: a complex of proteins places the RNA polymerase
    on the \emph{promoter}, a start signal in front of the gene.
  \item \textbf{Elongation}: the polymerase moves along the DNA, unwinds it,
    builds the RNA copy, and winds the DNA up again behind itself.
  \item \textbf{Termination}: at a stop signal (a stop sequence and/or a
    hairpin loop in the RNA) the polymerase stops and releases the mRNA.
\end{enumerate}
\end{defbox}

Only one of the two DNA strands is read. The polymerase reads the
\textbf{template strand} and builds the RNA complementary to it. The RNA
therefore has the same sequence as the other strand, the \textbf{coding
strand} --- with \seq{U} instead of \seq{T}. The RNA is built in the
$5' \to 3'$ direction. The real process is far more complex (dozens of
helper proteins are involved in initiation alone), but for sequence analysis
this level of detail is enough.

\section{Translation: From RNA to Protein}
\label{sec:translation}

The ribosome reads the mRNA in steps of three nucleotides. Each triplet,
called a \textbf{codon}, stands for one amino acid. The adapter between the
two languages is \textbf{transfer RNA (tRNA)}: each tRNA carries one specific
amino acid on one end and an \emph{anticodon} --- three bases complementary
to a codon --- on the other. Free tRNAs float around in the cell; only the one
whose anticodon matches the current codon can bind. The ribosome then attaches
its amino acid to the growing chain, shifts by one codon, and repeats.

\begin{defbox}[title={The genetic code}]
\begin{itemize}
  \item There are $4^3 = 64$ codons but only 20 amino acids, so most amino
    acids have several codons --- the code is \emph{redundant} (``degenerate'').
    Leucine, serine and arginine have six codons each; methionine and
    tryptophan only one.
  \item \textbf{Start codon:} \seq{AUG}, which also codes for
    \textbf{methionine}. That is why almost every protein starts with
    methionine (M).
  \item \textbf{Stop codons:} \seq{UAA}, \seq{UAG}, \seq{UGA}. They code for no
    amino acid; translation ends there.
\end{itemize}
\end{defbox}

\begin{table}[htbp]
\centering
\small
\renewcommand{\arraystretch}{1.08}
\begin{tabular}{@{}c|llll|c@{}}
\toprule
\textbf{1st} & \multicolumn{4}{c|}{\textbf{2nd base}} & \textbf{3rd} \\
 & \multicolumn{1}{c}{\textbf{U}} & \multicolumn{1}{c}{\textbf{C}} & \multicolumn{1}{c}{\textbf{A}} & \multicolumn{1}{c|}{\textbf{G}} & \\
\midrule
\multirow{4}{*}{\textbf{U}} & \texttt{UUU} Phe & \texttt{UCU} Ser & \texttt{UAU} Tyr & \texttt{UGU} Cys & \textbf{U} \\
 & \texttt{UUC} Phe & \texttt{UCC} Ser & \texttt{UAC} Tyr & \texttt{UGC} Cys & \textbf{C} \\
 & \texttt{UUA} Leu & \texttt{UCA} Ser & \texttt{UAA} \textbf{\color{thm@frame}Stop} & \texttt{UGA} \textbf{\color{thm@frame}Stop} & \textbf{A} \\
 & \texttt{UUG} Leu & \texttt{UCG} Ser & \texttt{UAG} \textbf{\color{thm@frame}Stop} & \texttt{UGG} Trp & \textbf{G} \\
\midrule
\multirow{4}{*}{\textbf{C}} & \texttt{CUU} Leu & \texttt{CCU} Pro & \texttt{CAU} His & \texttt{CGU} Arg & \textbf{U} \\
 & \texttt{CUC} Leu & \texttt{CCC} Pro & \texttt{CAC} His & \texttt{CGC} Arg & \textbf{C} \\
 & \texttt{CUA} Leu & \texttt{CCA} Pro & \texttt{CAA} Gln & \texttt{CGA} Arg & \textbf{A} \\
 & \texttt{CUG} Leu & \texttt{CCG} Pro & \texttt{CAG} Gln & \texttt{CGG} Arg & \textbf{G} \\
\midrule
\multirow{4}{*}{\textbf{A}} & \texttt{AUU} Ile & \texttt{ACU} Thr & \texttt{AAU} Asn & \texttt{AGU} Ser & \textbf{U} \\
 & \texttt{AUC} Ile & \texttt{ACC} Thr & \texttt{AAC} Asn & \texttt{AGC} Ser & \textbf{C} \\
 & \texttt{AUA} Ile & \texttt{ACA} Thr & \texttt{AAA} Lys & \texttt{AGA} Arg & \textbf{A} \\
 & \texttt{AUG} \textbf{\color{ex@frame}Met} & \texttt{ACG} Thr & \texttt{AAG} Lys & \texttt{AGG} Arg & \textbf{G} \\
\midrule
\multirow{4}{*}{\textbf{G}} & \texttt{GUU} Val & \texttt{GCU} Ala & \texttt{GAU} Asp & \texttt{GGU} Gly & \textbf{U} \\
 & \texttt{GUC} Val & \texttt{GCC} Ala & \texttt{GAC} Asp & \texttt{GGC} Gly & \textbf{C} \\
 & \texttt{GUA} Val & \texttt{GCA} Ala & \texttt{GAA} Glu & \texttt{GGA} Gly & \textbf{A} \\
 & \texttt{GUG} Val & \texttt{GCG} Ala & \texttt{GAG} Glu & \texttt{GGG} Gly & \textbf{G} \\
\bottomrule

\end{tabular}
\caption{The standard genetic code (mRNA codons). Read the first base from the
left, the second from the top, the third from the right. Start codon in
green, stop codons in red. You do not need to memorise it, but you should be
able to use it.}
\label{tab:codons}
\end{table}

\begin{exbox}[title={From gene to protein}]
Take the coding strand $5'$-\seq{ATG GCT TGG AAA GAC TAA}-$3'$ (spaces only
for readability).
\begin{enumerate}
  \item \textbf{Template strand} (what the polymerase reads), written
    $3' \to 5'$: \seq{TAC CGA ACC TTT CTG ATT}.
  \item \textbf{mRNA} (complementary to the template, i.e.\ equal to the coding
    strand with U for T): $5'$-\seq{AUG GCU UGG AAA GAC UAA}-$3'$.
  \item \textbf{Translation} with \cref{tab:codons}:
    \seq{AUG}~$\to$~Met, \seq{GCU}~$\to$~Ala, \seq{UGG}~$\to$~Trp,
    \seq{AAA}~$\to$~Lys, \seq{GAC}~$\to$~Asp, \seq{UAA}~$\to$~Stop.
  \item \textbf{Protein:} Met--Ala--Trp--Lys--Asp, in one-letter code
    \seq{MAWKD}, read from N- to C-terminus.
\end{enumerate}
\end{exbox}

\section{Reading Frames}
\label{sec:reading-frames}

If you are given a piece of DNA, you do not know where its codons start.
Shifting the start by one or two positions gives a completely different
protein; shifting by three gives the same one again. So each strand has
\textbf{three reading frames}, and since a gene can be on either strand,
double-stranded DNA has \textbf{six reading frames} in total.

\begin{exbox}[title={Six reading frames of one DNA snippet}]
\small
Strand: \seq{TCATGGCTTGGAAAGACTAAGC}. Its reverse complement:
\seq{GCTTAGTCTTTCCAAGCCATGA}. Translating from position 1, 2 and 3 of each
(\seq{*} = stop):

\medskip
\centering
\begin{tabular}{@{}lll@{}}
\toprule
frame & codons & protein \\
\midrule
$+1$ & \seq{TCA TGG CTT GGA AAG ACT AAG}~(\seq{C}) & \seq{SWLGKTK} \\
$+2$ & \seq{CAT GGC TTG GAA AGA CTA AGC} & \seq{HGLERLS} \\
$+3$ & \seq{ATG GCT TGG AAA GAC TAA}~(\seq{GC}) & \seq{MAWKD*} \\
$-1$ & \seq{GCT TAG TCT TTC CAA GCC ATG}~(\seq{A}) & \seq{A*SFQAM} \\
$-2$ & \seq{CTT AGT CTT TCC AAG CCA TGA} & \seq{LSLSKP*} \\
$-3$ & \seq{TTA GTC TTT CCA AGC CAT}~(\seq{GA}) & \seq{LVFPSH} \\
\bottomrule
\end{tabular}

\raggedright
Only frame $+3$ starts with \seq{ATG} (Met) and runs until a stop codon --- an
\emph{open reading frame}, the classic signature of a gene. It is the gene
from the previous example, hidden in a longer piece of DNA.
\end{exbox}

\begin{intbox}
Reading frames are like decoding a byte stream when you have lost the byte
boundaries: start one bit too late and every ``byte'' you read is garbage.
That is why a single inserted or deleted base inside a gene (a
\emph{frameshift}) usually destroys the whole protein downstream, while a
single substituted base changes at most one amino acid --- and often none,
thanks to the redundancy of the code.
\end{intbox}

\section{Protein Folding}

Why does a chain fold into one particular shape? Many effects play a role;
the lecture highlights one of the most important ones,
\textbf{hydrophobicity}. Hydrophobic (``water-fearing'', non-polar) side
chains tend to collapse towards the interior of the protein, where they are
shielded from the surrounding water; hydrophilic (polar and charged) side
chains face outward, towards the water. It is the same effect that makes oil
form droplets in water and lets lipids form cell membranes.

\begin{figure}[htbp]
\centering
\begin{tikzpicture}
\begin{axis}[bsa, width=0.62\linewidth, height=5.0cm, axis lines=left,
    xmin=0, xmax=10, ymin=0, ymax=10.5, xtick=\empty, ytick=\empty,
    xlabel={possible shapes (conformations)}, ylabel={energy}, grid=none,
    samples=200, domain=0.2:9.8]
  \addplot[very thick, accent] {9.5 - 7.5*exp(-((x-5.2)^2)/4.5) - 1.2*exp(-((x-3.0)^2)/0.3) - 0.8*exp(-((x-7.4)^2)/0.25) + 0.5*sin(deg(3*x))};
  \node[font=\scriptsize, anchor=south] at (axis cs:1.0,9.6) {unfolded};
  \node[font=\scriptsize, anchor=west, align=left] at (axis cs:6.9,6.2) {intermediate\\states};
  \node[font=\scriptsize, anchor=north] at (axis cs:5.2,1.3) {native state (lowest energy)};
  \draw[-{Stealth}, thick, figamber] (axis cs:1.6,9.0) to[bend right=20] (axis cs:4.6,2.6);
\end{axis}
\end{tikzpicture}
\caption{Folding as energy minimisation (schematic). The unfolded chain is
high in energy; folding moves it ``downhill'' through intermediate states
into the native state of lowest energy, in which the protein does its job.}
\label{fig:folding}
\end{figure}

All folding activity tends to drive the molecule into a state of lowest
energy (\cref{fig:folding}) --- folding is nature's version of an optimisation
algorithm. Predicting the result of this optimisation from the sequence is
exactly the problem AlphaFold solved.

\section{Splicing}
\label{sec:splicing}

In eukaryotes, genes are not continuous. They consist of \textbf{exons}
(the parts that end up in the protein) interrupted by \textbf{introns}
(parts that are cut out). After transcription, the introns are removed from
the RNA and the exons are glued together: this is \textbf{splicing}. It is
done by the \textbf{spliceosome}, a machine of about 60 proteins and pieces
of RNA that is assembled and disassembled for every single splicing event ---
and that has to work perfectly, otherwise the resulting protein is destroyed.

\begin{figure}[htbp]
\centering
\begin{tikzpicture}[x=1cm, y=1cm,
    ex/.style={draw, thick, minimum height=5mm, inner sep=0pt},
    intr/.style={thick, gray}]
  % pre-mRNA
  \node[font=\small, anchor=east] at (-0.2,0) {pre-mRNA};
  \node[ex, fill=figred!45, minimum width=1.2cm] (e1) at (0.6,0) {1};
  \draw[intr] (1.2,0) -- (1.8,0);
  \node[ex, fill=figamber!50, minimum width=1.2cm] (e2) at (2.4,0) {2};
  \draw[intr] (3.0,0) -- (3.6,0);
  \node[ex, fill=figgreen!45, minimum width=1.2cm] (e3) at (4.2,0) {3};
  \draw[intr] (4.8,0) -- (5.4,0);
  \node[ex, fill=figblue!40, minimum width=1.2cm] (e4) at (6.0,0) {4};
  \node[font=\scriptsize, gray] at (1.5,0.35) {intron};
  % isoform A
  \node[font=\small, anchor=east] at (-0.2,-1.3) {mRNA A};
  \node[ex, fill=figred!45, minimum width=1.2cm] at (0.6,-1.3) {1};
  \node[ex, fill=figamber!50, minimum width=1.2cm] at (1.8,-1.3) {2};
  \node[ex, fill=figgreen!45, minimum width=1.2cm] at (3.0,-1.3) {3};
  \node[ex, fill=figblue!40, minimum width=1.2cm] at (4.2,-1.3) {4};
  \node[font=\small, anchor=west] at (5.1,-1.3) {$\to$ protein isoform A};
  % isoform B
  \node[font=\small, anchor=east] at (-0.2,-2.2) {mRNA B};
  \node[ex, fill=figred!45, minimum width=1.2cm] at (0.6,-2.2) {1};
  \node[ex, fill=figgreen!45, minimum width=1.2cm] at (1.8,-2.2) {3};
  \node[ex, fill=figblue!40, minimum width=1.2cm] at (3.0,-2.2) {4};
  \node[font=\small, anchor=west] at (5.1,-2.2) {$\to$ protein isoform B};
  \draw[flow] (3.2,-0.4) -- (3.2,-0.9);
  \node[font=\scriptsize, anchor=west] at (3.3,-0.62) {splicing};
\end{tikzpicture}
\caption{Splicing removes the introns. With \textbf{alternative splicing}, the
same gene yields different mRNAs --- here, one with and one without exon 2 ---
and therefore different proteins.}
\label{fig:splicing}
\end{figure}

\begin{defbox}[title={Alternative splicing}]
Different combinations of exons from one gene lead to different proteins
(\emph{isoforms}). This contradicts the early idea from the times of Watson
and Crick of ``one gene --- one protein'', and it is the main reason why we
have many more proteins than genes. The current record holder is the
\emph{DSCAM} gene of the fruit fly \emph{Drosophila}, which can be spliced
into about 38,000 different mRNAs.
\end{defbox}

How the cell decides which exons to keep is still an active research topic;
introns are thought to play regulatory roles. In software terms, introns make
a gene \emph{modular}: the exons are reusable building blocks that can be
combined in different ways.

\section{Replication: From DNA to DNA}

Before a cell divides, its complete genome has to be copied, so that each
daughter cell gets one copy. An enzyme called \textbf{helicase} unwinds the
double helix, and \textbf{DNA polymerase} builds a new partner strand for
\emph{each} of the two old strands, using base pairing as the template. So,
unlike in transcription, two strands are produced, and each new double helix
consists of one old and one new strand. Helicase is fast: the slides compare
it with a jet engine, both rotating at around 10,000 revolutions per minute.

\section{Sequencing and the Human Genome Project}

\textbf{DNA sequencing} determines the order of the bases in a DNA molecule:
cells are sampled (for example from the liver), the DNA is extracted and
purified, and a sequencing machine reads it in many short pieces, which are
then put together computationally (genome assembly, see
\cref{sec:more-alignment}).

The \textbf{Human Genome Project} sequenced the whole human genome. It was
launched in October 1990, completed in April 2003, cost about 3 billion
dollars, and was funded by the US National Institutes of Health (NIH). The
result does not represent a single individual: it is a mosaic assembled from
several people. Since then, the cost of sequencing has dropped dramatically
--- from about 100 million dollars per genome in 2001 to roughly 1,000
dollars in 2022, much faster than Moore's law would predict.

\section{Mutations, Disease and Gene Expression}
\label{sec:expression}

A change of a single letter can have dramatic consequences. The lecture's
example is the \textbf{BRCA1} and \textbf{BRCA2} genes: more than 60\,\% of
women who inherit a harmful change in one of them develop breast cancer during
their lifetime, compared with about 13\,\% in the general population. For men,
7--26\,\% (BRCA1) and 19--61\,\% (BRCA2) develop prostate cancer by age 80,
compared with about 10.6\,\% overall. Mutations can also be beneficial ---
lactose tolerance (\cref{ch:bio1}) is an example.

Note also that every cell of your body contains the same DNA, but \emph{not
every gene is expressed in every cell}. A lung cell and a liver cell
produce very different amounts of specific proteins, and that is what makes
them different.

\section{Outlook: What This Course Is About}
\label{sec:alphafold-msa}

The rest of the course is about \textbf{comparing sequences}. A typical
picture is a \emph{multiple sequence alignment}: the same protein from
mouse, rat, human, pig, dog, chicken and frog written one below the other, so
that corresponding positions are in the same column. Some columns are
identical in all species --- these positions are \emph{conserved}, often over
hundreds of millions of years, because they are essential for the protein's
function. Other columns vary.

Such alignments are not just a classical tool. AlphaFold's success relies on
them: for a new sequence, it first searches databases for similar sequences
and builds a multiple sequence alignment, which enriches the information it
feeds into its neural network. Since this database search is expensive, a
current trend is to replace it with \emph{protein language models} --- neural
networks trained on huge amounts of sequence data, whose internal
representations are used instead. The institute behind this course
contributes to this line of work, for example with Bio-xLSTM, trained on DNA
and protein sequences (more in \cref{ch:ai}).

\section{Chapter Checklist}

\begin{itemize}
  \item what proteins are and examples of what they do;
  \item amino-acid structure (backbone + side chain), the 20 amino acids,
    their one-letter codes and chemical groups (non-polar, polar, charged);
  \item the peptide bond, why proteins are linear chains, N- and C-terminus;
  \item the four levels of protein structure ($\alpha$-helix, $\beta$-sheet,
    loops) and why structure determines function --- \emph{exam-relevant};
  \item transcription (initiation, elongation, termination; template vs.\
    coding strand) and translation (ribosome, tRNA, codons, start and stop
    codons, redundancy of the code); translate a short sequence by hand;
  \item reading frames (three per strand, six in total);
  \item folding by energy minimisation; hydrophobic in, hydrophilic out;
  \item exons, introns, (alternative) splicing;
  \item replication, sequencing, Human Genome Project in a few numbers.
\end{itemize}

% ==========================================================
%  CHAPTER 4 — ALIGNMENTS, SIMILARITY AND EVOLUTION
%  Lecture VO03 (19.10., A. Regl), slides 1--16, recording 0:00--27:00
% ==========================================================
\chapter{Alignments, Similarity and Evolution}
\label{ch:alignment}

\lectureref{VO03 (Scoring, part 1, Alois Regl), slides 1--16. Recording
0:00--27:00.}

With the biology in place, we can now ask the central question of the
course: \emph{how similar are two sequences, and what does that tell us?}
This chapter explains what an alignment is, why similarity between sequences
is evidence for evolutionary relationship, and where that reasoning reaches
its limits. The next chapters then turn ``similarity'' into numbers.

The scoring part of the course is given by Alois Regl, who taught the whole
course for many years. His outline: first why we need alignments and a
scoring system, then scoring for DNA (``the easiest of all \ldots{} very
primitive''), then scoring for proteins with the two classical methods PAM
and BLOSUM. In practice, he says, 99.99\,\% of all protein alignments use one
of those two. This lays the groundwork for the alignment algorithms later in
the course (\cref{ch:global,ch:local,ch:msa}).

\section{What Is an Alignment?}

Everybody has an intuition for it: write two sequences one above the other
and shift them, inserting gaps where necessary, until as many positions as
possible agree.

\alnblock{%
THIS~IS~---~AN~AL-IGNMENT\\
||~~~||~~~~~||~|~~|||||||\\
THAT~IS~NOT~AN~ASSIGNMENT}

\begin{defbox}[title={Pairwise alignment}]
An alignment of two sequences writes them one above the other, possibly with
\textbf{gap} characters (\seq{-}) inserted, so that both have the same length.
Each column then is one of:
\begin{itemize}
  \item a \textbf{match}: the same character in both sequences;
  \item a \textbf{mismatch}: two different characters;
  \item a \textbf{gap}: a character aligned to ``nothing'' --- a placeholder
    where one sequence has no corresponding character.
\end{itemize}
A column with two gaps is never allowed.
\end{defbox}

\section{How Sequences Change: Mutations}

Alignments are interesting because sequences change over time through
\textbf{mutations}. The slides illustrate the types with a word:

\begin{table}[htbp]
\centering
\small
\begin{tabular}{@{}llll@{}}
\toprule
\textbf{Mutation} & \textbf{Example} & \textbf{What happens} & \textbf{In an alignment} \\
\midrule
none          & \seq{BEAST} & --- & matches \\
substitution  & \seq{BEAST} $\to$ \seq{\mm{F}EAST}  & one letter replaced by another & mismatch \\
insertion     & \seq{BEAST} $\to$ \seq{B\mm{R}EAST} & a letter is added & gap in the other sequence \\
deletion      & \seq{BEAST} $\to$ \seq{BEST}  & a letter is removed & gap in this sequence \\
inversion     & \seq{BEAST} $\to$ \seq{BEA\mm{TS}} & a piece is reversed & not covered \\
\bottomrule
\end{tabular}
\caption{Types of mutations. Insertions and deletions are collectively called
\textbf{indels} (they show up as gaps). Inversions --- for example a whole
gene that appears elsewhere in reverse order --- are common in biology, but
standard alignments do not model them.}
\label{tab:mutations}
\end{table}

\section{Why Align? Similarity Means Relationship}

The underlying idea is simple. All living things --- us, animals, trees,
bacteria --- go back to a common ancestor (or a few). Every genome on Earth
evolved from that one by small mutations that accumulated over generations.
When a new species arises, it inherits the genome of its predecessor and
continues to accumulate its own mutations. Mutations that are passed down
the germline are the basis of evolution from the first organism to the
diversity we see today.

\begin{thmbox}[title={Similarity as evidence for relationship}]
By comparing sequences that serve the same purpose in different organisms
(typically genes or proteins), we can infer evolutionary relationships:
\textbf{the more similar two sequences, the ``younger'' their last common
ancestor}.
\end{thmbox}

The rate at which mutations accumulate is called the \textbf{molecular
clock}, and the simplest assumption is that it ``ticks'' at a constant speed.
Rough estimates from the slides: in mammals about 1 mutation per $10^9$ base
pairs per year; mitochondrial DNA about 100 times faster; viruses $10^3$ to
$10^5$ times faster; bacteria somewhere in between.

\begin{notebox}[title={The clock is not constant}]
The constant-rate assumption is a simplification and does not really hold:
there were periods with much higher mutation rates, for example the Cambrian
explosion, in which many new species appeared. For our purposes it is a
useful working assumption --- but keep it in mind when interpreting
``how long ago'' from sequence differences.
\end{notebox}

Comparing sequences is not only about evolution. To find the gene responsible
for a disease, one of the first steps is to align the sequences of healthy and
ill people and look for the differences.

\section{Example: Hemoglobin Across Species}

Hemoglobin is an ideal example: practically all animals have it, but with
slight variations accumulated during evolution. Hemoglobin from an Antarctic
fish (PDB entry \seq{1PBX}) and from humans (\seq{1A3N}) look almost the same
at first sight, but on closer inspection there are many small differences ---
and their number tells us how long ago fish and humans shared an ancestor.

The lecture compares the $\beta$ subunit of hemoglobin (gene \emph{HBB}) of
three organisms: human, chimpanzee (\emph{Pan troglodytes}) and chicken. What
do we expect? High similarity between human and chimpanzee (their common
ancestor lived only a few million years ago --- nothing compared with the
billions of years of life on Earth), and lower similarity between them and the
chicken. The sequences come in the standard \textbf{FASTA} format:

\begin{defbox}[title={FASTA format}]
A text format for sequences. A header line starts with \seq{>} and contains an
identifier and a description; the following lines contain the sequence (for
proteins in one-letter code):

\smallskip
{\ttfamily\footnotesize
>sp|P68871|HBB\_HUMAN Hemoglobin beta OS=Homo sapiens OX=9606 GN=HBB\\
MVHLTPEEKSAVTALWGKVNVDEVGGEALGRLLVVYPWTQRFFESFGDLSTPDAVMGNPK\\
VKAHGKKVLGAFSDGLAHLDNLKGTFATLSELHCDKLHVDPENFRLLGNVLVCVLAHHFG\\
KEFTPPVQAAYQKVVAGVANALAHKYH}
\end{defbox}

\begin{exbox}[title={Human, chimpanzee and chicken hemoglobin}]
All three $\beta$ chains have 147 amino acids. Comparing them position by
position (no gaps are needed here):
\begin{itemize}
  \item human vs.\ chimpanzee: \textbf{identical}, all 147 positions agree;
  \item human vs.\ chicken: 102 of 147 positions agree (69\,\%), 45 differ.
\end{itemize}
The first 30 positions of human and chicken:

\alnblock{%
MVHLTPEEKSAVTALWGKVNVDEVGGEALG\\
|||~|~|||~~~|~|||||||~|~|~|||~\\
MVHWTAEEKQLITGLWGKVNVAECGAEALA}

The multiple alignment tool \emph{Clustal Omega} (\cref{ch:msa}) turns such
comparisons into a tree. It puts human and chimpanzee at distance 0 from each
other and gives the chicken a branch of length 0.153 --- half of the 30.6\,\%
difference, because in a tree the distance between two leaves is split along
the path between them.
\end{exbox}

\paragraph{Why do all three start with M}
The lecturer asked this in the lecture, and nobody knew. The start codon
\seq{AUG} codes for methionine (\cref{sec:translation}), so every protein is
\emph{made} starting with M. Many proteins lose this first methionine
afterwards (\emph{post-translational modification}), but these three still
have it.

\section{What Does ``98\,\% of the Genome'' Mean?}

Headlines like ``Just 2.5\,\% of DNA turns mice into men'' (New Scientist,
2019) suggest that you could align the DNA of mice and humans and find only
2.5\,\% of the letters different. That is \emph{not} what such numbers mean.
Even seemingly unrelated organisms show striking similarity, simply because
cells and their proteins have to work in the same or a similar way in all of
them --- but the percentages are defined differently.

\begin{defbox}[title={Shared genome (the most widely accepted definition)}]
``A and B share $x\,\%$ of their genome'' means: $x\,\%$ of the \textbf{genes}
of A have a \textbf{homolog} in B. A homolog is a gene that is highly similar
(in this context typically 80--90\,\% similar, not identical) because it
derives from the same ancestral gene.
\end{defbox}

Three things follow:
\begin{enumerate}
  \item it is about \textbf{genes only} --- the non-coding regions (98\,\% of
    the DNA, \cref{fig:genome}) may differ much more;
  \item the homologs themselves are far from identical;
  \item the definition is \textbf{not symmetric} (\cref{fig:homologs}).
\end{enumerate}

\begin{figure}[htbp]
\centering
\begin{tikzpicture}[x=1cm, y=1cm]
  \node[font=\small\bfseries] at (-0.5,1.2) {A};
  \node[font=\small\bfseries] at (-0.5,0) {B};
  \draw[line width=3pt, figamber!80] (0,1.2) -- (2.2,1.2);
  \draw[line width=3pt, figgreen!60] (2.4,1.2) -- (4.4,1.2);
  \draw[line width=3pt, figblue!60] (4.6,1.2) -- (6.6,1.2);
  \draw[line width=3pt, figamber!80] (0,0) -- (2.2,0);
  \draw[line width=3pt, figblue!60] (3.4,0) -- (5.4,0);
  \draw[densely dashed, gray] (1.1,1.1) -- (1.1,0.1);
  \draw[densely dashed, gray] (5.6,1.1) -- (4.4,0.1);
  \node[font=\scriptsize, anchor=west] at (6.9,1.2) {3 genes};
  \node[font=\scriptsize, anchor=west] at (6.9,0) {2 genes};
  \node[font=\small, anchor=north, align=left] at (3.6,-0.45) {$2/3 = 67\,\%$ of A's genes have a homolog in B\\$2/2 = 100\,\%$ of B's genes have a homolog in A};
\end{tikzpicture}
\caption{Why ``shared genome'' is not symmetric. Organism A has three genes,
B has two; same colour = homologs. The percentage depends on whose genes
you count --- so always say what the 100\,\% are.}
\label{fig:homologs}
\end{figure}

\section{Sequence Is Not Everything}

Are the proteins PCNA (in eukaryotes) and the $\beta$ subunit of DNA
polymerase III (in bacteria) related? Both have the same job: they form a ring
that clamps around the DNA and keeps the DNA polymerase working while it copies.
Judged by sequence, the answer is a clear ``no'': after alignment they share
only about 10\,\% identical positions --- about what two random sequences
also reach. But their 3D structures (PDB \seq{1AXC} and \seq{1MMI}) look
almost the same: two different ways of building the same six-domain ring
(the bacterial clamp from two subunits with three domains each, PCNA from
three subunits with two domains each).

\begin{figure}[htbp]
\centering
\begin{tikzpicture}[node distance=9mm]
  \node[box, fill=figblue!10] (s1) {sequence 1};
  \node[box, fill=figblue!10, below=3mm of s1] (s2) {sequence 2};
  \node[box, fill=figgreen!10, right=12mm of $(s1)!0.5!(s2)$] (st) {structure};
  \node[box, fill=figamber!15, right=of st] (fn) {function};
  \draw[flow] (s1.east) -- (st.west);
  \draw[flow] (s2.east) -- (st.west);
  \draw[flow] (st) -- (fn);
\end{tikzpicture}
\caption{Proteins do their work based on \emph{structure}, not sequence. The
sequence determines the structure, but very different sequences can fold into
the same structure --- and then have the same function.}
\label{fig:seq-struct-fn}
\end{figure}

\figureplaceholder{4.0cm}{PCNA and the bacterial sliding clamp}%
{Side by side: the E.~coli DNA polymerase III $\beta$ clamp (PDB 1MMI) and
human/yeast PCNA (PDB 1AXC), both seen from the top as a ring around the DNA
--- two unrelated sequences with the same shape. \textit{Search:} ``PCNA beta
clamp structure comparison'' or the RCSB PDB pages of 1AXC and 1MMI.}%
{Same shape and job, almost no sequence similarity.}{fig:pcna}

\begin{notebox}[title={Be careful with sequence similarity}]
Sequence similarity is only part of the story. When you study evolution or
function, do not rely on sequence alignment alone --- look at the 3D
structure as well. Today that is easy, thanks to millions of predicted
structures (AlphaFold). How to \emph{quantify} the similarity of two
structures is shown in \cref{sec:chimerax}.
\end{notebox}

So there are two kinds of similarity: \textbf{sequence similarity} (compare the
strings) and \textbf{structural (3D) similarity}. This course covers the
first; structural similarity is the topic of structural bioinformatics.

\section{Chapter Checklist}

\begin{itemize}
  \item alignment terminology: match, mismatch, gap; mutation types
    (substitution, insertion, deletion = indel; inversion);
  \item the logic ``more similar $\Rightarrow$ more recent common ancestor'',
    the molecular clock, and why its constant rate is only an assumption;
  \item the FASTA format; why proteins start with methionine;
  \item what ``A and B share $x\,\%$ of their genome'' means (homologs; genes
    only; not symmetric);
  \item sequence $\to$ structure $\to$ function, and the PCNA example.
\end{itemize}

% ==========================================================
%  CHAPTER 5 — SCORING ALIGNMENTS: THE BASICS AND DNA
%  Lecture VO03 (A. Regl), slides 15--33, recording 25:00--71:00
% ==========================================================
\chapter{Scoring Alignments: The Basics and DNA}
\label{ch:scoring}

\lectureref{VO03 (Scoring, part 1), slides 15--33. Recording 25:00--71:00.}

To find the ``best'' alignment, we must first say what ``best'' means. This
chapter introduces scores for matches, mismatches and gaps, shows how they add
up to the score of an alignment, and why trying all alignments is hopeless.
It then looks at how DNA is scored in practice --- with a critical eye ---
and prepares the ground for the much more careful protein scores of the next
chapters.

\section{Which Alignment Is Better?}

Here are two ways to align \seq{ATGGCCAGC} and \seq{AGGCCAGCA}:

\noindent
\begin{minipage}[t]{0.48\linewidth}
\alnblock{%
ATGGCCAGC\\
|~|~|~~~~\\
AGGCCAGCA}
\centering\small without gaps: 3 matches, 6 mismatches
\end{minipage}\hfill
\begin{minipage}[t]{0.48\linewidth}
\alnblock{%
ATGGCCAGC-\\
|~|~|||||~\\
AGG-CCAGCA}
\centering\small with gaps: 7 matches, 1 mismatch, 2 gaps
\end{minipage}
\medskip

Which one is better?\footnote{The slide counts 2 matches and 7 mismatches for
the left alignment; position 5 (\seq{C}/\seq{C}) is also a match.} The right
one has many more matches --- and matches are good, because they mean the
sequence has not changed. But it needs gaps, and mutations that insert or
delete letters are \emph{rare} events compared with substitutions. Is one gap
worth more matches and fewer mismatches? We cannot decide without defining
how much each event is worth.

\begin{thmbox}[title={The idea of scoring}]
Give each column of an alignment a score: a \textbf{reward} for a match, a
\textbf{penalty} for a mismatch, and a (usually larger) penalty for a gap. The
score of the alignment is the sum over all columns, and the best alignment is
the one with the highest total score. The scores should \textbf{imitate
nature}: mutations that happen often in nature get a high (good) score,
mutations that happen rarely a low (bad) one.
\end{thmbox}

\begin{exbox}[title={The answer depends on the scores}]
With match $= +1$, mismatch $= -2$ and gap $= -3$:
\begin{itemize}
  \item left: $3 \cdot (+1) + 6 \cdot (-2) = -9$;
  \item right: $7 \cdot (+1) + 1 \cdot (-2) + 2 \cdot (-3) = -1$ $\;\Rightarrow$
    the alignment with gaps wins.
\end{itemize}
Make gaps much more expensive, say $-10$, and the right alignment drops to
$7 - 2 - 20 = -15$, worse than the left one. The scores encode our beliefs
about biology, and they decide which alignment is ``optimal''. (With the
first set of scores, the best possible alignment of these two sequences
scores $+2$: \seq{ATGGCCAGC-} over \seq{A-GGCCAGCA}, with 8 matches and 2
gaps.)
\end{exbox}

\paragraph{Reward or penalty?}
Some authors speak of rewards, others of penalties. They mean the same thing
with opposite signs: a penalty of 3 is a reward of $-3$. In practice, matches
usually get positive scores and mismatches and gaps negative ones, which is
where the word ``gap penalty'' comes from. Just be careful with the sign.

\section{Scoring Matrices}

Three numbers (match, mismatch, gap) are not enough in general: how likely a
substitution is depends on \emph{which} letters are involved. So we use a
\textbf{scoring matrix} $S$ that defines the score $S(a,b)$ for every pair
of letters $a, b$ (bases or amino acids).

The simplest approach ignores which letters are involved:
\[
  S(a,b) = \begin{cases} x & \text{if } a = b, \\ y & \text{if } a \ne b, \end{cases}
\]
with $x$ typically positive and $y$ zero or negative. The very simplest idea
just counts matches ($x = 1$, $y = 0$); a harder penalty for mismatches would
be, for example, $x = 1$, $y = -3$. For DNA, such heuristics are used in
practice; for proteins there are much more clever schemes (PAM and BLOSUM).

\subsection{From Columns to Alignments to the Optimum}

The scores are combined in three levels:
\begin{enumerate}
  \item a \textbf{single column} gets the score $S(a_i, b_i)$ of the two letters
    involved (or a gap penalty);
  \item a \textbf{whole alignment} gets the sum of its column scores,
    $\sum_{i} S(a_i, b_i)$;
  \item the \textbf{optimal alignment} is the one with the best score among all
    possible alignments $k$:
    \[
      \max_{k} \sum_{i} S\bigl(a_{i,k},\, b_{i,k}\bigr).
    \]
\end{enumerate}
The last line is ``what we are looking for''. Finding it efficiently is the
job of the alignment algorithms in \cref{ch:global,ch:local}.

\section{Two Small Examples}

\subsection{Shifting Without Gaps}

Are \seq{ATCGA} and \seq{GTTCG} similar? Without gaps, all we can do is shift
one sequence along the other and score each position. The lecture uses
match $= +1$, mismatch $= -2$, and --- implicitly --- charges $-1$ for every
letter that sticks out without a partner.

\begin{exbox}[title={All shifts of \seq{GTTCG} against \seq{ATCGA}}]
\centering
\small
\begin{tabular}{@{}lccc@{}}
\toprule
alignment & overlap & score of the overlap & total (overhang $-1$ each) \\
\midrule
\seq{\phantom{GTTC}ATCGA} / \seq{GTTCG\phantom{ATCG}}   & 1 & $-2$ & $-10$ \\
\seq{\phantom{GTT}ATCGA} / \seq{GTTCG\phantom{ATC}}     & 2 & $-4$ & $-10$ \\
\seq{\phantom{GT}ATCGA} / \seq{GTTCG\phantom{AT}}       & 3 & $-6$ & $-10$ \\
\seq{\phantom{G}ATCGA} / \seq{GTTCG\phantom{A}}         & 4 & $+1$ & $\mathbf{-1}$ \\
\seq{ATCGA} / \seq{GTTCG}                               & 5 & $-7$ & $-7$ \\
\seq{ATCGA\phantom{G}} / \seq{\phantom{A}GTTCG}         & 4 & $-8$ & $-10$ \\
\seq{ATCGA\phantom{GT}} / \seq{\phantom{AT}GTTCG}       & 3 & $-6$ & $-10$ \\
\seq{ATCGA\phantom{GTT}} / \seq{\phantom{ATC}GTTCG}     & 2 & $-1$ & $-7$ \\
\seq{ATCGA\phantom{GTTC}} / \seq{\phantom{ATCG}GTTCG}   & 1 & $-2$ & $-10$ \\
\bottomrule
\end{tabular}

\raggedright
\medskip
The best shift aligns \seq{TCG} with \seq{TCG}:

\alnblock{%
\phantom{G}ATCGA\\
\phantom{G}~|||~\\
GTTCG\phantom{A}}

Overlap: $-2 + 1 + 1 + 1 = +1$; minus 1 for each of the two overhanging
letters gives $-1$. Without gaps, there are only $5 + 5 - 1 = 9$ shifts to
try --- easy.
\end{exbox}

\subsection{Allowing Gaps}

Now compare \seq{ATCGA} and \seq{ACGA}, and allow insertions and deletions in
both sequences, with match $= +1$, mismatch $= -2$, gap $= -3$:

\begin{exbox}[title={Some alignments of \seq{ATCGA} and \seq{ACGA}}]
\centering
\small
\begin{tabular}{@{}lll@{}}
\toprule
alignment & columns & score \\
\midrule
\seq{ATCGA} / \seq{A-CGA} & 4 matches, 1 gap                  & $4 - 3 = \mathbf{+1}$ \\
\seq{ATCGA} / \seq{-ACGA} & 3 matches, 1 mismatch, 1 gap      & $3 - 2 - 3 = -2$ \\
\seq{ATCGA} / \seq{AC-GA} & 3 matches, 1 mismatch, 1 gap      & $-2$ \\
\seq{ATCGA} / \seq{ACG-A} & 2 matches, 2 mismatches, 1 gap    & $2 - 4 - 3 = -5$ \\
\seq{ATCGA} / \seq{ACGA-} & 1 match, 3 mismatches, 1 gap      & $1 - 6 - 3 = -8$ \\
\bottomrule
\end{tabular}

\raggedright
\medskip
The best of these is \seq{ATCGA} over \seq{A-CGA}: delete the \seq{T}, and
everything else matches. A dynamic-programming search over \emph{all}
alignments confirms that $+1$ is the optimum.\footnote{The slide lists $-7$ for
\seq{ACGA-}; with these scores it is $-8$.}
\end{exbox}

\section{How Many Alignments Are There?}

With gaps allowed anywhere in both sequences, the number of possible
alignments explodes. For two sequences of length $n$, the lecture counts
\[
  \binom{2n}{n} = \frac{(2n)!}{(n!)^2} \;\approx\; \frac{2^{2n}}{\sqrt{\pi n}}
\]
alignments (the approximation follows from Stirling's formula).

\begin{table}[htbp]
\centering
\small
\begin{tabular}{@{}rll@{}}
\toprule
$n$ & $\binom{2n}{n}$ (exact) & approximation $2^{2n}/\sqrt{\pi n}$ \\
\midrule
5    & 252                       & 258 \\
10   & 184{,}756                 & 186{,}968 \\
20   & $1.38 \times 10^{11}$     & $1.39 \times 10^{11}$ \\
100  & $9.05 \times 10^{58}$     & $9.07 \times 10^{58}$ \\
1000 & $2.05 \times 10^{600}$    & $2.05 \times 10^{600}$ \\
\bottomrule
\end{tabular}
\caption{The number of alignments grows exponentially. For comparison: the
whole universe contains about $10^{80}$ atoms.}
\label{tab:num-alignments}
\end{table}

For $n = 5$ you could write a program that tries all 252 alignments. For two
proteins of length 1000 --- a perfectly ordinary size --- there are about
$10^{600}$; even checking $10^{582}$ alignments per second would take
$10^{18}$ seconds, about twice the age of the universe. And sequences can be
thousands or millions of letters long.

\begin{notebox}[title={Brute force is out of the question}]
Computing all alignments to pick the best one is completely out of reach. We
need clever algorithms --- dynamic programming, which finds the optimum
without enumerating all alignments. That is the topic of
\cref{ch:global,ch:local}.
\end{notebox}

\section{Scoring for DNA}

For DNA, practitioners use simple schemes: three numbers for match, mismatch
and gap, ignoring which bases are involved and any biological reasoning. You
can see this in \textbf{BLAST} (\cref{sec:blast}), \emph{the} tool for
sequence similarity searches: under ``scoring parameters'' you choose a
match/mismatch pair such as $1, -2$ and the gap costs. Written as a matrix,
match/mismatch $= 1/{-2}$ looks like this:

\begin{center}
\small
\begin{tabular}{@{}c|rrrr@{}}
 & A & C & G & T \\
\midrule
A & 1 & $-2$ & $-2$ & $-2$ \\
C & $-2$ & 1 & $-2$ & $-2$ \\
G & $-2$ & $-2$ & 1 & $-2$ \\
T & $-2$ & $-2$ & $-2$ & 1 \\
\end{tabular}
\end{center}

Alois Regl calls this ``in my opinion, an unforgivable oversimplification''
--- the numbers come ``out of the gut'', without statistical or scientific
reasoning, and it would be interesting to see how the zillions of past
alignments would change with a better scheme.

\subsection{Transitions and Transversions}
\label{sec:transitions}

A first step towards a better world comes from chemistry. The four bases come
in two shapes:
\begin{itemize}
  \item \textbf{purines}: \seq{A} and \seq{G} --- larger, two-ring structure;
  \item \textbf{pyrimidines}: \seq{C} and \seq{T} --- smaller, single ring.
\end{itemize}

\begin{defbox}[title={Transition and transversion}]
A \textbf{transition} is a substitution within the same class (purine $\to$
purine, i.e.\ \seq{A} $\leftrightarrow$ \seq{G}, or pyrimidine $\to$
pyrimidine, \seq{C} $\leftrightarrow$ \seq{T}). A \textbf{transversion}
changes the class (purine $\leftrightarrow$ pyrimidine).
\end{defbox}

Structurally, replacing a base by one of the same shape disturbs the DNA much
less: \seq{A} is a better replacement for \seq{G} than for \seq{C} or
\seq{T}. And the data agree (\cref{fig:transitions}).

\begin{figure}[htbp]
\centering
\begin{tikzpicture}[x=1cm, y=1cm,
    nb/.style={circle, draw=accent, thick, minimum size=10mm, font=\large\bfseries}]
  \node[nb, fill=figblue!15] (A) at (0,2.6) {A};
  \node[nb, fill=figred!12] (C) at (4.2,2.6) {C};
  \node[nb, fill=figblue!15] (G) at (0,0) {G};
  \node[nb, fill=figred!12] (T) at (4.2,0) {T};
  \draw[very thick, figgreen, <->] (A) -- node[left, font=\small] {27.2\,\%} (G);
  \draw[very thick, figgreen, <->] (C) -- node[right, font=\small] {35.4\,\%} (T);
  \draw[thick, gray!70, <->] (A) -- node[above, font=\small] {9.4\,\%} (C);
  \draw[thick, gray!70, <->] (G) -- node[below, font=\small] {9.2\,\%} (T);
  \draw[thick, gray!70, <->] (A) -- node[pos=0.28, above right=-1pt, font=\small] {10.4\,\%} (T);
  \draw[thick, gray!70, <->] (G) -- node[pos=0.28, below right=-1pt, font=\small] {8.4\,\%} (C);
  \node[font=\small, figblue, align=center] at (-1.6,1.3) {purines\\(two rings)};
  \node[font=\small, figred, align=center] at (5.9,1.3) {pyrimidines\\(one ring)};
  \node[font=\small, anchor=west, align=left] at (7.0,1.8) {\textcolor{figgreen!80!black}{\textbf{transitions}} (vertical):\\ 62.6\,\% of all substitutions};
  \node[font=\small, anchor=west, align=left] at (7.0,0.8) {\textcolor{gray!60!black}{\textbf{transversions}} (others):\\ 37.4\,\%};
\end{tikzpicture}
\caption{Observed substitutions between the four bases in genes of the malaria
mosquito \emph{Anopheles} (data from the slides, Wondji et al.\ 2007). There
are two kinds of transitions but four kinds of transversions; per kind,
transitions are about $31.3 / 9.35 \approx 3.3$ times as frequent.}
\label{fig:transitions}
\end{figure}

So transitions are overrepresented by a factor of about 3, and transversions
should get a lower score (a higher penalty). Some DNA matrices take this into
account; others, including the standard BLAST settings, do not:

\begin{center}
\small
\begin{tabular}{@{}c|rrrr@{\hspace{2.5em}}c|rrrr@{\hspace{2.5em}}c|rrrr@{}}
 & A & C & G & T &  & A & C & G & T &  & A & C & G & T \\
\cmidrule(r{2.5em}){1-5}\cmidrule(r{2.5em}){6-10}\cmidrule{11-15}
A & 1 & $-5$ & $-1$ & $-5$ & A & 1 & $-3$ & $-3$ & $-3$ & A & 2 & $-3$ & $-3$ & $-3$ \\
C & $-5$ & 1 & $-5$ & $-1$ & C & $-3$ & 1 & $-3$ & $-3$ & C & $-3$ & 2 & $-3$ & $-3$ \\
G & $-1$ & $-5$ & 1 & $-5$ & G & $-3$ & $-3$ & 1 & $-3$ & G & $-3$ & $-3$ & 2 & $-3$ \\
T & $-5$ & $-1$ & $-5$ & 1 & T & $-3$ & $-3$ & $-3$ & 1 & T & $-3$ & $-3$ & $-3$ & 2 \\
\multicolumn{5}{c}{\footnotesize transitions $-1$, transversions $-5$} &
\multicolumn{10}{c}{\footnotesize both used by BLAST (no transition/transversion distinction)}
\end{tabular}
\end{center}

Even the left matrix only reflects a ``gut feeling'' about how much worse a
transversion is; there is no underlying theory behind its entries. (The
BLAST matrices do have some theoretical background, explained with protein
scoring in the next chapters --- interestingly, they ignore
transitions/transversions.)

\subsection{NCBI's Rule of Thumb}

The one guideline that exists comes from NCBI\footnote{The \emph{National
Center for Biotechnology Information}, part of the US National Institutes of
Health; it runs BLAST and many central databases. Its European counterpart is
the EMBL (European Molecular Biology Laboratory) with its bioinformatics
institute EMBL-EBI.}: \textbf{the more similar your sequences are, the higher
the ratio between the match and the mismatch score should be.}

\begin{center}
\small
\begin{tabular}{@{}lccc@{}}
\toprule
expected similarity & 99\,\% & 95\,\% & 75\,\% \\
match/mismatch & $1/{-3}$ & $1/{-2}$ & $1/{-1}$ \\
ratio & 0.33 & 0.50 & 1.00 \\
\bottomrule
\end{tabular}
\end{center}

\begin{intbox}
Why harsher mismatch penalties for \emph{similar} sequences? If two sequences
are expected to be 99\,\% identical (say, human and chimpanzee), mismatches
are rare and surprising --- an alignment that needs many of them is probably
the wrong alignment, so it should be punished. For distantly related
sequences (75\,\%), mismatches are everyday events and should not hurt much.
The same idea returns for proteins as the choice between PAM and BLOSUM
matrices for different evolutionary distances (\cref{sec:choose-matrix}).
\end{intbox}

\paragraph{A question from the lecture: are these matrices always symmetric}
In practice yes: $S(a,b) = S(b,a)$. In principle one could penalise
$\seq{A} \to \seq{C}$ differently from $\seq{C} \to \seq{A}$, but when we look
at two sequences we cannot tell which one was the original (\cref{sec:pam}).

\subsection{What About RNA?}

Why does nobody align RNA sequences in the same way? Because RNA is normally
just a copy of DNA: aligning the RNA gives the same result as aligning the
corresponding DNA. What is interesting about RNA is its \emph{structure}: most
RNA types (not mRNA) fold back onto themselves, because some stretches are
complementary to other stretches of the same molecule and pair up into
hairpins and loops. These structures influence transcription and translation,
so in RNA analysis one looks for regions of \emph{self-complementarity}.

\section{Scoring for Proteins: The Idea}

For proteins, a two-number scheme would be far too crude. The 20 amino
acids differ a lot: glycine is the smallest (its side chain is a single
hydrogen atom), while arginine and tryptophan are large (tryptophan is the
only one with a double ring). Some substitutions are therefore easy for a
protein to tolerate and others are not.

\paragraph{Geometric (conformational) considerations}
Mutations occur less frequently if they disturb the overall geometry of the
protein --- the same argument as for transitions and transversions. Valine
and threonine have side chains of almost the same shape: one can replace the
other without moving the chain. Replacing a small amino acid (glycine) by a
large one (tryptophan) means the chain has to make room for several positions
to the left and right --- unlikely. The opposite direction, large to small, is
easier.

\paragraph{Chemical considerations}
But structure is not everything. Valine is \emph{non-polar} (hydrophobic) and
threonine is \emph{polar} (hydrophilic) --- in Regl's words, ``like God and
the devil''. Polar amino acids sit on the outside of the protein, where the
water is; non-polar ones hide inside. So swapping valine for threonine causes
no conformational stress but a \emph{chemical} stress, and is unlikely. Changes
that tend to preserve the chemical properties of the residue, on the other
hand, are common and should get a low penalty:
\begin{itemize}
  \item non-polar $\leftrightarrow$ non-polar and polar $\leftrightarrow$ polar;
  \item aspartate $\leftrightarrow$ glutamate (both negatively charged);
  \item alanine $\leftrightarrow$ valine, leucine $\leftrightarrow$ isoleucine
    (similar hydrophobic side chains);
  \item plus a few others.
\end{itemize}

\subsection{Who Can Build a Substitution Matrix?}

For proteins we need a $20 \times 20$ matrix with a score for every pair of
amino acids --- 400 entries (210 different ones, since the matrix is
symmetric). Two ways to get it:
\begin{itemize}
  \item An informed \textbf{chemist or biologist} could set up a reasonable
    matrix by looking at the chemical and biological background, giving high
    penalties to changes that cause structural or chemical difficulties.
    But that means reasoning about 400 numbers by hand.
  \item As \textbf{bioinformaticians}, we can simply look at what nature did.
    Our databases contain millions of protein sequences; aligning related ones
    and counting which substitutions actually occurred gives the scores by
    statistics. The basic idea: \emph{if a mutation is rarely seen in nature,
    it must be unfavourable.}
\end{itemize}

\begin{thmbox}[title={Recipe for a protein scoring matrix}]
\begin{enumerate}
  \item Take a \textbf{careful selection of sequences} from one or several
    databases (not simply ``all of them'': mixing bacteria, viruses and
    humans leads to nothing, and the same protein appearing in many databases
    would be overrepresented).
  \item Build alignments and \textbf{count} how often each pair of amino acids
    is aligned (the pairwise mutation rates).
  \item \textbf{Convert} the counts into a statistical score --- for PAM and
    BLOSUM, a \emph{log-odds} score.
\end{enumerate}
The two established implementations of this recipe are \textbf{PAM} (old,
from 1974, but still widely used) and \textbf{BLOSUM} (younger, better data,
more sophisticated).
\end{thmbox}

\section{Chapter Checklist}

\begin{itemize}
  \item why we need scores; reward vs.\ penalty; scores should imitate nature;
  \item scoring matrix $S(a,b)$; score of an alignment = sum of column scores;
    optimal alignment = maximum over all alignments;
  \item score the small examples by hand (shifts without gaps; alignments
    with gaps);
  \item the number of alignments ($\binom{2n}{n} \approx 4^n/\sqrt{\pi n}$) and
    why brute force is impossible;
  \item DNA scoring in practice (BLAST match/mismatch); purines vs.\
    pyrimidines; transitions vs.\ transversions (factor $\approx 3$); NCBI's
    rule of thumb;
  \item why protein scoring needs a $20 \times 20$ matrix: geometric and
    chemical reasons; the three-step recipe.
\end{itemize}

% ==========================================================
%  CHAPTER 6 — LOG-ODDS SCORES
%  Lecture VO03 end (71:00--76:00) and VO04 (31:00--50:00)
%  Slides: VO04_scoring-2.pdf p. 1--6
% ==========================================================
\chapter{Log-Odds Scores: Turning Counts into Scores}
\label{ch:logodds}

\lectureref{End of VO03 (from 71:00) and VO04 (31:00--50:00), slides 34--39
of the scoring deck.}

The recipe from the last chapter ends with ``convert the counts into a
statistical score''. Both PAM and BLOSUM do this in the same way: with
\emph{log-odds} scores. This chapter explains what that means. Regl called
these slides the most important ones of his part of the course --- once you
understand the score of a single pair of amino acids, PAM and BLOSUM are just
two ways of collecting the data.

\section{Odds and Log-Odds}

Consider an experiment with only two possible outcomes: success or failure,
heads or tails.

\begin{defbox}[title={Odds and log-odds}]
For an event $A$ with probability $p(A)$, the \textbf{odds} are the ratio of
the probability of success to the probability of failure,
\[
  \text{odds}(A) = \frac{p(A)}{p(\neg A)} = \frac{p(A)}{1 - p(A)},
\]
and the \textbf{log-odds} are simply their logarithm,
$\log\bigl(p(A) / (1 - p(A))\bigr)$.
\end{defbox}

\begin{exbox}[title={Throwing a six}]
Let $A$ = ``throwing a 6'' with one die. Then $p(A) = 1/6$ and
$p(\neg A) = 5/6$ (throwing 1 to 5). The odds are
\[
  \frac{1/6}{5/6} = \frac{1}{5} = 0.2,
\]
and the log-odds are $\log_{10} 0.2 \approx -0.70$ (or $\ln 0.2 \approx -1.61$,
or $\log_2 0.2 \approx -2.32$ --- the base only changes the scale).
\end{exbox}

\section{The Score of a Pair of Amino Acids}

For a substitution matrix we want, for every pair of amino acids $(a,b)$, a
number that says: is it \emph{more or less likely than chance} to see $a$ and
$b$ aligned with each other in related proteins?

\begin{thmbox}[title={Log-odds score (the common core of PAM and BLOSUM)}]
\[
  \mathrm{Score}(a,b) = \mathrm{round}\!\left( \text{factor} \times
    \log \frac{\obs}{\expd} \right),
\]
where
\begin{itemize}
  \item \textbf{observed} is the relative frequency with which $a$ and $b$ are
    found aligned with each other in our set of related sequences;
  \item \textbf{expected} is the probability of seeing them aligned \emph{by
    chance}: $p(a) \times p(b)$, where $p(x)$ is the relative frequency of
    amino acid $x$ in the sequences.
\end{itemize}
The factor and the rounding only make the numbers easier to use; the base of
the logarithm does not matter (the original papers used $\log_2$, today
$\log_{10}$ is also common).
\end{thmbox}

Why $p(a) \times p(b)$? If the two positions had nothing to do with each
other, they would behave like two independent dice: the probability of
throwing two fours is $\frac16 \times \frac16$. Comparing what we
\emph{observe} with what independent chance would give is exactly the right
question: a pair that appears more often than chance must be favoured by
evolution.

\begin{notebox}[title={Strictly speaking, not ``odds''}]
Regl points out that observed/expected is a \emph{likelihood ratio}, not a
ratio of success to failure. Everybody calls it ``log-odds'' anyway, and so do
we. He is also not happy with the rounding: matrix entries are small integers
(mostly between about $-5$ and $+10$), so rounding 2.6 to 3 or 1.4 to 1 is a
substantial change --- but everybody does it.
\end{notebox}

\paragraph{PAM and BLOSUM differ only in the data}
Both use this formula. They differ in \emph{which} sequences they count:
\begin{itemize}
  \item \textbf{PAM} counts mutations along the family trees of closely related
    proteins (\cref{ch:pam});
  \item \textbf{BLOSUM} counts substitutions in short, highly conserved
    blocks without gaps, independent of family trees (\cref{ch:blosum}).
\end{itemize}

\section{A Complete Example}
\label{sec:logodds-example}

This is the example Regl called ``for sure the most important one'' of his
lectures. Our ``database'' is a tiny multiple alignment of four sequences over
a three-letter alphabet:

\alnblock{%
ABC\\
ABB\\
CBA\\
ABA}

\paragraph{Step 1: expected}
There are $4 \times 3 = 12$ characters: 5 A, 5 B, 2 C. So
\[
  p(A) = \tfrac{5}{12} \approx 0.417, \qquad
  p(B) = \tfrac{5}{12} \approx 0.417, \qquad
  p(C) = \tfrac{2}{12} \approx 0.167 .
\]
(The slides round these to 0.41, 0.41 and 0.17.) The expected frequency of a
pair is the product, e.g.\ $p(A)\,p(C) \approx 0.069$ and
$p(A)\,p(B) \approx 0.174$.

\paragraph{Step 2: observed}
In each column, compare every sequence with every sequence \emph{below} it
(not above --- otherwise every pair is counted twice), and count the pairs,
including the ``non-mutations'' such as A--A. Order does not matter: A--C and
C--A are the same pair.
\begin{itemize}
  \item column 1 (A, A, C, A): AA, AC, AA, AC, AA, CA $\;\Rightarrow\;$ 3 AA, 3 AC;
  \item column 2 (B, B, B, B): 6 BB;
  \item column 3 (C, B, A, A): CB, CA, CA, BA, BA, AA $\;\Rightarrow\;$ 1 BC,
    2 AC, 2 AB, 1 AA.
\end{itemize}
With $N$ sequences, each column contributes $\binom{N}{2} = N(N-1)/2$ pairs ---
here $4 \cdot 3 / 2 = 6$ --- so the total is $3 \times 6 = 18$ pairs.

\paragraph{Step 3: compare}

\begin{center}
\small
\begin{tabular}{@{}lrrrrr@{}}
\toprule
pair & count & observed & expected $p(a)p(b)$ & obs/exp & $\log_2$(obs/exp) \\
\midrule
AA & 4 & $4/18 = 0.222$ & 0.174 & 1.28 & $+0.36$ \\
AB & 2 & $2/18 = 0.111$ & 0.174 & 0.64 & $-0.64$ \\
AC & 5 & $5/18 = 0.278$ & 0.069 & 4.00 & $+2.00$ \\
BB & 6 & $6/18 = 0.333$ & 0.174 & 1.92 & $+0.94$ \\
BC & 1 & $1/18 = 0.056$ & 0.069 & 0.80 & $-0.32$ \\
CC & 0 & $0$            & 0.028 & 0    & $-\infty$ \\
\bottomrule
\end{tabular}
\end{center}

In this ``database'', A and C are aligned four times as often as chance
would predict: a strongly favoured substitution (score $+2$ bits). A and B are
aligned less often than chance: an unfavourable one. And C--C never occurs,
which gives $-\infty$ --- a sign that 4 sequences of length 3 are far too
little data.

\begin{notebox}[title={How much data is enough}]
A question from the lecture: two sequences would be ``crazy''; is there a
minimum? Regl knows no formal rule; his gut feeling is at least 50--100,
better 200 sequences --- and in general, use as many as you can, the software
does not care whether it processes 10 or 10,000. Real matrices also never
contain $-\infty$, because they are built from far more data. Note too that
BLOSUM computes the expected frequency of an unordered pair $a \ne b$ as
$2\,p(a)\,p(b)$ (the pair can arise as $ab$ or as $ba$), which lowers the
off-diagonal scores by one bit (\cref{sec:blosum-final}).
\end{notebox}

\section{Reading a Scoring Matrix}

Ignoring the factor and rounding, a score is $\log(\obs/\expd)$. Every
logarithm, whatever its base, is zero at 1, negative below 1 and positive
above 1 (\cref{fig:log}). So the \emph{sign} of a score tells you everything:

\begin{thmbox}[title={Interpreting a score}]
\begin{itemize}
  \item $\mathrm{Score}(a,b) > 0$: $a \leftrightarrow b$ is observed
    \textbf{more often} than by chance --- an acceptable, favoured
    substitution;
  \item $\mathrm{Score}(a,b) = 0$: observed \textbf{as often} as by chance;
  \item $\mathrm{Score}(a,b) < 0$: observed \textbf{less often} than by chance ---
    evolution resists this substitution.
\end{itemize}
\end{thmbox}

\begin{figure}[htbp]
\centering
\begin{tikzpicture}
\begin{axis}[bsa, width=0.55\linewidth, height=4.6cm, axis lines=middle,
    xmin=0, xmax=6.2, ymin=-3.2, ymax=2.8, samples=150, domain=0.12:6,
    xlabel={obs/exp}, xlabel style={at={(axis description cs:1,0.5)}, anchor=west},
    legend style={font=\scriptsize, draw=none, fill=white, fill opacity=0.85,
                  at={(0.98,0.03)}, anchor=south east}]
  \addplot[very thick, accent] {log2(x)}; \addlegendentry{$\log_2$}
  \addplot[very thick, figred] {ln(x)};   \addlegendentry{$\ln$}
  \addplot[very thick, figamber] {log10(x)}; \addlegendentry{$\log_{10}$}
  \fill[figgreen] (axis cs:1,0) circle (2pt);
  \node[font=\scriptsize, anchor=south east] at (axis cs:1,0.05) {1};
\end{axis}
\end{tikzpicture}
\caption{All logarithms are 0 at obs/exp $= 1$: the sign of a log-odds score
does not depend on the base.}
\label{fig:log}
\end{figure}

\needspace{16\baselineskip}
\begin{exbox}[title={Reading real scores}]
A few entries of a real matrix (BLOSUM62):

\begin{center}
\small
\begin{tabular}{@{}c|rrrr@{}}
 & A & C & D & E \\
\midrule
A & 4 & 0 & $-2$ & $-1$ \\
C & 0 & 9 & $-3$ & $-4$ \\
D & $-2$ & $-3$ & 6 & 2 \\
\end{tabular}
\end{center}

\begin{itemize}
  \item A $\to$ D: $-2$, less likely than by chance;
  \item D $\to$ E: $+2$, more likely than by chance (both are negatively
    charged --- a conservative substitution);
  \item A $\to$ C: 0, as likely as by chance;
  \item the diagonal is positive (an amino acid ``mutating'' into itself is
    the most common event), and C--C is particularly high: cysteine is rare
    and rarely replaced.
\end{itemize}
\end{exbox}

\paragraph{What time span does a matrix represent}
A student asked whether ``D $\to$ E is more likely than by chance'' refers to
some unit of time. Answer: it depends on the sequences used to build the
matrix. Built from human and fish proteins, whose last common ancestor lived
perhaps 200 million years ago, the matrix reflects that time span. PAM is
explicitly designed for a defined time span (\cref{ch:pam}); the standard
BLOSUM matrices were built from sequences from all over the databases, so they
represent a mixture of time spans.

\section{From a Score Back to a Probability}

The formula can be inverted. Ignoring factor and rounding, with
$\log_{10}$:
\[
  \mathrm{Score} = \log_{10}\frac{\obs}{p_a\, p_b}
  \quad\Longleftrightarrow\quad
  \obs = 10^{\mathrm{Score}} \times p_a \times p_b .
\]

\begin{exbox}[title={Score 3}]
With $\mathrm{Score} = 3$ and $p_a = p_b = 0.02$:
$\obs = 10^3 \times 0.02 \times 0.02 = 0.4$. (This is the slide's illustration
without scaling. In a real matrix the factor matters: BLOSUM uses
$2\log_2$, so there a score of 3 means obs/exp $= 2^{3/2} \approx 2.8$, not
$10^3$.)
\end{exbox}

\section{Chapter Checklist}

\begin{itemize}
  \item odds and log-odds; the dice example;
  \item the score formula $\mathrm{round}(\text{factor} \times
    \log(\obs/\expd))$, what ``observed'' and ``expected'' mean, and why
    expected $= p(a)\,p(b)$;
  \item compute observed and expected frequencies and the scores for a small
    multiple alignment by hand (\cref{sec:logodds-example}) ---
    \emph{the most important slide of this part};
  \item read a matrix: positive, zero, negative scores;
  \item invert a score into an observed frequency.
\end{itemize}

% ==========================================================
%  CHAPTER 7 — PAM MATRICES
%  Lecture VO04 (A. Regl), slides 40--63 of the deck (VO04_scoring-2 p. 7--29)
%  Recording 49:00--64:00 (many derivation slides were skipped in the lecture)
% ==========================================================
\chapter{PAM Matrices}
\label{ch:pam}

\lectureref{VO04 (Scoring, part 2), slides 40--63 of the deck. Recording
49:00--64:00. In the lecture, Regl skipped most of the derivation slides
(``you should believe me''); this chapter goes through them step by step.}

PAM was the first substitution matrix for proteins built from data, and it
is still in use half a century later. The idea: look at closely related
proteins, count which amino acids were replaced by which during evolution,
normalise the counts to a fixed amount of evolutionary change, and turn them
into log-odds scores.

\section{History and Name}
\label{sec:pam}

PAM goes back to \textbf{Margaret Dayhoff} (1925--1983), Georgetown University
Medical Center. Her activities spanned far beyond bioinformatics --- planetary
atmospheres, chemical bonding, phylogeny, even neural networks and pattern
recognition. She founded the \emph{Protein Information Resource} (PIR), one
of the first protein sequence databases.\footnote{The slide calls it ``the
precursor of today's PDB''. More precisely, PIR is a \emph{sequence} database
that later became part of UniProt; the PDB, the \emph{structure} database, was
founded independently in 1971.} In the 1970s (the slides say 1974) she
developed the substitution matrix for proteins that is still in use today.

\begin{defbox}[title={PAM = ``Point Accepted Mutation''}]
The name is not even proper English; ``accepted point mutation'' would be
more accurate, but PAM is easier to say. (For a while, ``percentage of accepted
mutations'' was also used, which comes close.)
\begin{itemize}
  \item \textbf{Point mutation}: one amino acid is replaced by another --- no
    insertions or deletions, no genomic rearrangements.
  \item \textbf{Accepted}: accepted by evolution, i.e.\ the mutation stays in the
    lineage of descendants. Many mutations die out after a generation or
    two; only those passed on to all following generations count.
\end{itemize}
\end{defbox}

\section{The Basic Idea}

``Look at evolutionary events and derive the log-odds for amino-acid
substitutions from them.'' The important twist: the mutation probabilities
are scaled so that they represent a fixed \emph{amount of evolution},
namely \textbf{one accepted mutation per 100 amino acids} (1\,\% change).
This unit is called \textbf{1 PAM}. Regl translates it, very roughly, into
about one million years --- a gross oversimplification, since the speed of
evolution is not constant, but it gives an idea.

Dayhoff's starting point: \textbf{1572 observed mutations} in evolutionary
trees of \textbf{71 families of closely related proteins} (more than 85\,\%
sequence identity).

\subsection{Why Closely Related Sequences?}

\begin{figure}[htbp]
\centering
\begin{tikzpicture}[x=1cm, y=1cm,
    st/.style={circle, draw=accent, thick, fill=accent!8, minimum size=7mm, font=\small\bfseries}]
  \draw[-{Stealth}, thick, gray] (-4.3,3.8) -- (-4.3,-0.2) node[below, font=\small] {time};
  \node[st] (t) at (1.2,3.6) {T};
  \node[st] (a) at (1.2,2.4) {A};
  \node[st] (i) at (1.2,1.2) {I};
  \node[st] (l) at (1.2,0.0) {L};
  \draw[flow] (t) -- (a); \draw[flow] (a) -- (i); \draw[flow] (i) -- (l);
  \draw[decorate, decoration={brace, amplitude=4pt}, thick, figgreen] (1.8,3.8) -- (1.8,2.2)
    node[midway, right=5pt, font=\small, align=left, figgreen] {closely related:\\1 mutation, and it is visible};
  \draw[decorate, decoration={brace, amplitude=6pt, mirror}, thick, figred] (0.5,3.8) -- (0.5,-0.2)
    node[midway, left=7pt, font=\small, align=right, figred] {remote relationship:\\T vs.\ L looks like 1 mutation,\\in reality there were 3};
\end{tikzpicture}
\caption{One position over time. Compared with each other, two sequences
that are far apart show only the net change T $\to$ L and hide the
intermediate steps. Closely related sequences make sure that each observed
difference is a single event.}
\label{fig:multiple-hits}
\end{figure}

We want to \emph{see} mutations, but we want to avoid \emph{unseen} ones. The
selection of sequences is a balance: not too closely related (some mutations
must have happened, otherwise the counts are all zero), and not too remotely
related (otherwise multiple mutations at the same position are miscounted as
one, \cref{fig:multiple-hits}).

\section{From Proteins to Counts}

\subsection{Counting Mutations Along a Tree}

For each protein family, Dayhoff built an evolutionary tree and counted, for
every branch, which amino acids were replaced by which.

\begin{exbox}[title={Three sequences, one tree}]
Take the three related sequences \seq{IAGCW}, \seq{IAGCT} and \seq{IIGCT}.
Among all trees that could explain them, choose the one that needs the
\emph{minimum number of mutations} (\cref{fig:pam-tree}): \seq{IIGCT} differs
from the other two in position 2 (A $\leftrightarrow$ I), and \seq{IAGCW}
differs from \seq{IAGCT} in position 5 (W $\leftrightarrow$ T). So the count
matrix $A$ gets one entry for A--I and one for W--T; all other entries stay
zero. Repeating this for all branches of all trees fills $A$.
\end{exbox}

\begin{figure}[htbp]
\centering
\begin{tikzpicture}[level distance=11mm,
    level 1/.style={sibling distance=36mm},
    level 2/.style={sibling distance=24mm},
    every node/.style={font=\small},
    mut/.style={draw=accent, ellipse, fill=accent!8, font=\small},
    leaf/.style={draw=accent, rounded corners=2pt, fill=figgreen!10, font=\small\ttfamily}]
  \node[mut] {A $\leftrightarrow$ I}
    child {node[mut] {W $\leftrightarrow$ T}
      child {node[leaf] {IAGCW}}
      child {node[leaf] {IAGCT}}}
    child {node[leaf] {IIGCT}};
\end{tikzpicture}
\caption{The most parsimonious tree for the three sequences: two mutations
explain all differences. The arrows are two-headed on purpose --- from the
sequences alone we cannot tell the direction of a mutation.}
\label{fig:pam-tree}
\end{figure}

\subsection{Making the Counts Symmetric}

Mutation counts are inherently \emph{not} symmetric: A $\to$ B and B $\to$ A
can (and will) occur with different frequencies. But when we look at two
sequences, say \seq{CCGTCAG} and \seq{CCGCCAG}, we cannot tell which one was
the original --- was it T $\to$ C or C $\to$ T? As the best guess, the count
matrix is made symmetric:
\[
  A_{ij} = A_{ji} = \frac{A_{ij} + A_{ji}}{2},
  \qquad\text{e.g.\ } \frac{8 + 7}{2} = 7.5 .
\]

\subsection{Frequencies and Mutabilities}

\begin{defbox}[title={Quantities from the counts}]
\begin{itemize}
  \item $n(j)$: total number of occurrences of amino acid $j$ in the data;
    $N = \sum_j n(j)$: total number of amino acids.
  \item $f(j) = n(j)/N$: the \textbf{frequency} of amino acid $j$ (sums to 1).
  \item $m(j) = \dfrac{\sum_{i \ne j} A(i,j)}{n(j)}$: the \textbf{mutability}
    of $j$ --- how often it mutated, relative to how often it occurs.
\end{itemize}
\end{defbox}

\section{The Mutation Probability Matrix}

Now we turn counts into probabilities. Let $M(a,b)$ be the probability that
amino acid $a$ is replaced by $b$ within the time span of 1 PAM.\footnote{The
slides write the same formulas with the roles of the indices exchanged
($M(i,j)$ as the probability that $j$ mutates into $i$, so that columns sum
to one). The convention used here --- rows sum to one --- matches the
Markov-chain view in \cref{sec:markov}; the resulting scores are the same.}

\begin{thmbox}[title={Mutation matrix $M$}]
\[
  M(a,b) = \lambda\,\frac{A(a,b)}{n(a)} \quad (a \ne b),
  \qquad
  M(a,a) = 1 - \sum_{b \ne a} M(a,b) = 1 - \lambda\, m(a).
\]
The diagonal has to be calculated, because ``non-mutations'' were never
counted: if an amino acid does not mutate, it stays the same, so each row of
$M$ sums to 1.
\end{thmbox}

\paragraph{Choosing $\lambda$}
$\lambda$ is a constant of proportionality that fixes the \emph{amount} of
evolution $M$ represents. We want 1 mutation per 100 amino acids: on average,
99\,\% of the positions stay the same and 1\,\% change. The expected fraction of
changing positions is $\sum_a f(a)\,\bigl(1 - M(a,a)\bigr) = \lambda \sum_a f(a)\,m(a)$,
so
\[
  0.99 = 1 - \lambda \sum_{j=1}^{20} f(j)\,m(j)
  \quad\Longrightarrow\quad
  \lambda = \frac{0.01}{\sum_{j=1}^{20} f(j)\,m(j)} .
\]
In matrix terms: 1\,\% of the probability mass sits off the diagonal and
99\,\% on it. Since Dayhoff's proteins were 85\,\% identical, $\lambda < 1$ ---
downweighting the observed mutations corresponds to moving towards sequences
that are 99\,\% instead of 85\,\% identical.

\section{From \texorpdfstring{$M$}{M} to PAM Scores}

A PAM score compares the probability of an accepted point mutation in nature
with the probability of seeing the same pair by chance:
\begin{itemize}
  \item probability of seeing $a$ replaced by $b$ in nature: $f(a)\,M(a,b)$
    (pick an $a$, then it mutates into $b$);
  \item probability of seeing the pair by chance: $f(a)\,f(b)$.
\end{itemize}
Their likelihood ratio is $\dfrac{f(a)\,M(a,b)}{f(a)\,f(b)} = \dfrac{M(a,b)}{f(b)}$.

\begin{thmbox}[title={The 1 PAM scoring matrix}]
\[
  S(a,b) = \log \frac{M(a,b)}{f(b)},
\]
then multiplied by a factor (e.g.\ 10) and rounded. Because we aimed for 1\,\%
total mutation, this matrix is called \textbf{1PAM}. It is symmetric: since
$f(a)\,M(a,b) = \lambda A(a,b)/N$ and $A$ is symmetric,
$M(a,b)/f(b) = M(b,a)/f(a)$.
\end{thmbox}

\begin{figure}[htbp]
\centering
\begin{tikzpicture}[node distance=12mm]
  \node[box] (n) {$n$\\{\scriptsize counts per AA}};
  \node[box, right=of n] (f) {$f$\\{\scriptsize frequencies}};
  \node[box, below=8mm of n] (A) {$A$\\{\scriptsize raw mutation counts,}\\[-2pt]{\scriptsize symmetrised}};
  \node[box, right=of A] (M) {$M$\\{\scriptsize mutation probabilities,}\\[-2pt]{\scriptsize 1\,\% total change}};
  \node[box, right=of M, fill=figred!10] (P) {$\mathrm{PAM}$\\{\scriptsize log-likelihood ratios}};
  \draw[flow] (n) -- (f); \draw[flow] (n) -- (M); \draw[flow] (A) -- (M);
  \draw[flow] (M) -- (P); \draw[flow] (f) -| (P);
\end{tikzpicture}
\caption{The matrices involved in PAM (the mutability $m$ and the constant
$\lambda$ are hidden inside the arrow from $A$ to $M$).}
\label{fig:pam-recap}
\end{figure}

\begin{exbox}[title={1PAM for a three-letter alphabet}]
Suppose our data contain $n = (50, 30, 20)$ occurrences of A, B and C
($N = 100$, so $f = (0.5, 0.3, 0.2)$), and the symmetrised count matrix has
$A(A,B) = 4$, $A(A,C) = 2$, $A(B,C) = 1$.
\begin{enumerate}
  \item Mutabilities: $m(A) = (4+2)/50 = 0.12$, $m(B) = (4+1)/30 \approx 0.167$,
    $m(C) = (2+1)/20 = 0.15$.
  \item $\sum_j f(j)\,m(j) = 0.06 + 0.05 + 0.03 = 0.14$, so
    $\lambda = 0.01/0.14 \approx 0.0714$.
  \item Mutation matrix (rows: from, columns: to):
    \[
      M = \begin{pmatrix}
        0.99143 & 0.00571 & 0.00286 \\
        0.00952 & 0.98810 & 0.00238 \\
        0.00714 & 0.00357 & 0.98929
      \end{pmatrix},
    \]
    e.g.\ $M(A,B) = 0.0714 \cdot 4/50 = 0.00571$ and
    $M(A,A) = 1 - 0.0714 \cdot 0.12 = 0.99143$. Check:
    $\sum_a f(a)\,(1 - M(a,a)) = 0.01$. \checkmark
  \item Scores $10 \log_{10}\bigl(M(a,b)/f(b)\bigr)$, rounded:
    \[
      \mathrm{1PAM} = \begin{pmatrix} 3 & -17 & -18 \\ -17 & 5 & -19 \\ -18 & -19 & 7 \end{pmatrix}.
    \]
    For example,
    \[
      S(A,A) = 10\log_{10}\tfrac{0.99143}{0.5} = 2.97 \approx 3, \qquad
      S(A,B) = 10\log_{10}\tfrac{0.00571}{0.3} = -17.2,
    \]
    and $S(B,A)$ gives the same value.
\end{enumerate}
The rare amino acid C gets the highest diagonal score: seeing C aligned with
C is the most ``surprising'' match. After only 1\,\% of change, all mutations
are very unlikely --- hence the strongly negative off-diagonal entries.
\end{exbox}

\section{Sidestep: Markov Matrices}
\label{sec:markov}

To get matrices for longer time spans, we need one property of PAM: it is a
\textbf{Markov matrix}.

\begin{defbox}[title={Markov (transition) matrix}]
Given a graph with transition probabilities along its edges, written as an
$n \times n$ matrix $T$ with $T_{ij}$ = probability of going from node $i$ to
node $j$ in one step: each row sums to 1 (``you have to go somewhere, even if
you stay''). The probabilities for $k$ steps are given by the $k$-th matrix
power, $T^k$.
\end{defbox}

\begin{figure}[htbp]
\centering
\begin{tikzpicture}[node distance=24mm,
    st/.style={circle, draw=accent, thick, fill=accent!8, minimum size=8mm, font=\small\bfseries}]
  \node[st] (S) {S};
  \node[st, right=of S] (R) {R};
  \node[st, right=of R] (F) {F};
  \draw[flow] (S) -- node[above, lbl] {0.77} (R);
  \draw[flow] (R) -- node[above, lbl] {0.35} (F);
  \draw[flow] (F) to[bend left=35] node[below, lbl] {0.95} (S);
  \draw[flow] (S) to[loop above, looseness=6] node[above, lbl] {0.23} (S);
  \draw[flow] (R) to[loop above, looseness=6] node[above, lbl] {0.65} (R);
  \draw[flow] (F) to[loop above, looseness=6] node[above, lbl] {0.05} (F);
  \node[anchor=west, font=\small] at (8.0,0.2) {%
    $T = \begin{pmatrix} 0.23 & 0.77 & 0 \\ 0 & 0.65 & 0.35 \\ 0.95 & 0 & 0.05 \end{pmatrix}$};
\end{tikzpicture}
\caption{A three-state Markov chain (the slide's example) and its transition
matrix; rows and columns in the order S, R, F.}
\label{fig:markov}
\end{figure}

\begin{exbox}[title={Two steps}]
\[
  T^2 = \begin{pmatrix} 0.053 & 0.678 & 0.270 \\ 0.333 & 0.423 & 0.245 \\ 0.266 & 0.732 & 0.003 \end{pmatrix}.
\]
For example, $T^2_{SR} = 0.23 \cdot 0.77 + 0.77 \cdot 0.65 + 0 \cdot 0 = 0.678$:
either stay in S and then go to R, or go to R and stay there. After many
steps, every row converges to the same vector ($T^{50}$ has all rows equal to
$(0.249, 0.549, 0.202)$): the chain ``forgets'' where it started and reaches
its \emph{stationary distribution}.
\end{exbox}

\section{PAMs for Longer Time Spans}

\begin{thmbox}[title={$n$PAM}]
Since 1PAM describes one step of evolution as a Markov matrix, $n$ steps are
described by the $n$-th power,
\[
  M^{(n)} = M^{n} \qquad (\text{``}n\text{PAM} = 1\text{PAM}^n\text{''}),
\]
and the $n$PAM scores are computed from $M^n$ exactly as before,
\[
  S_n(a,b) = \log\frac{M^n(a,b)}{f(b)} .
\]
Matrices up to 250PAM are in use.
Rule of thumb for small $n$: with $n$PAM you can expect about $n\,\%$
mutations.
\end{thmbox}

\begin{notebox}[title={Power of the probabilities, not of the scores}]
``$n\mathrm{PAM} = 1\mathrm{PAM}^n$'' on the slides refers to the \emph{mutation
probability matrix} $M$. Raising the integer \emph{score} matrix to the $n$-th
power would be meaningless --- scores are logarithms. Compute $M^n$ first,
then take logs.
\end{notebox}

\paragraph{Does 100PAM mean 100\,\% changed? And 250PAM}
No. Over long time spans, some residues change several times, some change back
to their original, and some do not change at all. The higher $n$, the more the
mutation probabilities approach 100\,\% --- but only \emph{asymptotically}.
(In practice, 250PAM corresponds to sequences that are still about 20\,\%
identical, \cref{tab:rule-of-thumb}.) The transition probabilities ``walk
away'' from the diagonal as the time span grows (\cref{fig:pam-decay}).

\begin{figure}[htbp]
\centering
\begin{tikzpicture}
\begin{axis}[bsa, width=0.66\linewidth, height=5cm, xmin=0, xmax=1000, ymin=0.3, ymax=1.02,
    xlabel={$n$ in $n$PAM}, ylabel={expected identity}, ytick={0.38,0.5,0.6,0.7,0.8,0.9,1},
    yticklabel style={/pgf/number format/fixed}]
  \addplot[very thick, accent, mark=*, mark size=1.2pt] coordinates {
    (0,1.0) (10,0.9071) (25,0.7935) (50,0.6567) (75,0.5657) (100,0.5051) (150,0.4373)
    (200,0.4066) (250,0.3925) (300,0.3859) (400,0.3814) (500,0.3803) (750,0.38) (1000,0.38)};
  \addplot[dashed, figred, thick, domain=0:1000] {0.38};
  \node[font=\scriptsize, figred, anchor=south east] at (axis cs:1000,0.38) {random: $\sum_a f(a)^2 = 0.38$};
\end{axis}
\end{tikzpicture}
\caption{The three-letter example: probability that a position is unchanged
after $n$PAM, $\sum_a f(a)\,M^n(a,a)$. It starts at 1, drops by about 1\,\% per
PAM at first, and then levels off at the value two \emph{random} sequences
would have. After 100PAM, about half of the positions are still identical,
not zero.}
\label{fig:pam-decay}
\end{figure}

\begin{exbox}[title={The three-letter example over time}]
\small
\centering
\begin{tabular}{@{}lccl@{}}
\toprule
 & identity & $M^n$ diagonal & scores $10\log_{10}(M^n(a,b)/f(b))$ \\
\midrule
1PAM   & 0.99  & $(0.991, 0.988, 0.989)$ & diag $(3,5,7)$, off-diag $-17 \ldots -19$ \\
10PAM  & 0.91  & $(0.921, 0.890, 0.899)$ & diag $(3,5,7)$, off-diag $-8, -9$ \\
50PAM  & 0.66  & $(0.711, 0.598, 0.608)$ & diag $(2,3,5)$, off-diag $-2, -3$ \\
100PAM & 0.51  & $(0.590, 0.428, 0.408)$ & diag $(1,2,3)$, off-diag $-1$ \\
250PAM & 0.39  & $(0.507, 0.311, 0.228)$ & diag $(0,0,1)$, off-diag $0$ \\
$\infty$PAM & 0.38 & $(0.5, 0.3, 0.2) = f$ & all $0$ \\
\bottomrule
\end{tabular}

\raggedright
\medskip
The longer the time span, the less a match is worth and the less a mismatch
hurts --- exactly the NCBI intuition from \cref{sec:transitions}.
\end{exbox}

\begin{notebox}[title={Exam-style question: what does 0PAM look like? And 10,000,000PAM}]
Regl used to give extra points for a quick answer:
\begin{itemize}
  \item \textbf{0PAM}: no time has passed, nothing has changed --- $M^0$ is the
    \textbf{identity matrix} (1 on the diagonal, 0 elsewhere). Every amino
    acid is only ever aligned with itself.
  \item \textbf{10,000,000PAM}: after a practically infinite time the chain has
    forgotten where it started: every row of $M^n$ equals the background
    frequencies $f$, so ``observed = expected'' for every pair and all scores
    become 0. The distribution is random --- there is no difference between
    diagonal and off-diagonal any more, and an alignment carries no
    information.
\end{itemize}
\end{notebox}

\section{PAM in Practice}

Always use a PAM that fits your project: to compare humans with chimpanzees,
use a low PAM (the slides suggest 20PAM or 50PAM; Regl also mentions 5PAM or
10PAM); to compare humans with fish, use 250PAM.

\paragraph{What 250PAM looks like}
In the 250PAM matrix, all diagonal entries are positive. Cysteine (C, 12),
tryptophan (W, 17) and tyrosine (Y, 10) have particularly high diagonal values:
they strongly resist mutation. Phenylalanine and tyrosine (F--Y, $+7$) replace
each other easily, both being large aromatic amino acids.

\paragraph{Minitask from the slides: find the error}
The slides show a short MATLAB program that computes 1PAM for six amino acids
and ask who finds the error. The relevant line is
\begin{center}
\verb|p(i,j) = tsym(i,j) / 100 * frel(i) * ss;|
\end{center}
where \verb|ss| is the sum of all off-diagonal counts. From the formulas above,
$\lambda = 0.01\,N/\verb|ss|$, so the intended value is
$\lambda A(i,j)/n(i) = A(i,j) / (100 \cdot f(i) \cdot \verb|ss|)$. Because of
operator precedence, the code \emph{multiplies} by \verb|frel(i) * ss|
instead of dividing --- the parentheses are missing. (The diagonal also needs
special treatment, $1 - \lambda m$, rather than the same formula.) A full
worked PAM example is in the lecture notes (``Bioinformatics I'' script,
Appendix C.1.1).

\section{Problems with PAM}

\begin{itemize}
  \item \textbf{Old data.} In 1974 only few sequences were available, so PAM is
    certainly biased.
  \item \textbf{Error amplification.} Small errors in 1PAM are blown up by
    repeatedly multiplying the matrix for $n$PAM --- you cannot really trust
    250PAM.
  \item \textbf{Missing rare mutations.} Using only closely related sequences,
    some rare substitutions are underrepresented or missing altogether.
  \item \textbf{A devil's circle.} A good PAM needs good phylogenetic trees,
    while good phylogenetic trees need good PAMs.
  \item \textbf{Position dependence.} The mutation rate is not constant along a
    protein; it depends on where an amino acid sits in the structure. This
    observation leads directly to BLOSUM.
\end{itemize}

\section{Chapter Checklist}

\begin{itemize}
  \item what PAM stands for (accepted point mutation) and the 1PAM unit (1
    mutation per 100 amino acids);
  \item Dayhoff's data (71 families, $>85\,\%$ identity, 1572 mutations) and why
    closely related sequences are needed;
  \item counting along a parsimonious tree; symmetrising $A$;
    $n$, $f$, $m$; the mutation matrix $M$ and the choice of $\lambda$;
  \item $S(a,b) = \log(M(a,b)/f(b))$; compute a small example;
  \item Markov matrices and $n\mathrm{PAM} = M^n$; why 100PAM is not 100\,\%
    change; 0PAM and $\infty$PAM;
  \item choosing a PAM by evolutionary distance; the problems of PAM.
\end{itemize}

% ==========================================================
%  CHAPTER 8 — BLOSUM MATRICES AND CHOOSING A MATRIX
%  Lecture VO04 (A. Regl), slides 64--120 (deck p. 30--65)
%  Recording 64:50--78:00 (derivation slides were skipped in the lecture)
% ==========================================================
\chapter{BLOSUM Matrices and How to Choose a Matrix}
\label{ch:blosum}

\lectureref{VO04 (Scoring, part 2), slides 64--120 of the deck. Recording
64:50--78:00. As with PAM, the lecture skipped the derivation (``these are all
the details how you would do that; we will skip them''); this chapter works
through them with a complete small example.}

BLOSUM is the second big family of substitution matrices and the one you will
use in practice: it is the default in BLAST and most other tools. It uses the
same log-odds idea as PAM (\cref{ch:logodds}); what changes is the data it
counts.

\section{Why Blocks? Similarity Is Not Uniform}

When two related proteins are compared, the similarity is usually \emph{not}
spread evenly along the sequence. In many cases it is concentrated in a few
short regions --- \emph{blocks} --- separated by stretches that are much less
similar. The slides show this for the two DNA clamps from \cref{ch:alignment}
(BLASTing 1MMI against 1AXC) and for the human TATA-box binding protein
compared with other organisms: high similarity appears as isolated red blocks.

\paragraph{Why blocks appear}
A protein can roughly be divided into its \textbf{functional site} --- the
handful of amino acids that actually do the work, for example the ones that
bind to DNA --- and the rest, which is mostly there for structural reasons, to
hold everything in place. A mutation in the structural part often has little
or no effect. A mutation in the functional site can be disastrous: the protein
loses its function, and the organism may die. Such mutations are not
``accepted'' by evolution (\cref{ch:pam}). So functional sites
\textbf{mutate much more slowly}, and they stay similar across species
(\cref{fig:functional-blocks}).

\begin{figure}[htbp]
\centering
\begin{tikzpicture}[x=1cm, y=1cm, font=\small]
  \fill[gray!25] (0,0) rectangle (4.2,0.6);
  \fill[figgreen!55] (4.2,0) rectangle (6.0,0.6);
  \fill[gray!25] (6.0,0) rectangle (10.4,0.6);
  \draw[thick] (0,0) rectangle (10.4,0.6);
  \node at (2.1,0.3) {structural part};
  \node at (5.1,0.3) {functional};
  \node at (8.2,0.3) {structural part};
  \node[anchor=east, font=\small\bfseries] at (-0.2,0.3) {protein};
  \node[anchor=east] at (-0.2,-0.5) {mutation rate};
  \node at (2.1,-0.5) {normal / high};
  \node[figgreen!60!black] at (5.1,-0.5) {\textbf{low}};
  \node at (8.2,-0.5) {normal / high};
  \node[anchor=east] at (-0.2,-1.1) {similarity};
  \node at (2.1,-1.1) {low};
  \node[figgreen!60!black] at (5.1,-1.1) {\textbf{high}};
  \node at (8.2,-1.1) {low};
\end{tikzpicture}

\bigskip
\begin{tikzpicture}[x=0.36cm, y=-0.45cm, font=\ttfamily\small]
\fill[figgreen!18] (8,-0.1) rectangle (14,5.1);
\node[anchor=east, font=\scriptsize\rmfamily] at (-0.3,0) {seq 1};
\node at (0.5,0) {F};
\node at (1.5,0) {W};
\node at (2.5,0) {D};
\node at (3.5,0) {K};
\node at (4.5,0) {E};
\node at (5.5,0) {S};
\node at (6.5,0) {R};
\node at (7.5,0) {S};
\node at (8.5,0) {G};
\node at (9.5,0) {D};
\node at (10.5,0) {S};
\node at (11.5,0) {G};
\node at (12.5,0) {G};
\node at (13.5,0) {P};
\node at (14.5,0) {Q};
\node at (15.5,0) {Y};
\node at (16.5,0) {A};
\node at (17.5,0) {R};
\node at (18.5,0) {K};
\node at (19.5,0) {I};
\node at (20.5,0) {W};
\node at (21.5,0) {E};
\node[anchor=east, font=\scriptsize\rmfamily] at (-0.3,1) {seq 2};
\node at (0.5,1) {M};
\node at (1.5,1) {A};
\node at (2.5,1) {A};
\node at (3.5,1) {A};
\node at (4.5,1) {V};
\node at (5.5,1) {A};
\node at (6.5,1) {P};
\node at (7.5,1) {H};
\node at (8.5,1) {G};
\node at (9.5,1) {D};
\node at (10.5,1) {S};
\node at (11.5,1) {G};
\node at (12.5,1) {G};
\node at (13.5,1) {P};
\node at (14.5,1) {V};
\node at (15.5,1) {I};
\node at (16.5,1) {N};
\node at (17.5,1) {I};
\node at (18.5,1) {I};
\node at (19.5,1) {R};
\node at (20.5,1) {L};
\node at (21.5,1) {A};
\node[anchor=east, font=\scriptsize\rmfamily] at (-0.3,2) {seq 3};
\node at (0.5,2) {Q};
\node at (1.5,2) {V};
\node at (2.5,2) {E};
\node at (3.5,2) {G};
\node at (4.5,2) {L};
\node at (5.5,2) {E};
\node at (6.5,2) {M};
\node at (7.5,2) {T};
\node at (8.5,2) {G};
\node at (9.5,2) {D};
\node at (10.5,2) {S};
\node at (11.5,2) {A};
\node at (12.5,2) {G};
\node at (13.5,2) {P};
\node at (14.5,2) {S};
\node at (15.5,2) {T};
\node at (16.5,2) {P};
\node at (17.5,2) {W};
\node at (18.5,2) {C};
\node at (19.5,2) {S};
\node at (20.5,2) {I};
\node at (21.5,2) {P};
\node[anchor=east, font=\scriptsize\rmfamily] at (-0.3,3) {seq 4};
\node at (0.5,3) {Q};
\node at (1.5,3) {G};
\node at (2.5,3) {N};
\node at (3.5,3) {V};
\node at (4.5,3) {N};
\node at (5.5,3) {D};
\node at (6.5,3) {R};
\node at (7.5,3) {T};
\node at (8.5,3) {G};
\node at (9.5,3) {D};
\node at (10.5,3) {S};
\node at (11.5,3) {G};
\node at (12.5,3) {G};
\node at (13.5,3) {P};
\node at (14.5,3) {A};
\node at (15.5,3) {S};
\node at (16.5,3) {C};
\node at (17.5,3) {L};
\node at (18.5,3) {Y};
\node at (19.5,3) {W};
\node at (20.5,3) {W};
\node at (21.5,3) {P};
\node[anchor=east, font=\scriptsize\rmfamily] at (-0.3,4) {seq 5};
\node at (0.5,4) {G};
\node at (1.5,4) {G};
\node at (2.5,4) {T};
\node at (3.5,4) {I};
\node at (4.5,4) {A};
\node at (5.5,4) {H};
\node at (6.5,4) {V};
\node at (7.5,4) {V};
\node at (8.5,4) {G};
\node at (9.5,4) {D};
\node at (10.5,4) {S};
\node at (11.5,4) {G};
\node at (12.5,4) {G};
\node at (13.5,4) {A};
\node at (14.5,4) {R};
\node at (15.5,4) {K};
\node at (16.5,4) {V};
\node at (17.5,4) {Y};
\node at (18.5,4) {A};
\node at (19.5,4) {P};
\node at (20.5,4) {T};
\node at (21.5,4) {F};
\node[anchor=east, font=\scriptsize\rmfamily] at (-0.3,5) {seq 6};
\node at (0.5,5) {T};
\node at (1.5,5) {V};
\node at (2.5,5) {H};
\node at (3.5,5) {Q};
\node at (4.5,5) {C};
\node at (5.5,5) {S};
\node at (6.5,5) {N};
\node at (7.5,5) {W};
\node at (8.5,5) {G};
\node at (9.5,5) {D};
\node at (10.5,5) {S};
\node at (11.5,5) {G};
\node at (12.5,5) {G};
\node at (13.5,5) {P};
\node at (14.5,5) {Q};
\node at (15.5,5) {N};
\node at (16.5,5) {A};
\node at (17.5,5) {V};
\node at (18.5,5) {V};
\node at (19.5,5) {Y};
\node at (20.5,5) {Y};
\node at (21.5,5) {M};
\node[font=\scriptsize\rmfamily, figgreen!60!black, anchor=north] at (11.0,5.3) {conserved block};
\end{tikzpicture}

\caption{Top: the functional site of a protein mutates slowly, so it stays
similar across related proteins. Bottom: in a multiple alignment of related
sequences this shows up as a \emph{conserved block} (a made-up illustration;
only two positions in the block have changed).}
\label{fig:functional-blocks}
\end{figure}

\begin{intbox}[title={The software analogy}]
Think of a widely used library. Its public API --- function names and
signatures that thousands of callers depend on --- practically never changes,
because any change breaks everything. Internal helper code and comments are
rewritten all the time. If you compared ten versions of the library, the API
would be a ``conserved block''. Functional sites are the public API of a
protein. Histones, the proteins DNA is wrapped around (\cref{sec:packing}), are
an extreme case: almost their entire sequence is ``API'', and they are among
the most highly conserved proteins known.
\end{intbox}

The slides' example is \textbf{vitronectin}, an essential part of human blood
serum that is also found in the bone matrix: a mutation in one region has
little or no effect on its function, a mutation in another possibly a
disastrous one.

\section{The BLOSUM Idea}
\label{sec:blosum}

\begin{defbox}[title={BLOSUM = BLOcks SUbstitution Matrix}]
A substitution matrix computed from \textbf{blocks}: short, highly conserved
regions of multiple alignments of related proteins, \emph{without gaps}. The
blocks come from the BLOCKS database.\footnote{BLOSUM was introduced by
Steven and Jorja Henikoff in 1992.} \textbf{BLOSUM$x$} is the matrix built
from blocks in which sequences that are more than $x\,\%$ identical have been
merged, so that no two remaining (groups of) sequences are more than
$x\,\%$ similar.
\end{defbox}

\begin{center}
\small
\begin{tabular}{@{}p{0.45\linewidth}p{0.45\linewidth}@{}}
\toprule
\textbf{How PAMs are made} & \textbf{How BLOSUMs are made} \\
\midrule
Take closely related sequences (whole proteins). &
Take highly conserved subsequences (blocks without gaps) from groups of
proteins. \\
Build family trees, count mutations along the branches. &
Eliminate or merge sequences that are too similar (more than $x\,\%$). \\
Scale to 1\,\% change (1PAM), extrapolate by matrix powers ($n$PAM). &
Count substitutions directly in the columns of the blocks --- one
calculation per matrix, no extrapolation. \\
Log-odds, scale, round. & Log-odds, scale (factor 2), round. \\
\bottomrule
\end{tabular}
\end{center}

\subsection{Blocks}

A block is a set of aligned subsequences in which the columns are highly
similar (or identical). Blocks
\begin{enumerate}
  \item all have the \textbf{same length} (they are truncated if necessary), and
  \item contain \textbf{no insertions or deletions} --- only substitutions.
\end{enumerate}
If one sequence has a gap inside a region, it is simply left out of that block
(on the slide, the beetle genera \emph{Pseudomorpha} and \emph{Carenum} are
dropped from one block each because of their gaps). How the underlying multiple
alignment is made in the first place is the topic of \cref{ch:msa}.

\subsection{``Closely Related, but Not Too Close''}

The guiding principle when selecting sequences is the same as for PAM. But
BLOSUM must also solve a different problem: databases are full of near
duplicates.

\begin{exbox}[title={12,000 human insulins}]
Searching UniProt for ``insulin'' gives about 292,000 results; restricted to
human insulin there are still nearly 12,000 entries --- insulin bound to many
different molecules, in many configurations, variants and fragments. If we
counted all of them, insulin would be grossly overrepresented, and the
matrix would mostly describe insulin rather than proteins in general.
\end{exbox}

So first, compute all pairwise similarities (the relative number of matches
after alignment, in \%) and mark every pair that is at least $x\,\%$ similar.
Groups of such close relatives form \emph{clusters}. Then there are three
ways to handle a cluster:

\begin{itemize}
  \item \textbf{Elimination}: keep one representative (e.g.\ the first) and
    remove all others;
  \item \textbf{Similarization}: replace the cluster by a single \emph{consensus}
    sequence (majority vote in each column);
  \item \textbf{Weighting}: keep all $n$ sequences of the cluster, but give each
    a weight of $1/n$.
\end{itemize}

\begin{exbox}[title={Elimination, similarization and weighting at 75\,\%}]
Seven sequences of length 4; ``75\,\% identity'' means 3 of 4 positions agree.
\begin{center}
\small
\begin{tabular}{@{}lll@{\qquad}l@{\qquad}l@{\qquad}l@{}}
\toprule
\multicolumn{3}{@{}l}{\textbf{Input}} & \textbf{Elimination} & \textbf{Similarization} & \textbf{Weighting} \\
\midrule
1 & \seq{VIIL} & \rdelim\}{2}{*}[\,cluster] & \seq{VIIL} & \seq{VIIL} & \seq{VIIL} (50\,\%) \\
2 & \seq{VIII} & & --- & & \seq{VIII} (50\,\%) \\
3 & \seq{LIVV} & & \seq{LIVV} & \seq{LIVV} & \seq{LIVV} (100\,\%) \\
4 & \seq{LLVI} & & \seq{LLVI} & \seq{LLVI} & \seq{LLVI} (100\,\%) \\
5 & \seq{VLLI} & \rdelim\}{3}{*}[\,cluster] & \seq{VLLI} & \seq{VLLV} & \seq{VLLI} (33\,\%) \\
6 & \seq{VLLV} & & --- & & \seq{VLLV} (33\,\%) \\
7 & \seq{VLLV} & & --- & & \seq{VLLV} (33\,\%) \\
\bottomrule
\end{tabular}
\end{center}
Sequences 1 and 2 differ only in the last position, and so do 5, 6 and 7.
Elimination keeps the first one of each cluster. Similarization takes the
majority in each column: \seq{VLLV} wins 2:1 in the second cluster; in the
first, the last position is a tie, and the slide keeps \seq{L}. Weighting can
also be read as one sequence per cluster with a ``fractional'' last letter:
\seq{VII}[50\,\% L, 50\,\% I] and \seq{VLL}[33\,\% I, 67\,\% V].
\end{exbox}

\begin{notebox}[title={The BLOSUM number runs the ``wrong'' way}]
In BLOSUM$x$, sequences up to $x\,\%$ similarity are kept as separate
sequences. A \textbf{high} $x$ (BLOSUM90) keeps very similar sequences: the
matrix describes \textbf{short} evolutionary distances. A \textbf{low} $x$
(BLOSUM45) merges everything above 45\,\%, so only distant relationships are
counted: the matrix is for \textbf{distant} relatives. For PAM it is the
opposite: a high number means a long distance (\cref{fig:pam-blosum-axis}).
\end{notebox}

\section{Building a BLOSUM, Step by Step}
\label{sec:blosum-steps}

The procedure on the slides, for sequences from the BLOCKS database with
about $x\,\%$ identity:
\begin{enumerate}
  \item cluster the sequences and merge those with more than $x\,\%$ identity;
  \item count the amino-acid pairs in each column and sum them up: matrix $C$;
  \item compute the normalising constant $T$ and the observed pair
    probabilities $Q$;
  \item compute the probabilities $p$ of the single amino acids and from them
    the expected pair probabilities $E$;
  \item compute the log-odds scores $S = \log_2(Q/E)$, scale them and round
    them to the nearest integer.
\end{enumerate}
Everything we are aiming for is the familiar formula
$S(a,b) = \log \bigl(q_{a,b}/e_{a,b}\bigr)$: observed mutation probability over
the probability by chance.

\subsection{Step 1: Counting Pairs, Giving \texorpdfstring{$C$}{C}}

Take each block, look at each column separately, and count every pair of
amino acids in that column --- \emph{each sequence with each sequence below
it}, including identical pairs, exactly as in \cref{sec:logodds-example}.
Add up the counts over all columns (and all blocks) to get the symmetric count
matrix $C = (C_{ij})$.

\begin{exbox}[title={Possible pairings in one column}]
\begin{minipage}[c]{0.14\linewidth}
\centering\ttfamily
A\\A\\B\\A\\C\\A
\end{minipage}%
\begin{minipage}[c]{0.84\linewidth}
A column with four A, one B and one C. All possible pairs:
\begin{itemize}
  \item A--A: choose 2 of the 4 A's: $\binom{4}{2} = 6$;
  \item A--B: each of the 4 A's with the B: $4 \times 1 = 4$;
  \item A--C: $4 \times 1 = 4$;
  \item B--C: $1 \times 1 = 1$.
\end{itemize}
Check: $6 + 4 + 4 + 1 = 15 = \binom{6}{2} = \frac{6 \cdot 5}{2}$, the
number of pairs of 6 sequences. \checkmark
\end{minipage}
\end{exbox}

The example already contains the shortcut. If $n_i^{(k)}$ is the number of
times amino acid $i$ occurs in column $k$:

\begin{thmbox}[title={Shortcut for $C$}]
\[
  C_{ij} = \sum_{k} c_{ij}^{(k)},
  \qquad
  c_{ij}^{(k)} =
  \begin{cases}
    \dbinom{n_i^{(k)}}{2} = \dfrac{n_i^{(k)}\bigl(n_i^{(k)} - 1\bigr)}{2} & \text{if } i = j,\\[2ex]
    n_i^{(k)}\, n_j^{(k)} & \text{otherwise.}
  \end{cases}
\]
Same letter: choose two of them. Different letters: every $i$ with every $j$.
\end{thmbox}

\subsection{Step 2: Normalising, Giving \texorpdfstring{$Q$}{Q}}

The total number of pairs is the normalising constant. With $L$ columns in the
block and $N$ sequences, every column contributes $\binom{N}{2}$ pairs, so
\[
  T = \sum_{i \ge j} C_{ij} = L \times \frac{N(N-1)}{2},
  \qquad
  Q_{ij} = \frac{C_{ij}}{T}.
\]
$Q_{ij}$ is the \textbf{observed} probability of the pair $i$--$j$. (Sum over
$i \ge j$ only: the pair A--C is the same as C--A and must not be counted
twice.)

\subsection{Step 3: Single Amino-Acid Probabilities \texorpdfstring{$p$}{p}}

\[
  p_i = Q_{ii} + \sum_{j \ne i} \frac{Q_{ij}}{2}.
\]
Intuition: every pair consists of two letters. In a pair $i$--$i$ both letters
are $i$; in a pair $i$--$j$ only half of the letters are $i$. So $p_i$ is the
fraction of all letters in all pairs that are $i$.

\subsection{Step 4: Expected Probabilities \texorpdfstring{$E$}{E}}

If the two amino acids in a pair were independent (two dice), we would expect
\[
  E_{ab} =
  \begin{cases}
    p_a^2 & \text{if } a = b,\\
    p_a p_b + p_b p_a = 2\, p_a p_b & \text{otherwise.}
  \end{cases}
\]
The factor 2 is there because an unordered pair $a$--$b$ can arise in two ways
($a$ on top and $b$ below, or the other way round), just as ``a 1 and a 2'' with
two dice has probability $2/36$, but ``two 6'' only $1/36$.

\subsection{Step 5: Scores and the Final Matrix}
\label{sec:blosum-final}

\begin{thmbox}[title={From $Q$ and $E$ to BLOSUM}]
\[
  S(a,b) = \log_2 \frac{Q_{ab}}{E_{ab}},
  \qquad
  \mathrm{BLOSUM}(a,b) = \mathrm{round}\bigl(2 \times S(a,b)\bigr).
\]
The factor 2 is the usual one for BLOSUM (scores in half-bit units), but any
factor would work. As always:
$\mathrm{BLOSUM}(a,b) > 0$ means the substitution appears more often than by
chance, $\approx 0$ equally often, $< 0$ less often.
\end{thmbox}

\begin{figure}[htbp]
\centering
\begin{tikzpicture}[node distance=7mm]
  \node[box] (C) {$C$\\{\scriptsize pair counts over}\\[-2pt]{\scriptsize all columns, all blocks}};
  \node[box, right=of C] (Q) {$Q$\\{\scriptsize observed pair}\\[-2pt]{\scriptsize probabilities}};
  \node[box, right=of Q] (E) {$E$\\{\scriptsize pair probabilities}\\[-2pt]{\scriptsize by chance}};
  \node[box, right=of E] (S) {$S$\\{\scriptsize log-likelihood}\\[-2pt]{\scriptsize ratios $\log_2(Q/E)$}};
  \node[box, right=of S, fill=figred!10] (B) {BLOSUM\\{\scriptsize scaled ($\times 2$)}\\[-2pt]{\scriptsize and rounded}};
  \node[box, below=6mm of Q] (p) {$p$\\{\scriptsize single amino-acid}\\[-2pt]{\scriptsize probabilities}};
  \draw[flow] (C) -- node[above, lbl] {$/T$} (Q);
  \draw[flow] (Q) -- (p);
  \draw[flow] (p) -| (E);
  \draw[flow] (Q.north) to[bend left=25] (S.north);
  \draw[flow] (E) -- (S);
  \draw[flow] (S) -- (B);
\end{tikzpicture}
\caption{The matrices involved in BLOSUM. Compare with PAM
(\cref{fig:pam-recap}): only the way of collecting the data differs.}
\label{fig:blosum-recap}
\end{figure}

\begin{exbox}[title={A complete BLOSUM from a tiny block}]
Our ``block'' has $L = 2$ columns and $N = 3$ sequences over the letters N, D
and F:
\alnblock{%
NF\\
NF\\
DF}

\textbf{Step 1, counts.} Column 1 has $n_N = 2$, $n_D = 1$:
$c_{NN} = \binom{2}{2} = 1$, $c_{ND} = 2 \times 1 = 2$, $c_{DD} = \binom{1}{2} = 0$.
Column 2 has $n_F = 3$: $c_{FF} = \binom{3}{2} = 3$. Summed:
$C_{NN} = 1$, $C_{ND} = 2$, $C_{FF} = 3$, all others 0. (Check by raw
counting: N--N once, N--D twice in column 1; F--F three times in column 2.)

\textbf{Step 2, normalise.} $T = 2 \times \frac{3 \cdot 2}{2} = 6$, so
$Q_{NN} = \tfrac16$, $Q_{ND} = \tfrac13$, $Q_{FF} = \tfrac12$.

\textbf{Step 3, single probabilities.}
$p_N = \tfrac16 + \tfrac12 \cdot \tfrac13 = \tfrac13$,\;
$p_D = 0 + \tfrac12 \cdot \tfrac13 = \tfrac16$,\;
$p_F = \tfrac12$ (sum: 1 \checkmark; indeed 2 of the 6 letters are N, 1 is D,
3 are F).

\textbf{Steps 4 and 5.}
\begin{center}
\small
\begin{tabular}{@{}lccccc@{}}
\toprule
pair & $Q$ & $E$ & $Q/E$ & $S = \log_2(Q/E)$ & BLOSUM $= \mathrm{round}(2S)$ \\
\midrule
N--N & $1/6$ & $p_N^2 = 1/9$            & 1.5 & 0.585 & 1 \\
N--D & $1/3$ & $2p_Np_D = 1/9$          & 3   & 1.585 & 3 \\
F--F & $1/2$ & $p_F^2 = 1/4$            & 2   & 1.000 & 2 \\
D--D & 0     & $1/36$                   & 0   & $-\infty$ & $-\infty$ \\
N--F & 0     & $2p_Np_F = 1/3$          & 0   & $-\infty$ & $-\infty$ \\
D--F & 0     & $2p_Dp_F = 1/6$          & 0   & $-\infty$ & $-\infty$ \\
\bottomrule
\end{tabular}
\end{center}
N--D gets the highest score: in our data, N and D are aligned three times as
often as chance predicts. The $-\infty$ entries show again that a real matrix
needs vastly more data --- with many blocks, the counts $C$ are simply summed
over all of them (the course script, formulas 3.22--3.24, covers several
blocks; a larger worked example is in its Appendix C.1.2).
\end{exbox}

\section{What BLOSUM62 Looks Like}
\label{sec:blosum62}

\begin{figure}[htbp]
\centering
\begin{tikzpicture}[x=0.52cm, y=-0.52cm, font=\scriptsize]
\fill[figblue!84] (0,0) rectangle ++(1,1);
\node[white] at (0.5,0.5) {9};
\fill[figred!24] (1,0) rectangle ++(1,1);
\node[black] at (1.5,0.5) {$-1$};
\fill[figred!24] (2,0) rectangle ++(1,1);
\node[black] at (2.5,0.5) {$-1$};
\fill[figred!52] (3,0) rectangle ++(1,1);
\node[black] at (3.5,0.5) {$-3$};
\fill[white] (4,0) rectangle ++(1,1);
\node[black] at (4.5,0.5) {0};
\fill[figred!52] (5,0) rectangle ++(1,1);
\node[black] at (5.5,0.5) {$-3$};
\fill[figred!52] (6,0) rectangle ++(1,1);
\node[black] at (6.5,0.5) {$-3$};
\fill[figred!52] (7,0) rectangle ++(1,1);
\node[black] at (7.5,0.5) {$-3$};
\fill[figred!66] (8,0) rectangle ++(1,1);
\node[black] at (8.5,0.5) {$-4$};
\fill[figred!52] (9,0) rectangle ++(1,1);
\node[black] at (9.5,0.5) {$-3$};
\fill[figred!52] (10,0) rectangle ++(1,1);
\node[black] at (10.5,0.5) {$-3$};
\fill[figred!52] (11,0) rectangle ++(1,1);
\node[black] at (11.5,0.5) {$-3$};
\fill[figred!52] (12,0) rectangle ++(1,1);
\node[black] at (12.5,0.5) {$-3$};
\fill[figred!24] (13,0) rectangle ++(1,1);
\node[black] at (13.5,0.5) {$-1$};
\fill[figred!24] (14,0) rectangle ++(1,1);
\node[black] at (14.5,0.5) {$-1$};
\fill[figred!24] (15,0) rectangle ++(1,1);
\node[black] at (15.5,0.5) {$-1$};
\fill[figred!24] (16,0) rectangle ++(1,1);
\node[black] at (16.5,0.5) {$-1$};
\fill[figred!38] (17,0) rectangle ++(1,1);
\node[black] at (17.5,0.5) {$-2$};
\fill[figred!38] (18,0) rectangle ++(1,1);
\node[black] at (18.5,0.5) {$-2$};
\fill[figred!38] (19,0) rectangle ++(1,1);
\node[black] at (19.5,0.5) {$-2$};
\node[font=\scriptsize\bfseries] at (-0.5,0.5) {C};
\node[font=\scriptsize\bfseries] at (0.5,-0.5) {C};
\fill[figred!24] (0,1) rectangle ++(1,1);
\node[black] at (0.5,1.5) {$-1$};
\fill[figblue!44] (1,1) rectangle ++(1,1);
\node[black] at (1.5,1.5) {4};
\fill[figblue!20] (2,1) rectangle ++(1,1);
\node[black] at (2.5,1.5) {1};
\fill[figred!24] (3,1) rectangle ++(1,1);
\node[black] at (3.5,1.5) {$-1$};
\fill[figblue!20] (4,1) rectangle ++(1,1);
\node[black] at (4.5,1.5) {1};
\fill[white] (5,1) rectangle ++(1,1);
\node[black] at (5.5,1.5) {0};
\fill[figblue!20] (6,1) rectangle ++(1,1);
\node[black] at (6.5,1.5) {1};
\fill[white] (7,1) rectangle ++(1,1);
\node[black] at (7.5,1.5) {0};
\fill[white] (8,1) rectangle ++(1,1);
\node[black] at (8.5,1.5) {0};
\fill[white] (9,1) rectangle ++(1,1);
\node[black] at (9.5,1.5) {0};
\fill[figred!24] (10,1) rectangle ++(1,1);
\node[black] at (10.5,1.5) {$-1$};
\fill[figred!24] (11,1) rectangle ++(1,1);
\node[black] at (11.5,1.5) {$-1$};
\fill[white] (12,1) rectangle ++(1,1);
\node[black] at (12.5,1.5) {0};
\fill[figred!24] (13,1) rectangle ++(1,1);
\node[black] at (13.5,1.5) {$-1$};
\fill[figred!38] (14,1) rectangle ++(1,1);
\node[black] at (14.5,1.5) {$-2$};
\fill[figred!38] (15,1) rectangle ++(1,1);
\node[black] at (15.5,1.5) {$-2$};
\fill[figred!38] (16,1) rectangle ++(1,1);
\node[black] at (16.5,1.5) {$-2$};
\fill[figred!38] (17,1) rectangle ++(1,1);
\node[black] at (17.5,1.5) {$-2$};
\fill[figred!38] (18,1) rectangle ++(1,1);
\node[black] at (18.5,1.5) {$-2$};
\fill[figred!52] (19,1) rectangle ++(1,1);
\node[black] at (19.5,1.5) {$-3$};
\node[font=\scriptsize\bfseries] at (-0.5,1.5) {S};
\node[font=\scriptsize\bfseries] at (1.5,-0.5) {S};
\fill[figred!24] (0,2) rectangle ++(1,1);
\node[black] at (0.5,2.5) {$-1$};
\fill[figblue!20] (1,2) rectangle ++(1,1);
\node[black] at (1.5,2.5) {1};
\fill[figblue!52] (2,2) rectangle ++(1,1);
\node[black] at (2.5,2.5) {5};
\fill[figred!24] (3,2) rectangle ++(1,1);
\node[black] at (3.5,2.5) {$-1$};
\fill[white] (4,2) rectangle ++(1,1);
\node[black] at (4.5,2.5) {0};
\fill[figred!38] (5,2) rectangle ++(1,1);
\node[black] at (5.5,2.5) {$-2$};
\fill[white] (6,2) rectangle ++(1,1);
\node[black] at (6.5,2.5) {0};
\fill[figred!24] (7,2) rectangle ++(1,1);
\node[black] at (7.5,2.5) {$-1$};
\fill[figred!24] (8,2) rectangle ++(1,1);
\node[black] at (8.5,2.5) {$-1$};
\fill[figred!24] (9,2) rectangle ++(1,1);
\node[black] at (9.5,2.5) {$-1$};
\fill[figred!38] (10,2) rectangle ++(1,1);
\node[black] at (10.5,2.5) {$-2$};
\fill[figred!24] (11,2) rectangle ++(1,1);
\node[black] at (11.5,2.5) {$-1$};
\fill[figred!24] (12,2) rectangle ++(1,1);
\node[black] at (12.5,2.5) {$-1$};
\fill[figred!24] (13,2) rectangle ++(1,1);
\node[black] at (13.5,2.5) {$-1$};
\fill[figred!24] (14,2) rectangle ++(1,1);
\node[black] at (14.5,2.5) {$-1$};
\fill[figred!24] (15,2) rectangle ++(1,1);
\node[black] at (15.5,2.5) {$-1$};
\fill[white] (16,2) rectangle ++(1,1);
\node[black] at (16.5,2.5) {0};
\fill[figred!38] (17,2) rectangle ++(1,1);
\node[black] at (17.5,2.5) {$-2$};
\fill[figred!38] (18,2) rectangle ++(1,1);
\node[black] at (18.5,2.5) {$-2$};
\fill[figred!38] (19,2) rectangle ++(1,1);
\node[black] at (19.5,2.5) {$-2$};
\node[font=\scriptsize\bfseries] at (-0.5,2.5) {T};
\node[font=\scriptsize\bfseries] at (2.5,-0.5) {T};
\fill[figred!52] (0,3) rectangle ++(1,1);
\node[black] at (0.5,3.5) {$-3$};
\fill[figred!24] (1,3) rectangle ++(1,1);
\node[black] at (1.5,3.5) {$-1$};
\fill[figred!24] (2,3) rectangle ++(1,1);
\node[black] at (2.5,3.5) {$-1$};
\fill[figblue!68] (3,3) rectangle ++(1,1);
\node[white] at (3.5,3.5) {7};
\fill[figred!24] (4,3) rectangle ++(1,1);
\node[black] at (4.5,3.5) {$-1$};
\fill[figred!38] (5,3) rectangle ++(1,1);
\node[black] at (5.5,3.5) {$-2$};
\fill[figred!38] (6,3) rectangle ++(1,1);
\node[black] at (6.5,3.5) {$-2$};
\fill[figred!24] (7,3) rectangle ++(1,1);
\node[black] at (7.5,3.5) {$-1$};
\fill[figred!24] (8,3) rectangle ++(1,1);
\node[black] at (8.5,3.5) {$-1$};
\fill[figred!24] (9,3) rectangle ++(1,1);
\node[black] at (9.5,3.5) {$-1$};
\fill[figred!38] (10,3) rectangle ++(1,1);
\node[black] at (10.5,3.5) {$-2$};
\fill[figred!38] (11,3) rectangle ++(1,1);
\node[black] at (11.5,3.5) {$-2$};
\fill[figred!24] (12,3) rectangle ++(1,1);
\node[black] at (12.5,3.5) {$-1$};
\fill[figred!38] (13,3) rectangle ++(1,1);
\node[black] at (13.5,3.5) {$-2$};
\fill[figred!52] (14,3) rectangle ++(1,1);
\node[black] at (14.5,3.5) {$-3$};
\fill[figred!52] (15,3) rectangle ++(1,1);
\node[black] at (15.5,3.5) {$-3$};
\fill[figred!38] (16,3) rectangle ++(1,1);
\node[black] at (16.5,3.5) {$-2$};
\fill[figred!66] (17,3) rectangle ++(1,1);
\node[black] at (17.5,3.5) {$-4$};
\fill[figred!52] (18,3) rectangle ++(1,1);
\node[black] at (18.5,3.5) {$-3$};
\fill[figred!66] (19,3) rectangle ++(1,1);
\node[black] at (19.5,3.5) {$-4$};
\node[font=\scriptsize\bfseries] at (-0.5,3.5) {P};
\node[font=\scriptsize\bfseries] at (3.5,-0.5) {P};
\fill[white] (0,4) rectangle ++(1,1);
\node[black] at (0.5,4.5) {0};
\fill[figblue!20] (1,4) rectangle ++(1,1);
\node[black] at (1.5,4.5) {1};
\fill[white] (2,4) rectangle ++(1,1);
\node[black] at (2.5,4.5) {0};
\fill[figred!24] (3,4) rectangle ++(1,1);
\node[black] at (3.5,4.5) {$-1$};
\fill[figblue!44] (4,4) rectangle ++(1,1);
\node[black] at (4.5,4.5) {4};
\fill[white] (5,4) rectangle ++(1,1);
\node[black] at (5.5,4.5) {0};
\fill[figred!38] (6,4) rectangle ++(1,1);
\node[black] at (6.5,4.5) {$-2$};
\fill[figred!38] (7,4) rectangle ++(1,1);
\node[black] at (7.5,4.5) {$-2$};
\fill[figred!24] (8,4) rectangle ++(1,1);
\node[black] at (8.5,4.5) {$-1$};
\fill[figred!24] (9,4) rectangle ++(1,1);
\node[black] at (9.5,4.5) {$-1$};
\fill[figred!38] (10,4) rectangle ++(1,1);
\node[black] at (10.5,4.5) {$-2$};
\fill[figred!24] (11,4) rectangle ++(1,1);
\node[black] at (11.5,4.5) {$-1$};
\fill[figred!24] (12,4) rectangle ++(1,1);
\node[black] at (12.5,4.5) {$-1$};
\fill[figred!24] (13,4) rectangle ++(1,1);
\node[black] at (13.5,4.5) {$-1$};
\fill[figred!24] (14,4) rectangle ++(1,1);
\node[black] at (14.5,4.5) {$-1$};
\fill[figred!24] (15,4) rectangle ++(1,1);
\node[black] at (15.5,4.5) {$-1$};
\fill[white] (16,4) rectangle ++(1,1);
\node[black] at (16.5,4.5) {0};
\fill[figred!38] (17,4) rectangle ++(1,1);
\node[black] at (17.5,4.5) {$-2$};
\fill[figred!38] (18,4) rectangle ++(1,1);
\node[black] at (18.5,4.5) {$-2$};
\fill[figred!52] (19,4) rectangle ++(1,1);
\node[black] at (19.5,4.5) {$-3$};
\node[font=\scriptsize\bfseries] at (-0.5,4.5) {A};
\node[font=\scriptsize\bfseries] at (4.5,-0.5) {A};
\fill[figred!52] (0,5) rectangle ++(1,1);
\node[black] at (0.5,5.5) {$-3$};
\fill[white] (1,5) rectangle ++(1,1);
\node[black] at (1.5,5.5) {0};
\fill[figred!38] (2,5) rectangle ++(1,1);
\node[black] at (2.5,5.5) {$-2$};
\fill[figred!38] (3,5) rectangle ++(1,1);
\node[black] at (3.5,5.5) {$-2$};
\fill[white] (4,5) rectangle ++(1,1);
\node[black] at (4.5,5.5) {0};
\fill[figblue!60] (5,5) rectangle ++(1,1);
\node[black] at (5.5,5.5) {6};
\fill[white] (6,5) rectangle ++(1,1);
\node[black] at (6.5,5.5) {0};
\fill[figred!24] (7,5) rectangle ++(1,1);
\node[black] at (7.5,5.5) {$-1$};
\fill[figred!38] (8,5) rectangle ++(1,1);
\node[black] at (8.5,5.5) {$-2$};
\fill[figred!38] (9,5) rectangle ++(1,1);
\node[black] at (9.5,5.5) {$-2$};
\fill[figred!38] (10,5) rectangle ++(1,1);
\node[black] at (10.5,5.5) {$-2$};
\fill[figred!38] (11,5) rectangle ++(1,1);
\node[black] at (11.5,5.5) {$-2$};
\fill[figred!38] (12,5) rectangle ++(1,1);
\node[black] at (12.5,5.5) {$-2$};
\fill[figred!52] (13,5) rectangle ++(1,1);
\node[black] at (13.5,5.5) {$-3$};
\fill[figred!66] (14,5) rectangle ++(1,1);
\node[black] at (14.5,5.5) {$-4$};
\fill[figred!66] (15,5) rectangle ++(1,1);
\node[black] at (15.5,5.5) {$-4$};
\fill[figred!52] (16,5) rectangle ++(1,1);
\node[black] at (16.5,5.5) {$-3$};
\fill[figred!52] (17,5) rectangle ++(1,1);
\node[black] at (17.5,5.5) {$-3$};
\fill[figred!52] (18,5) rectangle ++(1,1);
\node[black] at (18.5,5.5) {$-3$};
\fill[figred!38] (19,5) rectangle ++(1,1);
\node[black] at (19.5,5.5) {$-2$};
\node[font=\scriptsize\bfseries] at (-0.5,5.5) {G};
\node[font=\scriptsize\bfseries] at (5.5,-0.5) {G};
\fill[figred!52] (0,6) rectangle ++(1,1);
\node[black] at (0.5,6.5) {$-3$};
\fill[figblue!20] (1,6) rectangle ++(1,1);
\node[black] at (1.5,6.5) {1};
\fill[white] (2,6) rectangle ++(1,1);
\node[black] at (2.5,6.5) {0};
\fill[figred!38] (3,6) rectangle ++(1,1);
\node[black] at (3.5,6.5) {$-2$};
\fill[figred!38] (4,6) rectangle ++(1,1);
\node[black] at (4.5,6.5) {$-2$};
\fill[white] (5,6) rectangle ++(1,1);
\node[black] at (5.5,6.5) {0};
\fill[figblue!60] (6,6) rectangle ++(1,1);
\node[black] at (6.5,6.5) {6};
\fill[figblue!20] (7,6) rectangle ++(1,1);
\node[black] at (7.5,6.5) {1};
\fill[white] (8,6) rectangle ++(1,1);
\node[black] at (8.5,6.5) {0};
\fill[white] (9,6) rectangle ++(1,1);
\node[black] at (9.5,6.5) {0};
\fill[figblue!20] (10,6) rectangle ++(1,1);
\node[black] at (10.5,6.5) {1};
\fill[white] (11,6) rectangle ++(1,1);
\node[black] at (11.5,6.5) {0};
\fill[white] (12,6) rectangle ++(1,1);
\node[black] at (12.5,6.5) {0};
\fill[figred!38] (13,6) rectangle ++(1,1);
\node[black] at (13.5,6.5) {$-2$};
\fill[figred!52] (14,6) rectangle ++(1,1);
\node[black] at (14.5,6.5) {$-3$};
\fill[figred!52] (15,6) rectangle ++(1,1);
\node[black] at (15.5,6.5) {$-3$};
\fill[figred!52] (16,6) rectangle ++(1,1);
\node[black] at (16.5,6.5) {$-3$};
\fill[figred!52] (17,6) rectangle ++(1,1);
\node[black] at (17.5,6.5) {$-3$};
\fill[figred!38] (18,6) rectangle ++(1,1);
\node[black] at (18.5,6.5) {$-2$};
\fill[figred!66] (19,6) rectangle ++(1,1);
\node[black] at (19.5,6.5) {$-4$};
\node[font=\scriptsize\bfseries] at (-0.5,6.5) {N};
\node[font=\scriptsize\bfseries] at (6.5,-0.5) {N};
\fill[figred!52] (0,7) rectangle ++(1,1);
\node[black] at (0.5,7.5) {$-3$};
\fill[white] (1,7) rectangle ++(1,1);
\node[black] at (1.5,7.5) {0};
\fill[figred!24] (2,7) rectangle ++(1,1);
\node[black] at (2.5,7.5) {$-1$};
\fill[figred!24] (3,7) rectangle ++(1,1);
\node[black] at (3.5,7.5) {$-1$};
\fill[figred!38] (4,7) rectangle ++(1,1);
\node[black] at (4.5,7.5) {$-2$};
\fill[figred!24] (5,7) rectangle ++(1,1);
\node[black] at (5.5,7.5) {$-1$};
\fill[figblue!20] (6,7) rectangle ++(1,1);
\node[black] at (6.5,7.5) {1};
\fill[figblue!60] (7,7) rectangle ++(1,1);
\node[black] at (7.5,7.5) {6};
\fill[figblue!28] (8,7) rectangle ++(1,1);
\node[black] at (8.5,7.5) {2};
\fill[white] (9,7) rectangle ++(1,1);
\node[black] at (9.5,7.5) {0};
\fill[figred!24] (10,7) rectangle ++(1,1);
\node[black] at (10.5,7.5) {$-1$};
\fill[figred!38] (11,7) rectangle ++(1,1);
\node[black] at (11.5,7.5) {$-2$};
\fill[figred!24] (12,7) rectangle ++(1,1);
\node[black] at (12.5,7.5) {$-1$};
\fill[figred!52] (13,7) rectangle ++(1,1);
\node[black] at (13.5,7.5) {$-3$};
\fill[figred!52] (14,7) rectangle ++(1,1);
\node[black] at (14.5,7.5) {$-3$};
\fill[figred!66] (15,7) rectangle ++(1,1);
\node[black] at (15.5,7.5) {$-4$};
\fill[figred!52] (16,7) rectangle ++(1,1);
\node[black] at (16.5,7.5) {$-3$};
\fill[figred!52] (17,7) rectangle ++(1,1);
\node[black] at (17.5,7.5) {$-3$};
\fill[figred!52] (18,7) rectangle ++(1,1);
\node[black] at (18.5,7.5) {$-3$};
\fill[figred!66] (19,7) rectangle ++(1,1);
\node[black] at (19.5,7.5) {$-4$};
\node[font=\scriptsize\bfseries] at (-0.5,7.5) {D};
\node[font=\scriptsize\bfseries] at (7.5,-0.5) {D};
\fill[figred!66] (0,8) rectangle ++(1,1);
\node[black] at (0.5,8.5) {$-4$};
\fill[white] (1,8) rectangle ++(1,1);
\node[black] at (1.5,8.5) {0};
\fill[figred!24] (2,8) rectangle ++(1,1);
\node[black] at (2.5,8.5) {$-1$};
\fill[figred!24] (3,8) rectangle ++(1,1);
\node[black] at (3.5,8.5) {$-1$};
\fill[figred!24] (4,8) rectangle ++(1,1);
\node[black] at (4.5,8.5) {$-1$};
\fill[figred!38] (5,8) rectangle ++(1,1);
\node[black] at (5.5,8.5) {$-2$};
\fill[white] (6,8) rectangle ++(1,1);
\node[black] at (6.5,8.5) {0};
\fill[figblue!28] (7,8) rectangle ++(1,1);
\node[black] at (7.5,8.5) {2};
\fill[figblue!52] (8,8) rectangle ++(1,1);
\node[black] at (8.5,8.5) {5};
\fill[figblue!28] (9,8) rectangle ++(1,1);
\node[black] at (9.5,8.5) {2};
\fill[white] (10,8) rectangle ++(1,1);
\node[black] at (10.5,8.5) {0};
\fill[white] (11,8) rectangle ++(1,1);
\node[black] at (11.5,8.5) {0};
\fill[figblue!20] (12,8) rectangle ++(1,1);
\node[black] at (12.5,8.5) {1};
\fill[figred!38] (13,8) rectangle ++(1,1);
\node[black] at (13.5,8.5) {$-2$};
\fill[figred!52] (14,8) rectangle ++(1,1);
\node[black] at (14.5,8.5) {$-3$};
\fill[figred!52] (15,8) rectangle ++(1,1);
\node[black] at (15.5,8.5) {$-3$};
\fill[figred!38] (16,8) rectangle ++(1,1);
\node[black] at (16.5,8.5) {$-2$};
\fill[figred!52] (17,8) rectangle ++(1,1);
\node[black] at (17.5,8.5) {$-3$};
\fill[figred!38] (18,8) rectangle ++(1,1);
\node[black] at (18.5,8.5) {$-2$};
\fill[figred!52] (19,8) rectangle ++(1,1);
\node[black] at (19.5,8.5) {$-3$};
\node[font=\scriptsize\bfseries] at (-0.5,8.5) {E};
\node[font=\scriptsize\bfseries] at (8.5,-0.5) {E};
\fill[figred!52] (0,9) rectangle ++(1,1);
\node[black] at (0.5,9.5) {$-3$};
\fill[white] (1,9) rectangle ++(1,1);
\node[black] at (1.5,9.5) {0};
\fill[figred!24] (2,9) rectangle ++(1,1);
\node[black] at (2.5,9.5) {$-1$};
\fill[figred!24] (3,9) rectangle ++(1,1);
\node[black] at (3.5,9.5) {$-1$};
\fill[figred!24] (4,9) rectangle ++(1,1);
\node[black] at (4.5,9.5) {$-1$};
\fill[figred!38] (5,9) rectangle ++(1,1);
\node[black] at (5.5,9.5) {$-2$};
\fill[white] (6,9) rectangle ++(1,1);
\node[black] at (6.5,9.5) {0};
\fill[white] (7,9) rectangle ++(1,1);
\node[black] at (7.5,9.5) {0};
\fill[figblue!28] (8,9) rectangle ++(1,1);
\node[black] at (8.5,9.5) {2};
\fill[figblue!52] (9,9) rectangle ++(1,1);
\node[black] at (9.5,9.5) {5};
\fill[white] (10,9) rectangle ++(1,1);
\node[black] at (10.5,9.5) {0};
\fill[figblue!20] (11,9) rectangle ++(1,1);
\node[black] at (11.5,9.5) {1};
\fill[figblue!20] (12,9) rectangle ++(1,1);
\node[black] at (12.5,9.5) {1};
\fill[white] (13,9) rectangle ++(1,1);
\node[black] at (13.5,9.5) {0};
\fill[figred!52] (14,9) rectangle ++(1,1);
\node[black] at (14.5,9.5) {$-3$};
\fill[figred!38] (15,9) rectangle ++(1,1);
\node[black] at (15.5,9.5) {$-2$};
\fill[figred!38] (16,9) rectangle ++(1,1);
\node[black] at (16.5,9.5) {$-2$};
\fill[figred!52] (17,9) rectangle ++(1,1);
\node[black] at (17.5,9.5) {$-3$};
\fill[figred!24] (18,9) rectangle ++(1,1);
\node[black] at (18.5,9.5) {$-1$};
\fill[figred!38] (19,9) rectangle ++(1,1);
\node[black] at (19.5,9.5) {$-2$};
\node[font=\scriptsize\bfseries] at (-0.5,9.5) {Q};
\node[font=\scriptsize\bfseries] at (9.5,-0.5) {Q};
\fill[figred!52] (0,10) rectangle ++(1,1);
\node[black] at (0.5,10.5) {$-3$};
\fill[figred!24] (1,10) rectangle ++(1,1);
\node[black] at (1.5,10.5) {$-1$};
\fill[figred!38] (2,10) rectangle ++(1,1);
\node[black] at (2.5,10.5) {$-2$};
\fill[figred!38] (3,10) rectangle ++(1,1);
\node[black] at (3.5,10.5) {$-2$};
\fill[figred!38] (4,10) rectangle ++(1,1);
\node[black] at (4.5,10.5) {$-2$};
\fill[figred!38] (5,10) rectangle ++(1,1);
\node[black] at (5.5,10.5) {$-2$};
\fill[figblue!20] (6,10) rectangle ++(1,1);
\node[black] at (6.5,10.5) {1};
\fill[figred!24] (7,10) rectangle ++(1,1);
\node[black] at (7.5,10.5) {$-1$};
\fill[white] (8,10) rectangle ++(1,1);
\node[black] at (8.5,10.5) {0};
\fill[white] (9,10) rectangle ++(1,1);
\node[black] at (9.5,10.5) {0};
\fill[figblue!76] (10,10) rectangle ++(1,1);
\node[white] at (10.5,10.5) {8};
\fill[white] (11,10) rectangle ++(1,1);
\node[black] at (11.5,10.5) {0};
\fill[figred!24] (12,10) rectangle ++(1,1);
\node[black] at (12.5,10.5) {$-1$};
\fill[figred!38] (13,10) rectangle ++(1,1);
\node[black] at (13.5,10.5) {$-2$};
\fill[figred!52] (14,10) rectangle ++(1,1);
\node[black] at (14.5,10.5) {$-3$};
\fill[figred!52] (15,10) rectangle ++(1,1);
\node[black] at (15.5,10.5) {$-3$};
\fill[figred!52] (16,10) rectangle ++(1,1);
\node[black] at (16.5,10.5) {$-3$};
\fill[figred!24] (17,10) rectangle ++(1,1);
\node[black] at (17.5,10.5) {$-1$};
\fill[figblue!28] (18,10) rectangle ++(1,1);
\node[black] at (18.5,10.5) {2};
\fill[figred!38] (19,10) rectangle ++(1,1);
\node[black] at (19.5,10.5) {$-2$};
\node[font=\scriptsize\bfseries] at (-0.5,10.5) {H};
\node[font=\scriptsize\bfseries] at (10.5,-0.5) {H};
\fill[figred!52] (0,11) rectangle ++(1,1);
\node[black] at (0.5,11.5) {$-3$};
\fill[figred!24] (1,11) rectangle ++(1,1);
\node[black] at (1.5,11.5) {$-1$};
\fill[figred!24] (2,11) rectangle ++(1,1);
\node[black] at (2.5,11.5) {$-1$};
\fill[figred!38] (3,11) rectangle ++(1,1);
\node[black] at (3.5,11.5) {$-2$};
\fill[figred!24] (4,11) rectangle ++(1,1);
\node[black] at (4.5,11.5) {$-1$};
\fill[figred!38] (5,11) rectangle ++(1,1);
\node[black] at (5.5,11.5) {$-2$};
\fill[white] (6,11) rectangle ++(1,1);
\node[black] at (6.5,11.5) {0};
\fill[figred!38] (7,11) rectangle ++(1,1);
\node[black] at (7.5,11.5) {$-2$};
\fill[white] (8,11) rectangle ++(1,1);
\node[black] at (8.5,11.5) {0};
\fill[figblue!20] (9,11) rectangle ++(1,1);
\node[black] at (9.5,11.5) {1};
\fill[white] (10,11) rectangle ++(1,1);
\node[black] at (10.5,11.5) {0};
\fill[figblue!52] (11,11) rectangle ++(1,1);
\node[black] at (11.5,11.5) {5};
\fill[figblue!28] (12,11) rectangle ++(1,1);
\node[black] at (12.5,11.5) {2};
\fill[figred!24] (13,11) rectangle ++(1,1);
\node[black] at (13.5,11.5) {$-1$};
\fill[figred!52] (14,11) rectangle ++(1,1);
\node[black] at (14.5,11.5) {$-3$};
\fill[figred!38] (15,11) rectangle ++(1,1);
\node[black] at (15.5,11.5) {$-2$};
\fill[figred!52] (16,11) rectangle ++(1,1);
\node[black] at (16.5,11.5) {$-3$};
\fill[figred!52] (17,11) rectangle ++(1,1);
\node[black] at (17.5,11.5) {$-3$};
\fill[figred!38] (18,11) rectangle ++(1,1);
\node[black] at (18.5,11.5) {$-2$};
\fill[figred!52] (19,11) rectangle ++(1,1);
\node[black] at (19.5,11.5) {$-3$};
\node[font=\scriptsize\bfseries] at (-0.5,11.5) {R};
\node[font=\scriptsize\bfseries] at (11.5,-0.5) {R};
\fill[figred!52] (0,12) rectangle ++(1,1);
\node[black] at (0.5,12.5) {$-3$};
\fill[white] (1,12) rectangle ++(1,1);
\node[black] at (1.5,12.5) {0};
\fill[figred!24] (2,12) rectangle ++(1,1);
\node[black] at (2.5,12.5) {$-1$};
\fill[figred!24] (3,12) rectangle ++(1,1);
\node[black] at (3.5,12.5) {$-1$};
\fill[figred!24] (4,12) rectangle ++(1,1);
\node[black] at (4.5,12.5) {$-1$};
\fill[figred!38] (5,12) rectangle ++(1,1);
\node[black] at (5.5,12.5) {$-2$};
\fill[white] (6,12) rectangle ++(1,1);
\node[black] at (6.5,12.5) {0};
\fill[figred!24] (7,12) rectangle ++(1,1);
\node[black] at (7.5,12.5) {$-1$};
\fill[figblue!20] (8,12) rectangle ++(1,1);
\node[black] at (8.5,12.5) {1};
\fill[figblue!20] (9,12) rectangle ++(1,1);
\node[black] at (9.5,12.5) {1};
\fill[figred!24] (10,12) rectangle ++(1,1);
\node[black] at (10.5,12.5) {$-1$};
\fill[figblue!28] (11,12) rectangle ++(1,1);
\node[black] at (11.5,12.5) {2};
\fill[figblue!52] (12,12) rectangle ++(1,1);
\node[black] at (12.5,12.5) {5};
\fill[figred!24] (13,12) rectangle ++(1,1);
\node[black] at (13.5,12.5) {$-1$};
\fill[figred!52] (14,12) rectangle ++(1,1);
\node[black] at (14.5,12.5) {$-3$};
\fill[figred!38] (15,12) rectangle ++(1,1);
\node[black] at (15.5,12.5) {$-2$};
\fill[figred!38] (16,12) rectangle ++(1,1);
\node[black] at (16.5,12.5) {$-2$};
\fill[figred!52] (17,12) rectangle ++(1,1);
\node[black] at (17.5,12.5) {$-3$};
\fill[figred!38] (18,12) rectangle ++(1,1);
\node[black] at (18.5,12.5) {$-2$};
\fill[figred!52] (19,12) rectangle ++(1,1);
\node[black] at (19.5,12.5) {$-3$};
\node[font=\scriptsize\bfseries] at (-0.5,12.5) {K};
\node[font=\scriptsize\bfseries] at (12.5,-0.5) {K};
\fill[figred!24] (0,13) rectangle ++(1,1);
\node[black] at (0.5,13.5) {$-1$};
\fill[figred!24] (1,13) rectangle ++(1,1);
\node[black] at (1.5,13.5) {$-1$};
\fill[figred!24] (2,13) rectangle ++(1,1);
\node[black] at (2.5,13.5) {$-1$};
\fill[figred!38] (3,13) rectangle ++(1,1);
\node[black] at (3.5,13.5) {$-2$};
\fill[figred!24] (4,13) rectangle ++(1,1);
\node[black] at (4.5,13.5) {$-1$};
\fill[figred!52] (5,13) rectangle ++(1,1);
\node[black] at (5.5,13.5) {$-3$};
\fill[figred!38] (6,13) rectangle ++(1,1);
\node[black] at (6.5,13.5) {$-2$};
\fill[figred!52] (7,13) rectangle ++(1,1);
\node[black] at (7.5,13.5) {$-3$};
\fill[figred!38] (8,13) rectangle ++(1,1);
\node[black] at (8.5,13.5) {$-2$};
\fill[white] (9,13) rectangle ++(1,1);
\node[black] at (9.5,13.5) {0};
\fill[figred!38] (10,13) rectangle ++(1,1);
\node[black] at (10.5,13.5) {$-2$};
\fill[figred!24] (11,13) rectangle ++(1,1);
\node[black] at (11.5,13.5) {$-1$};
\fill[figred!24] (12,13) rectangle ++(1,1);
\node[black] at (12.5,13.5) {$-1$};
\fill[figblue!52] (13,13) rectangle ++(1,1);
\node[black] at (13.5,13.5) {5};
\fill[figblue!20] (14,13) rectangle ++(1,1);
\node[black] at (14.5,13.5) {1};
\fill[figblue!28] (15,13) rectangle ++(1,1);
\node[black] at (15.5,13.5) {2};
\fill[figblue!20] (16,13) rectangle ++(1,1);
\node[black] at (16.5,13.5) {1};
\fill[white] (17,13) rectangle ++(1,1);
\node[black] at (17.5,13.5) {0};
\fill[figred!24] (18,13) rectangle ++(1,1);
\node[black] at (18.5,13.5) {$-1$};
\fill[figred!24] (19,13) rectangle ++(1,1);
\node[black] at (19.5,13.5) {$-1$};
\node[font=\scriptsize\bfseries] at (-0.5,13.5) {M};
\node[font=\scriptsize\bfseries] at (13.5,-0.5) {M};
\fill[figred!24] (0,14) rectangle ++(1,1);
\node[black] at (0.5,14.5) {$-1$};
\fill[figred!38] (1,14) rectangle ++(1,1);
\node[black] at (1.5,14.5) {$-2$};
\fill[figred!24] (2,14) rectangle ++(1,1);
\node[black] at (2.5,14.5) {$-1$};
\fill[figred!52] (3,14) rectangle ++(1,1);
\node[black] at (3.5,14.5) {$-3$};
\fill[figred!24] (4,14) rectangle ++(1,1);
\node[black] at (4.5,14.5) {$-1$};
\fill[figred!66] (5,14) rectangle ++(1,1);
\node[black] at (5.5,14.5) {$-4$};
\fill[figred!52] (6,14) rectangle ++(1,1);
\node[black] at (6.5,14.5) {$-3$};
\fill[figred!52] (7,14) rectangle ++(1,1);
\node[black] at (7.5,14.5) {$-3$};
\fill[figred!52] (8,14) rectangle ++(1,1);
\node[black] at (8.5,14.5) {$-3$};
\fill[figred!52] (9,14) rectangle ++(1,1);
\node[black] at (9.5,14.5) {$-3$};
\fill[figred!52] (10,14) rectangle ++(1,1);
\node[black] at (10.5,14.5) {$-3$};
\fill[figred!52] (11,14) rectangle ++(1,1);
\node[black] at (11.5,14.5) {$-3$};
\fill[figred!52] (12,14) rectangle ++(1,1);
\node[black] at (12.5,14.5) {$-3$};
\fill[figblue!20] (13,14) rectangle ++(1,1);
\node[black] at (13.5,14.5) {1};
\fill[figblue!44] (14,14) rectangle ++(1,1);
\node[black] at (14.5,14.5) {4};
\fill[figblue!28] (15,14) rectangle ++(1,1);
\node[black] at (15.5,14.5) {2};
\fill[figblue!36] (16,14) rectangle ++(1,1);
\node[black] at (16.5,14.5) {3};
\fill[white] (17,14) rectangle ++(1,1);
\node[black] at (17.5,14.5) {0};
\fill[figred!24] (18,14) rectangle ++(1,1);
\node[black] at (18.5,14.5) {$-1$};
\fill[figred!52] (19,14) rectangle ++(1,1);
\node[black] at (19.5,14.5) {$-3$};
\node[font=\scriptsize\bfseries] at (-0.5,14.5) {I};
\node[font=\scriptsize\bfseries] at (14.5,-0.5) {I};
\fill[figred!24] (0,15) rectangle ++(1,1);
\node[black] at (0.5,15.5) {$-1$};
\fill[figred!38] (1,15) rectangle ++(1,1);
\node[black] at (1.5,15.5) {$-2$};
\fill[figred!24] (2,15) rectangle ++(1,1);
\node[black] at (2.5,15.5) {$-1$};
\fill[figred!52] (3,15) rectangle ++(1,1);
\node[black] at (3.5,15.5) {$-3$};
\fill[figred!24] (4,15) rectangle ++(1,1);
\node[black] at (4.5,15.5) {$-1$};
\fill[figred!66] (5,15) rectangle ++(1,1);
\node[black] at (5.5,15.5) {$-4$};
\fill[figred!52] (6,15) rectangle ++(1,1);
\node[black] at (6.5,15.5) {$-3$};
\fill[figred!66] (7,15) rectangle ++(1,1);
\node[black] at (7.5,15.5) {$-4$};
\fill[figred!52] (8,15) rectangle ++(1,1);
\node[black] at (8.5,15.5) {$-3$};
\fill[figred!38] (9,15) rectangle ++(1,1);
\node[black] at (9.5,15.5) {$-2$};
\fill[figred!52] (10,15) rectangle ++(1,1);
\node[black] at (10.5,15.5) {$-3$};
\fill[figred!38] (11,15) rectangle ++(1,1);
\node[black] at (11.5,15.5) {$-2$};
\fill[figred!38] (12,15) rectangle ++(1,1);
\node[black] at (12.5,15.5) {$-2$};
\fill[figblue!28] (13,15) rectangle ++(1,1);
\node[black] at (13.5,15.5) {2};
\fill[figblue!28] (14,15) rectangle ++(1,1);
\node[black] at (14.5,15.5) {2};
\fill[figblue!44] (15,15) rectangle ++(1,1);
\node[black] at (15.5,15.5) {4};
\fill[figblue!20] (16,15) rectangle ++(1,1);
\node[black] at (16.5,15.5) {1};
\fill[white] (17,15) rectangle ++(1,1);
\node[black] at (17.5,15.5) {0};
\fill[figred!24] (18,15) rectangle ++(1,1);
\node[black] at (18.5,15.5) {$-1$};
\fill[figred!38] (19,15) rectangle ++(1,1);
\node[black] at (19.5,15.5) {$-2$};
\node[font=\scriptsize\bfseries] at (-0.5,15.5) {L};
\node[font=\scriptsize\bfseries] at (15.5,-0.5) {L};
\fill[figred!24] (0,16) rectangle ++(1,1);
\node[black] at (0.5,16.5) {$-1$};
\fill[figred!38] (1,16) rectangle ++(1,1);
\node[black] at (1.5,16.5) {$-2$};
\fill[white] (2,16) rectangle ++(1,1);
\node[black] at (2.5,16.5) {0};
\fill[figred!38] (3,16) rectangle ++(1,1);
\node[black] at (3.5,16.5) {$-2$};
\fill[white] (4,16) rectangle ++(1,1);
\node[black] at (4.5,16.5) {0};
\fill[figred!52] (5,16) rectangle ++(1,1);
\node[black] at (5.5,16.5) {$-3$};
\fill[figred!52] (6,16) rectangle ++(1,1);
\node[black] at (6.5,16.5) {$-3$};
\fill[figred!52] (7,16) rectangle ++(1,1);
\node[black] at (7.5,16.5) {$-3$};
\fill[figred!38] (8,16) rectangle ++(1,1);
\node[black] at (8.5,16.5) {$-2$};
\fill[figred!38] (9,16) rectangle ++(1,1);
\node[black] at (9.5,16.5) {$-2$};
\fill[figred!52] (10,16) rectangle ++(1,1);
\node[black] at (10.5,16.5) {$-3$};
\fill[figred!52] (11,16) rectangle ++(1,1);
\node[black] at (11.5,16.5) {$-3$};
\fill[figred!38] (12,16) rectangle ++(1,1);
\node[black] at (12.5,16.5) {$-2$};
\fill[figblue!20] (13,16) rectangle ++(1,1);
\node[black] at (13.5,16.5) {1};
\fill[figblue!36] (14,16) rectangle ++(1,1);
\node[black] at (14.5,16.5) {3};
\fill[figblue!20] (15,16) rectangle ++(1,1);
\node[black] at (15.5,16.5) {1};
\fill[figblue!44] (16,16) rectangle ++(1,1);
\node[black] at (16.5,16.5) {4};
\fill[figred!24] (17,16) rectangle ++(1,1);
\node[black] at (17.5,16.5) {$-1$};
\fill[figred!24] (18,16) rectangle ++(1,1);
\node[black] at (18.5,16.5) {$-1$};
\fill[figred!52] (19,16) rectangle ++(1,1);
\node[black] at (19.5,16.5) {$-3$};
\node[font=\scriptsize\bfseries] at (-0.5,16.5) {V};
\node[font=\scriptsize\bfseries] at (16.5,-0.5) {V};
\fill[figred!38] (0,17) rectangle ++(1,1);
\node[black] at (0.5,17.5) {$-2$};
\fill[figred!38] (1,17) rectangle ++(1,1);
\node[black] at (1.5,17.5) {$-2$};
\fill[figred!38] (2,17) rectangle ++(1,1);
\node[black] at (2.5,17.5) {$-2$};
\fill[figred!66] (3,17) rectangle ++(1,1);
\node[black] at (3.5,17.5) {$-4$};
\fill[figred!38] (4,17) rectangle ++(1,1);
\node[black] at (4.5,17.5) {$-2$};
\fill[figred!52] (5,17) rectangle ++(1,1);
\node[black] at (5.5,17.5) {$-3$};
\fill[figred!52] (6,17) rectangle ++(1,1);
\node[black] at (6.5,17.5) {$-3$};
\fill[figred!52] (7,17) rectangle ++(1,1);
\node[black] at (7.5,17.5) {$-3$};
\fill[figred!52] (8,17) rectangle ++(1,1);
\node[black] at (8.5,17.5) {$-3$};
\fill[figred!52] (9,17) rectangle ++(1,1);
\node[black] at (9.5,17.5) {$-3$};
\fill[figred!24] (10,17) rectangle ++(1,1);
\node[black] at (10.5,17.5) {$-1$};
\fill[figred!52] (11,17) rectangle ++(1,1);
\node[black] at (11.5,17.5) {$-3$};
\fill[figred!52] (12,17) rectangle ++(1,1);
\node[black] at (12.5,17.5) {$-3$};
\fill[white] (13,17) rectangle ++(1,1);
\node[black] at (13.5,17.5) {0};
\fill[white] (14,17) rectangle ++(1,1);
\node[black] at (14.5,17.5) {0};
\fill[white] (15,17) rectangle ++(1,1);
\node[black] at (15.5,17.5) {0};
\fill[figred!24] (16,17) rectangle ++(1,1);
\node[black] at (16.5,17.5) {$-1$};
\fill[figblue!60] (17,17) rectangle ++(1,1);
\node[black] at (17.5,17.5) {6};
\fill[figblue!36] (18,17) rectangle ++(1,1);
\node[black] at (18.5,17.5) {3};
\fill[figblue!20] (19,17) rectangle ++(1,1);
\node[black] at (19.5,17.5) {1};
\node[font=\scriptsize\bfseries] at (-0.5,17.5) {F};
\node[font=\scriptsize\bfseries] at (17.5,-0.5) {F};
\fill[figred!38] (0,18) rectangle ++(1,1);
\node[black] at (0.5,18.5) {$-2$};
\fill[figred!38] (1,18) rectangle ++(1,1);
\node[black] at (1.5,18.5) {$-2$};
\fill[figred!38] (2,18) rectangle ++(1,1);
\node[black] at (2.5,18.5) {$-2$};
\fill[figred!52] (3,18) rectangle ++(1,1);
\node[black] at (3.5,18.5) {$-3$};
\fill[figred!38] (4,18) rectangle ++(1,1);
\node[black] at (4.5,18.5) {$-2$};
\fill[figred!52] (5,18) rectangle ++(1,1);
\node[black] at (5.5,18.5) {$-3$};
\fill[figred!38] (6,18) rectangle ++(1,1);
\node[black] at (6.5,18.5) {$-2$};
\fill[figred!52] (7,18) rectangle ++(1,1);
\node[black] at (7.5,18.5) {$-3$};
\fill[figred!38] (8,18) rectangle ++(1,1);
\node[black] at (8.5,18.5) {$-2$};
\fill[figred!24] (9,18) rectangle ++(1,1);
\node[black] at (9.5,18.5) {$-1$};
\fill[figblue!28] (10,18) rectangle ++(1,1);
\node[black] at (10.5,18.5) {2};
\fill[figred!38] (11,18) rectangle ++(1,1);
\node[black] at (11.5,18.5) {$-2$};
\fill[figred!38] (12,18) rectangle ++(1,1);
\node[black] at (12.5,18.5) {$-2$};
\fill[figred!24] (13,18) rectangle ++(1,1);
\node[black] at (13.5,18.5) {$-1$};
\fill[figred!24] (14,18) rectangle ++(1,1);
\node[black] at (14.5,18.5) {$-1$};
\fill[figred!24] (15,18) rectangle ++(1,1);
\node[black] at (15.5,18.5) {$-1$};
\fill[figred!24] (16,18) rectangle ++(1,1);
\node[black] at (16.5,18.5) {$-1$};
\fill[figblue!36] (17,18) rectangle ++(1,1);
\node[black] at (17.5,18.5) {3};
\fill[figblue!68] (18,18) rectangle ++(1,1);
\node[white] at (18.5,18.5) {7};
\fill[figblue!28] (19,18) rectangle ++(1,1);
\node[black] at (19.5,18.5) {2};
\node[font=\scriptsize\bfseries] at (-0.5,18.5) {Y};
\node[font=\scriptsize\bfseries] at (18.5,-0.5) {Y};
\fill[figred!38] (0,19) rectangle ++(1,1);
\node[black] at (0.5,19.5) {$-2$};
\fill[figred!52] (1,19) rectangle ++(1,1);
\node[black] at (1.5,19.5) {$-3$};
\fill[figred!38] (2,19) rectangle ++(1,1);
\node[black] at (2.5,19.5) {$-2$};
\fill[figred!66] (3,19) rectangle ++(1,1);
\node[black] at (3.5,19.5) {$-4$};
\fill[figred!52] (4,19) rectangle ++(1,1);
\node[black] at (4.5,19.5) {$-3$};
\fill[figred!38] (5,19) rectangle ++(1,1);
\node[black] at (5.5,19.5) {$-2$};
\fill[figred!66] (6,19) rectangle ++(1,1);
\node[black] at (6.5,19.5) {$-4$};
\fill[figred!66] (7,19) rectangle ++(1,1);
\node[black] at (7.5,19.5) {$-4$};
\fill[figred!52] (8,19) rectangle ++(1,1);
\node[black] at (8.5,19.5) {$-3$};
\fill[figred!38] (9,19) rectangle ++(1,1);
\node[black] at (9.5,19.5) {$-2$};
\fill[figred!38] (10,19) rectangle ++(1,1);
\node[black] at (10.5,19.5) {$-2$};
\fill[figred!52] (11,19) rectangle ++(1,1);
\node[black] at (11.5,19.5) {$-3$};
\fill[figred!52] (12,19) rectangle ++(1,1);
\node[black] at (12.5,19.5) {$-3$};
\fill[figred!24] (13,19) rectangle ++(1,1);
\node[black] at (13.5,19.5) {$-1$};
\fill[figred!52] (14,19) rectangle ++(1,1);
\node[black] at (14.5,19.5) {$-3$};
\fill[figred!38] (15,19) rectangle ++(1,1);
\node[black] at (15.5,19.5) {$-2$};
\fill[figred!52] (16,19) rectangle ++(1,1);
\node[black] at (16.5,19.5) {$-3$};
\fill[figblue!20] (17,19) rectangle ++(1,1);
\node[black] at (17.5,19.5) {1};
\fill[figblue!28] (18,19) rectangle ++(1,1);
\node[black] at (18.5,19.5) {2};
\fill[figblue!100] (19,19) rectangle ++(1,1);
\node[white] at (19.5,19.5) {11};
\node[font=\scriptsize\bfseries] at (-0.5,19.5) {W};
\node[font=\scriptsize\bfseries] at (19.5,-0.5) {W};
\draw[gray!40, very thin, step=1] (0,0) grid (20,20);
\draw[black, thick] (0,0) rectangle (1,1);
\draw[black, thick] (1,1) rectangle (6,6);
\draw[black, thick] (6,6) rectangle (10,10);
\draw[black, thick] (10,10) rectangle (13,13);
\draw[black, thick] (13,13) rectangle (17,17);
\draw[black, thick] (17,17) rectangle (20,20);
\draw[black] (0,0) rectangle (20,20);
\end{tikzpicture}

\caption{BLOSUM62, the most widely used substitution matrix. The amino acids
are ordered by chemical similarity; the black squares mark the groups
(C \,|\, small: S T P A G \,|\, acidic and amide: N D E Q \,|\, basic: H R K
\,|\, hydrophobic: M I L V \,|\, aromatic: F Y W). Blue = positive (more often
than by chance), red = negative.}
\label{fig:blosum62}
\end{figure}

\Cref{fig:blosum62} shows the full BLOSUM62 matrix. A few things to read off:
\begin{itemize}
  \item The \textbf{diagonal} is always positive. It is largest for rare,
    hard-to-replace amino acids: tryptophan (W, 11) and cysteine (C, 9).
    Common amino acids with many acceptable substitutes have small diagonal
    values (A, S: 4).
  \item Positive off-diagonal entries cluster in the \textbf{group squares}:
    I--L--V--M, F--Y--W, D--E--N--Q, K--R. These are the conservative
    substitutions from \cref{ch:scoring}: similar size and chemistry.
  \item Replacing a hydrophobic by a charged amino acid is strongly penalised
    (the red areas between the groups).
\end{itemize}

\paragraph{Conformational aspects}
The slides check some entries against chemistry:
\begin{itemize}
  \item leucine $\to$ isoleucine: $+2$ --- both large hydrophobic, almost the
    same shape;
  \item aspartate $\to$ glutamate: $+2$ --- both negatively charged, glutamate
    is just one $\mathrm{CH_2}$ group longer;
  \item tyrosine $\to$ phenylalanine: $+3$ --- both aromatic rings; tyrosine has
    only an extra $\mathrm{OH}$ group;
  \item alanine $\to$ valine: $0$ --- both small and hydrophobic, yet only ``as
    often as chance''. The matrix is learned from data, not from chemistry
    textbooks, and it does not always match our intuition.
\end{itemize}

\section{PAM versus BLOSUM, and How to Choose}
\label{sec:choose-matrix}

\begin{center}
\small
\begin{tabular}{@{}p{0.45\linewidth}p{0.45\linewidth}@{}}
\toprule
\textbf{PAM} & \textbf{BLOSUM} \\
\midrule
made from \emph{global} alignments of closely related proteins &
made from \emph{local} alignments of only highly conserved subsequences \\
long-range PAMs are \emph{extrapolated} from 1PAM &
every matrix is based on \emph{observations}, no extrapolation \\
higher number = more distant &
higher number = more closely related \\
\bottomrule
\end{tabular}
\end{center}

Both have the same goal, and for every PAM there is a BLOSUM ``partner'' that
leads to similar results. Different sources pair them differently; the
slides list three versions. Regl recommends the following one, because it is
based on matrices with similar \emph{relative entropy}:

\begin{figure}[htbp]
\centering
\begin{tikzpicture}[x=1cm, y=1cm, font=\small]
  \draw[-{Stealth}, very thick, gray] (0,0) -- (12.4,0);
  \node[anchor=south west, font=\scriptsize, gray] at (0,0.05) {low divergence (close relatives)};
  \node[anchor=south east, font=\scriptsize, gray] at (12.4,0.05) {high divergence (distant relatives)};
  \foreach \x/\p/\b in {1.2/100/90, 3.6/120/80, 6.0/160/60, 8.4/200/52, 10.8/250/45} {
    \draw[thick] (\x,-0.12) -- (\x,0.12);
    \node[figblue, font=\small\bfseries] at (\x,-0.55) {PAM\p};
    \node[figred, font=\small\bfseries] at (\x,-1.05) {BLOSUM\b};
  }
  \node[anchor=east, figblue, font=\scriptsize] at (-0.1,-0.55) {number grows $\rightarrow$};
  \node[anchor=east, figred, font=\scriptsize] at (-0.1,-1.05) {number shrinks $\rightarrow$};
\end{tikzpicture}
\caption{Corresponding PAM and BLOSUM matrices along the evolutionary
distance axis (the pairing recommended in the lecture). PAM numbers increase
with distance, BLOSUM numbers decrease.}
\label{fig:pam-blosum-axis}
\end{figure}

\begin{notebox}[title={Background: what ``relative entropy'' measures}]
The relative entropy $H = \sum_{a \le b} q_{ab}\, S(a,b)$ (with $S$ in bits) is
the average score per aligned pair in \emph{truly related} sequences --- how
much information, on average, one aligned position carries to distinguish
``related'' from ``random''. A matrix for close relatives is ``sharp'' (high
$H$: matches are expected, mismatches are suspicious); a matrix for distant
relatives is ``blunt'' (low $H$). Matrices with the same $H$ behave similarly,
which is why they make good partners: BLOSUM62 and PAM160, for example, both
carry about 0.7 bits per position.
\end{notebox}

\subsection{A Rule of Thumb}

First check the \textbf{sequence similarity} of the sequences you want to
align, then pick the matrix (\cref{tab:rule-of-thumb}). For example, at about
30\,\% similarity, take BLOSUM60 or PAM160; the two give roughly the same
answers.

\begin{table}[htbp]
\centering
\small
\begin{tabular}{@{}ccc@{}}
\toprule
\textbf{Similarity (\%)} & \textbf{PAM} & \textbf{BLOSUM} \\
\midrule
43 & 100 & 90 \\
38 & 120 & 80 \\
30 & 160 & 60 \\
25 & 200 & 50 \\
20 & 250 & 45 \\
\bottomrule
\end{tabular}
\caption{Rule of thumb for selecting a matrix (it corresponds to the pairing in
\cref{fig:pam-blosum-axis}).}
\label{tab:rule-of-thumb}
\end{table}

\subsection{Which One? BLOSUM}

\begin{notebox}[title={Regl's advice --- likely exam material}]
``Clearly BLOSUM!'' Both have pros and cons, but PAM is around mainly for
historical reasons (``maybe people use PAM because they were always using
PAM'') and has the drawbacks listed in \cref{ch:pam}. NCBI agrees: ``A number
of tests \ldots{} suggest that the BLOSUM matrices \ldots{} are superior to
the PAM matrices for detecting biological relationships'' (NCBI's BLAST
tutorial).
\end{notebox}

\paragraph{In the tools}
In BLAST (\cref{sec:blast}), the matrix is chosen under ``Algorithm parameters
$\to$ Scoring parameters''. The choice is surprisingly small: PAM30, PAM70,
PAM250, BLOSUM45, BLOSUM50, BLOSUM62 (the default), BLOSUM80 and BLOSUM90.
Regl would drop the PAMs, offer more BLOSUMs and allow uploading your own
matrix. ChimeraX (\cref{sec:chimerax}) offers many more PAMs and BLOSUMs: its
3D ``matchmaker'' starts with a \emph{sequence} alignment, because regions with
high sequence similarity are the ones it tries to superimpose first --- and a
sequence alignment needs a scoring matrix. The standard matrices can be
downloaded from \url{ftp://ftp.ncbi.nih.gov/blast/matrices/}.

\subsection{Standard Matrix or ``Bring Your Own''?}

All standard matrices were calculated from ``general'' proteins from across
the databases. If you work on a special class of proteins, the honest answer is:
calculate your own.
\begin{itemize}
  \item \textbf{PAM}: possible in theory, but in practice nobody does it ---
    dividing all sequences into families and finding the minimum-mutation
    trees is cumbersome, and PAM is disfavoured anyway.
  \item \textbf{BLOSUM}: possible in theory and in practice. Selecting the
    sequences is the main work; after that, computing the matrix is ``a piece of
    cake'' --- programs are easily found or written (\cref{sec:blosum-steps}).
\end{itemize}

\begin{figure}[htbp]
\centering
\begin{tikzpicture}[font=\small,
    q/.style={box, fill=figamber!12, text width=3.4cm},
    a/.style={box, fill=figgreen!12, text width=7.2cm, align=left}]
  \node[q] (type) at (0,0) {What kind of proteins?};
  \node[a, anchor=west] (own) at (3.4,0) {special class (e.g.\ transmembrane proteins, DNA clamps): take blocks from \emph{that} class, build your own BLOSUM};
  \node[q] (sim) at (0,-2.7) {How related are the organisms?};
  \node[a, anchor=west] (close) at (3.4,-1.7) {closely related (human, chimp, orangutan): BLOSUM90 or 100};
  \node[a, anchor=west] (far) at (3.4,-2.7) {distantly related: BLOSUM50 or 30};
  \node[a, anchor=west] (dk) at (3.4,-3.7) {a mixture of both, or don't know: \mbox{\textbf{BLOSUM62}}};
  \draw[flow] (type) -- (own);
  \draw[flow] (type) -- node[right, lbl] {then} (sim);
  \draw[flow] (sim.east) -- ++(0.35,0) |- (close.west);
  \draw[flow] (sim.east) -- (far.west);
  \draw[flow] (sim.east) -- ++(0.35,0) |- (dk.west);
\end{tikzpicture}
\caption{Choosing a scoring matrix: first orient it on the sequence type of
your project, then on the sequence similarity.}
\label{fig:choose-matrix}
\end{figure}

BLOSUM62 is what people use when they are not really sure which one to take ---
it is ``somehow in the middle'' and was built from a collection of general
proteins.

\subsection{Alternatives}

There are many more matrices, each mainly or only used for special purposes:
\begin{center}
\small
\begin{tabular}{@{}lp{0.72\linewidth}@{}}
\toprule
\textbf{Matrix} & \textbf{Idea} \\
\midrule
MD   & re-calculated PAM using up-to-date sequences \\
VTML & enhanced PAM with more sophisticated statistics behind the log-odds \\
PfSSM & special BLOSUM for the AT-rich genome of \emph{Plasmodium falciparum}
  (the malaria parasite) \\
CBM, CCF & special BLOSUMs for transmembrane proteins \\
JTT, PHAT, SLIM, \ldots & for special classes of proteins, mostly derived from BLOSUM \\
\bottomrule
\end{tabular}
\end{center}
All of them perform better in their field of application --- but none of them
is available in BLAST. Everybody uses BLOSUM (and some still use PAM).

\section{Outlook: More Ideas on Scoring and Alignment}
\label{sec:more-alignment}

So far, every column of an alignment is scored with the \emph{same} matrix,
and alignments compare two sequences. The last slides of the scoring part
point beyond that; the details follow in later chapters.

\begin{itemize}
  \item \textbf{Profile-based scoring}: the alignment concentrates on
    functional sites across the whole sequence. This needs highly specialised
    algorithms and clever gap-scoring schemes.
  \item \textbf{Aligning groups of sequences}: like profile-based scoring, but
    first a profile (usually via majority vote, as in similarization) is
    constructed and used as a representative sequence (\cref{ch:msa}).
  \item \textbf{Position-specific scoring}: frequencies and log-odds are
    collected for \emph{each position} of the alignment separately --- the
    functional site gets its own, much stricter scores. This is used in some
    multiple sequence alignment software (and is the idea behind HMM
    profiles, \cref{ch:hmm}).
\end{itemize}

\paragraph{One more type of alignment: overlaps}
Sequencing machines cannot read a whole genome in one go; they produce
millions of short \emph{reads} (\cref{ch:bio2}). In \textbf{genome assembly},
the reads are put together like a puzzle by finding pairs of reads whose ends
overlap (\cref{fig:overlap}). Neither a global alignment (the whole of both
reads) nor a local one (any region) is right here: we need an alignment that
\emph{ignores gaps at the start of read $i$ and at the end of read $k$}.
The changes to the algorithms are quite straightforward (\cref{ch:global}).

\begin{figure}[htbp]
\centering
\begin{tikzpicture}[x=0.36cm, y=1cm, font=\ttfamily\small]
  \node[anchor=east, font=\small\rmfamily] at (-0.5,0.5) {read $i$};
  \node[anchor=east, font=\small\rmfamily] at (-0.5,0) {read $k$};
  \foreach \c [count=\j] in {T,G,A,C,C,T,A,G,G,A,T,C,C,A,G,T} \node at (\j,0.5) {\c};
  \foreach \c [count=\j] in {G,G,A,T,C,C,A,G,T,T,A,C,G,A,A,C} \node at (\j+7,0) {\c};
  \fill[figgreen, opacity=0.18] (7.5,-0.25) rectangle (16.5,0.75);
  \draw[decorate, decoration={brace, amplitude=4pt, mirror}] (0.5,-0.35) -- (7.4,-0.35)
    node[midway, below=4pt, font=\scriptsize\rmfamily] {free end gap in read $k$};
  \draw[decorate, decoration={brace, amplitude=4pt}] (16.6,0.85) -- (23.5,0.85)
    node[midway, above=4pt, font=\scriptsize\rmfamily] {free end gap in read $i$};
  \node[font=\scriptsize\rmfamily, figgreen!60!black] at (12,-0.65) {matching overlap};
\end{tikzpicture}
\caption{Overlap alignment of two sequencing reads: the end of read $i$
matches the start of read $k$. The overhanging parts must not be penalised as
gaps.}
\label{fig:overlap}
\end{figure}

``This is not the end'': there are also \emph{repeated matches} and
\emph{nonlinear matches}, all implemented by small changes to the alignment
algorithms.

\section{Chapter Checklist}

\begin{itemize}
  \item why similarity comes in blocks (functional sites mutate slowly);
  \item what BLOSUM stands for, what a block is (same length, no indels) and
    what the $x$ in BLOSUM$x$ means --- and that it runs opposite to PAM;
  \item elimination, similarization and weighting of close relatives;
  \item compute $C$ (by counting and with the shortcut), $T$, $Q$, $p$, $E$,
    $S$ and the rounded BLOSUM for a small block;
  \item read BLOSUM62: diagonal, conservative substitutions, examples;
  \item PAM vs BLOSUM: global/extrapolated vs local/observed; the
    correspondence and the rule of thumb; why BLOSUM is recommended;
  \item how to choose a matrix (sequence type, similarity, BLOSUM62 if unsure);
    when to build your own;
  \item profile-based and position-specific scoring; overlap alignment for
    genome assembly.
\end{itemize}

% ==========================================================
%  CHAPTER 9 — TOOLS IN PRACTICE (ADD-ONS)
%  Lecture VO04 (A. Regl), add-on slides 1--18 (VO04_scoring-2_add-ons.pdf)
%  Recording VO04 00:00--31:00 (live demo of UniProt, BLAST, AlphaFold, ChimeraX)
% ==========================================================
\chapter{Tools in Practice: From Sequence to Structure}
\label{ch:tools}

\lectureref{VO04, ``add-ons'' deck (18 slides). Recording 00:00--31:00: a live
demo of UniProt, BLAST, AlphaFold and ChimeraX with human and bovine insulin,
followed by a real research case (METTL3 in bladder cancer).}

The previous chapters explained \emph{how} sequences are compared. This chapter
shows the tools you actually click on. They come up again and again in the rest
of the course, and Regl's question about BLAST is a classic in exams. The
running example asks a simple question: \emph{how similar are human and bovine
insulin?} Insulin has the same function in humans and in cows, but over time
mutations have appeared.

\begin{figure}[htbp]
\centering
\begin{tikzpicture}[node distance=9mm, font=\small,
    tool/.style={box, text width=2.5cm, minimum height=1.25cm}]
  \node[tool] (u) {\textbf{UniProt}\\{\scriptsize get the sequences}\\[-2pt]{\scriptsize (FASTA)}};
  \node[tool, right=of u] (b) {\textbf{BLAST}\\{\scriptsize compare the}\\[-2pt]{\scriptsize sequences}};
  \node[tool, right=of b] (a) {\textbf{AlphaFold}\\{\scriptsize predict the 3D}\\[-2pt]{\scriptsize structures}};
  \node[tool, right=of a] (c) {\textbf{ChimeraX}\\{\scriptsize superimpose and}\\[-2pt]{\scriptsize compare in 3D}};
  \node[tool, below=7mm of a, fill=gray!10] (p) {\textbf{PDB}\\{\scriptsize or download}\\[-2pt]{\scriptsize known structures}};
  \draw[flow] (u) -- (b); \draw[flow] (b) -- (a); \draw[flow] (a) -- (c);
  \draw[flow, dashed] (p) -| (c);
\end{tikzpicture}
\caption{The workflow of the lecture demo --- and, as we will see, of a real
research project. Sequence first, then structure.}
\label{fig:tool-workflow}
\end{figure}

\section{Databases: Where Sequences and Structures Live}
\label{sec:databases}

\begin{center}
\small
\begin{tabular}{@{}p{0.2\linewidth}p{0.38\linewidth}p{0.34\linewidth}@{}}
\toprule
\textbf{Database} & \textbf{Contains} & \textbf{Notes} \\
\midrule
\textbf{UniProt} (\url{uniprot.org}) &
\textbf{sequences}, mainly proteins, with everything known about them &
run by EMBL-EBI together with the Swiss Institute of Bioinformatics and PIR \\
\textbf{PDB} (Protein Data Bank) &
experimentally determined \textbf{3D structures} &
entries have 4-character IDs such as \texttt{1AXC} (\cref{ch:alignment}) \\
\textbf{NCBI} &
the US National Center for Biotechnology Information: sequence databases
(e.g.\ GenBank) and tools &
home of \textbf{BLAST} \\
\textbf{AlphaFold DB} &
about 200 million \textbf{predicted} structures &
computed from UniProt sequences (\cref{sec:alphafold}) \\
\bottomrule
\end{tabular}
\end{center}

\begin{notebox}[title={Don't mix them up}]
UniProt contains \emph{sequences}, the PDB contains \emph{structures}. (And
neither of them is run by NCBI.)
\end{notebox}

\paragraph{Getting the insulin sequences}
Searching UniProt for ``insulin'' gives about 292,000 results --- restricted to
human, still nearly 12,000, because insulin appears in many configurations,
bound to many different molecules (the same overrepresentation problem BLOSUM
has to deal with, \cref{ch:blosum}). The first hit is the one we want: human
insulin, 110 amino acids. On its page you find everything that is known about
it; we only need the sequence, so we download it in FASTA format
(\cref{ch:alignment}):

\alnblock{%
>sp|P01308|INS\_HUMAN Insulin OS=Homo sapiens OX=9606 GN=INS PE=1 SV=1\\
MALWMRLLPLLALLALWGPDPAAAFVNQHLCGSHLVEALYLVCGERGFFYTPKTRREAED\\
LQVGQVELGGGPGAGSLQPLALEGSLQKRGIVEQCCTSICSLYQLENYCN}

The first line is the identifier and some information (here: the UniProt
accession \texttt{P01308}, the organism, the gene name), the following lines
the sequence, 60 letters per line. The same for bovine insulin (105 amino acids).

\section{BLAST}
\label{sec:blast}

\begin{defbox}[title={BLAST = Basic Local Alignment Search Tool}]
NCBI's tool for comparing a sequence (the \textbf{query}) against all
sequences of a database, or against a second sequence (the \textbf{subject}).
It returns the best \textbf{local} alignments (\cref{ch:local}) together with
statistics on how significant they are. Variants: \texttt{blastp} for proteins,
\texttt{blastn} for nucleotides, and a few more.
\end{defbox}

\begin{notebox}[title={Exam classic}]
Regl once asked a student in a bachelor exam whether BLAST produces local or
global alignments, and she did not know. The answer is in the name: Basic
\textbf{Local} Alignment Search Tool.
\end{notebox}

Normally BLAST checks your sequence against a whole database (``which known
proteins look like mine?''). For two given sequences, tick ``Align two or more
sequences'', paste or upload the two FASTA files and run protein BLAST. Under
``Algorithm parameters'' you can choose the scoring matrix (default BLOSUM62,
\cref{sec:choose-matrix}).

\subsection{Reading the Output}

\begin{center}
\ttfamily\footnotesize
\begin{tabular}{@{}l@{\;}r@{\;\;}l@{\;\;}r@{}}
Query & 1  & MALWMRLLPLLALLALWGPDPAAAFVNQHLCGSHLVEALYLVCGERGFFYTPKTRREAED & 60 \\
      &    & MALW~RL~PLLALLALW~P~PA~AFVNQHLCGSHLVEALYLVCGERGFFYTPK~RRE~E~ & \\
Sbjct & 1  & MALWTRLRPLLALLALWPPPPARAFVNQHLCGSHLVEALYLVCGERGFFYTPKARREVEG & 60 \\[1ex]
Query & 61 & LQVGQVELGGGPGAGSLQPLALEGSLQKRGIVEQCCTSICSLYQLENYCN & 110 \\
      &    & ~QVG~+EL~~~~~AG~~~~~~LEG~~QKRGIVEQCC~S+CSLYQLENYCN & \\
Sbjct & 61 & PQVGALEL-{}-{}-{}-{}-AGGPGAGGLEGPPQKRGIVEQCCASVCSLYQLENYCN & 105 \\
\end{tabular}
\end{center}

\begin{itemize}
  \item The first line is the query, the third line the subject; the numbers
    are positions in each sequence.
  \item The \textbf{middle line} shows the letter if both are identical, a
    \textbf{+} if they are different but \emph{similar} (positive score in the
    scoring matrix, e.g.\ V--L and I--V: $+1$ and $+3$ in BLOSUM62), and a
    \textbf{blank} if they are dissimilar.
  \item A \textbf{minus sign} in a sequence is a gap (insertion or deletion).
\end{itemize}
One look at the middle line tells you how similar the sequences are: lots of
blanks, not similar. Here: 84 of 110 positions identical (76\,\%), 86 similar
(78\,\%, ``positives''), and one gap of 5 (BLAST reports exactly these three
numbers above every alignment).

\begin{exbox}[title={Where are the differences --- and why there}]
The 26 differences are not spread evenly (\cref{fig:insulin}). To understand
the pattern you need one piece of biology: the 110 amino acids are
\emph{preproinsulin}, the precursor. The first 24 (the \emph{signal peptide})
only direct the new protein into the cell's export route and are cut off
right away. Then the middle part,
the \textbf{C-peptide}, is cut out too. The actual hormone consists only of the
\textbf{B chain} (30 amino acids) and the \textbf{A chain} (21), held together by
disulfide bonds between their cysteines.

The differences sit where they do not matter: 18 in the C-peptide (which is
thrown away) and 5 in the signal peptide, but only \textbf{3 in the 51 amino
acids of the hormone itself}. Functional parts are conserved
(\cref{fig:functional-blocks}) --- and that is why insulin from cattle (and
pigs) could be used to treat diabetes for about 60 years, until human insulin
could be produced by genetically modified bacteria in the 1980s.
\end{exbox}

\begin{figure}[htbp]
\centering
\begin{tikzpicture}[x=0.12cm, y=1cm, font=\small]
\fill[gray!30] (0.0,0) rectangle (24.0,0.55);
\node[font=\scriptsize] at (12.0,0.275) {signal peptide};
\fill[figblue!45] (24.0,0) rectangle (54.0,0.55);
\node[font=\scriptsize] at (39.0,0.275) {B chain};
\fill[white] (54.0,0) rectangle (56.0,0.55);
\fill[figamber!40] (56.0,0) rectangle (87.0,0.55);
\node[font=\scriptsize] at (71.5,0.275) {C-peptide};
\fill[white] (87.0,0) rectangle (89.0,0.55);
\fill[figblue!45] (89.0,0) rectangle (110.0,0.55);
\node[font=\scriptsize] at (99.5,0.275) {A chain};
\draw[thick] (0,0) rectangle (110,0.55);
\draw[figred, very thick] (4.5,0.65) -- (4.5,1.0);
\draw[figred, very thick] (7.5,0.65) -- (7.5,1.0);
\draw[figred, very thick] (17.5,0.65) -- (17.5,1.0);
\draw[figred, very thick] (19.5,0.65) -- (19.5,1.0);
\draw[figred, very thick] (22.5,0.65) -- (22.5,1.0);
\draw[figred, very thick] (53.5,0.65) -- (53.5,1.0);
\draw[figred, very thick] (57.5,0.65) -- (57.5,1.0);
\draw[figred, very thick] (59.5,0.65) -- (59.5,1.0);
\draw[figred, very thick] (60.5,0.65) -- (60.5,1.0);
\draw[figred, very thick] (64.5,0.65) -- (64.5,1.0);
\draw[figred, very thick] (65.5,0.65) -- (65.5,1.0);
\draw[figred, very thick] (68.5,0.65) -- (68.5,1.0);
\fill[figred] (68.5,1.0) circle (1.3pt);
\draw[figred, very thick] (69.5,0.65) -- (69.5,1.0);
\fill[figred] (69.5,1.0) circle (1.3pt);
\draw[figred, very thick] (70.5,0.65) -- (70.5,1.0);
\fill[figred] (70.5,1.0) circle (1.3pt);
\draw[figred, very thick] (71.5,0.65) -- (71.5,1.0);
\fill[figred] (71.5,1.0) circle (1.3pt);
\draw[figred, very thick] (72.5,0.65) -- (72.5,1.0);
\fill[figred] (72.5,1.0) circle (1.3pt);
\draw[figred, very thick] (75.5,0.65) -- (75.5,1.0);
\draw[figred, very thick] (76.5,0.65) -- (76.5,1.0);
\draw[figred, very thick] (77.5,0.65) -- (77.5,1.0);
\draw[figred, very thick] (78.5,0.65) -- (78.5,1.0);
\draw[figred, very thick] (79.5,0.65) -- (79.5,1.0);
\draw[figred, very thick] (80.5,0.65) -- (80.5,1.0);
\draw[figred, very thick] (84.5,0.65) -- (84.5,1.0);
\draw[figred, very thick] (85.5,0.65) -- (85.5,1.0);
\draw[figred, very thick] (96.5,0.65) -- (96.5,1.0);
\draw[figred, very thick] (98.5,0.65) -- (98.5,1.0);
\node[font=\scriptsize, gray] at (0.5,-0.22) {1};
\node[font=\scriptsize, gray] at (23.5,-0.22) {24};
\node[font=\scriptsize, gray] at (53.5,-0.22) {54};
\node[font=\scriptsize, gray] at (86.5,-0.22) {87};
\node[font=\scriptsize, gray] at (109.5,-0.22) {110};
\node[anchor=west, font=\scriptsize, figred] at (111,0.85) {differences};
\draw[decorate, decoration={brace, amplitude=4pt, mirror}, thick, figblue] (24,-0.45) -- (54,-0.45) node[midway, below=4pt, font=\scriptsize, figblue] {1 difference};
\draw[decorate, decoration={brace, amplitude=4pt, mirror}, thick, figamber!80!black] (56,-0.45) -- (87,-0.45) node[midway, below=4pt, font=\scriptsize, figamber!80!black] {18 differences incl.\ 5 gaps};
\draw[decorate, decoration={brace, amplitude=4pt, mirror}, thick, figblue] (89,-0.45) -- (110,-0.45) node[midway, below=4pt, font=\scriptsize, figblue] {2 differences};
\end{tikzpicture}

\caption{Human vs bovine preproinsulin (positions of the human sequence). Each
red tick is a position where the BLAST alignment shows a difference (dots: the
5 gaps). The mature hormone (blue: B and A chain) is almost identical.}
\label{fig:insulin}
\end{figure}

\section{AlphaFold}
\label{sec:alphafold}

If you search for AlphaFold, you get two different things:
\begin{itemize}
  \item the \textbf{AlphaFold Protein Structure Database}: pre-computed
    structures for practically all sequences in UniProt and other sources ---
    about 200 million. Look here first;
  \item the \textbf{AlphaFold Server} (\url{alphafoldserver.com}): paste your
    own sequence and get a prediction from scratch. Needed for sequences that
    are not in the database --- for example a mutant (\cref{sec:mettl3}).
\end{itemize}
The predicted structure has the same file format as an experimental PDB entry,
so all structure tools can use both.

\subsection{How Much to Trust a Prediction}

AlphaFold colours its prediction by a per-residue confidence score (pLDDT,
from 0 to 100, \cref{fig:plddt}).

\begin{figure}[htbp]
\centering
\begin{tikzpicture}[x=1cm, y=1cm, font=\small]
  \definecolor{afvh}{HTML}{0053D6}\definecolor{afh}{HTML}{65CBF3}
  \definecolor{afl}{HTML}{FFDB13}\definecolor{afvl}{HTML}{FF7D45}
  \fill[afvl] (0,0) rectangle (5,0.5);
  \fill[afl] (5,0) rectangle (7,0.5);
  \fill[afh] (7,0) rectangle (9,0.5);
  \fill[afvh] (9,0) rectangle (10,0.5);
  \draw (0,0) rectangle (10,0.5);
  \foreach \x/\l in {0/0, 5/50, 7/70, 9/90, 10/100} \node[below, font=\scriptsize] at (\x,0) {\l};
  \node[above, font=\scriptsize, align=center] at (2.5,0.5) {very low\\(likely disordered)};
  \node[above, font=\scriptsize] at (6,0.5) {low};
  \node[above, font=\scriptsize] at (8,0.5) {confident};
  \node[above, font=\scriptsize, align=center] at (9.5,0.5) {very\\high};
  \node[anchor=east, font=\scriptsize] at (-0.2,0.25) {pLDDT};
\end{tikzpicture}
\caption{AlphaFold's confidence colour code: dark blue = high confidence,
yellow and orange = very unsure.}
\label{fig:plddt}
\end{figure}

The typical picture --- for insulin as for almost any protein --- is a
confident (blue) \textbf{compact core} and uncertain (orange)
\textbf{loops} sticking out. This makes physical sense: in a compact
structure, every amino acid is held in place by many hydrogen bonds and van der
Waals interactions with its neighbours, so its position is well determined. A
loop on the surface could be in many positions (``this one, or this one, or
this one''), and the predictor is honestly unsure. In the insulin prediction,
even a small $\alpha$-helix is coloured orange: AlphaFold does not really know
whether it is a helix at all.

Each run produces \textbf{five predictions}, which are similar but slightly
different. Comparing them is a good way to see how stable a prediction is.

\figureplaceholder{4.5cm}{AlphaFold prediction of human insulin, coloured by confidence}%
{Search the AlphaFold DB (\texttt{alphafold.ebi.ac.uk}) for \texttt{P01308}
(human insulin) and take a screenshot of the 3D view: blue core, orange loops.
Alternatively paste the sequence into alphafoldserver.com.}%
{AlphaFold's prediction of human (pre-pro)insulin: confident core, uncertain
loops.}{fig:af-insulin}

\section{ChimeraX}
\label{sec:chimerax}

\textbf{ChimeraX} (\url{cgl.ucsf.edu/chimerax}), maintained by UCSF, is
state-of-the-art molecular visualisation and analysis software: 3D matching of
several proteins and ``tons of chemical calculations''. It is free --- Regl's
recommendation: ``Just download it, you will use it over and over again.''
Don't confuse it with \emph{Chimera}, its predecessor.

\paragraph{Superimposing two structures with matchmaker}
Open both structures (from AlphaFold or the PDB). They are already roughly
aligned, but not well. The \emph{matchmaker} tool superimposes them. It asks
which structure is the reference (with two it does not matter; with ten, all
others are oriented on the one you pick). Internally, matchmaker first computes
a \emph{sequence} alignment --- which is why it needs a scoring matrix
(\cref{sec:choose-matrix}) --- and then fits the 3D coordinates, starting with
the regions where the sequences agree best.

For the two insulins, the best fit is achieved in the middle part of the
prediction (a ``curly'' stretch that is neither a clear helix nor a sheet); the
two $\alpha$-helices at the ends end up in different positions, and the loops
differ as well.

\figureplaceholder{4.5cm}{ChimeraX matchmaker: human vs bovine insulin superimposed}%
{To make it yourself in ChimeraX: \texttt{alphafold fetch P01308},
\texttt{alphafold fetch P01317} (bovine insulin), then
\texttt{matchmaker \#2 to \#1}. Or search for ``ChimeraX matchmaker
superposition'' and take an example image.}%
{Two predicted insulin structures superimposed with ChimeraX's matchmaker.}%
{fig:chimerax-insulin}

\paragraph{How close is close}
Superpositions are judged by the average distance between matching residues,
in \textbf{\aa ngstr\"om} (1\,\AA{} $= 10^{-10}$\,m; a carbon--carbon bond is
about 1.5\,\AA{} long). Below 1\,\AA{}, two structures are practically
identical at that position.

\paragraph{Alternatives}
\textbf{PyMOL} (\url{pymol.org}): powerful, originates from the UNIX world
(keyboard shortcuts!), recently got a modern user interface. \textbf{TM-align}
(\url{zhanggroup.org/TM-align}): 3D alignment only, also works with low
sequence similarity, but only two proteins at a time.

\section{A Real Case: METTL3 and Bladder Cancer}
\label{sec:mettl3}

The insulin comparison is a toy example. The same workflow is used in a
current research project in which Regl is involved: a group of medical doctors
from the UK and Austria researching bladder cancer.

\paragraph{The starting point}
The doctors found two families in which bladder cancer had been inherited over
several generations (\cref{fig:pedigree}). Such a pedigree immediately suggests
a genetic predisposition. Moreover, in all (or most) cancer cases, the same
protein carried the same mutation, and the healthy family members did not.

\begin{figure}[htbp]
\centering
\begin{tikzpicture}[x=1cm, y=1cm,
    m/.style={draw, thick, minimum size=5mm, inner sep=0pt},
    f/.style={draw, thick, circle, minimum size=5.4mm, inner sep=0pt},
    aff/.style={fill=accent!80}]
  % generation I
  \node[m, aff] (g1m) at (0,0) {}; \node[f] (g1f) at (1.4,0) {};
  \draw[thick] (g1m) -- (g1f);
  \draw[thick] (0.7,0) -- (0.7,-0.6) -- (-0.7,-0.6) -- (-0.7,-0.9);
  \draw[thick] (0.7,-0.6) -- (2.1,-0.6) -- (2.1,-0.9);
  \draw[thick] (0.7,-0.6) -- (0.7,-0.9);
  % generation II
  \node[f, aff, anchor=north] (a) at (-0.7,-0.9) {};
  \node[m, anchor=north] (b) at (0.7,-0.9) {};
  \node[m, aff, anchor=north] (c) at (2.1,-0.9) {};
  \node[f, anchor=north] (cw) at (3.5,-0.9) {};
  \draw[thick] (c) -- (cw);
  \draw[thick] (2.8,-1.17) -- (2.8,-1.75) -- (2.1,-1.75) -- (2.1,-2.05);
  \draw[thick] (2.8,-1.75) -- (3.5,-1.75) -- (3.5,-2.05);
  % generation III
  \node[f, aff, anchor=north] at (2.1,-2.05) {};
  \node[m, anchor=north] at (3.5,-2.05) {};
  % legend
  \node[m, aff] at (5.6,-0.3) {}; \node[anchor=west, font=\scriptsize] at (5.9,-0.3) {affected male};
  \node[f, aff] at (5.6,-0.9) {}; \node[anchor=west, font=\scriptsize] at (5.9,-0.9) {affected female};
  \node[m] at (5.6,-1.5) {}; \node[f] at (6.4,-1.5) {}; \node[anchor=west, font=\scriptsize] at (6.7,-1.5) {unaffected};
  \node[anchor=east, font=\scriptsize, gray] at (-1.3,0) {I};
  \node[anchor=east, font=\scriptsize, gray] at (-1.3,-1.17) {II};
  \node[anchor=east, font=\scriptsize, gray] at (-1.3,-2.32) {III};
\end{tikzpicture}
\caption{How to read a pedigree (a made-up schematic, not the real family):
squares are men, circles women, filled symbols affected; a horizontal line
between two people is a couple, the lines below lead to their children. A
disease that appears generation after generation points to an inherited
cause.}
\label{fig:pedigree}
\end{figure}

\paragraph{The proteins}
METTL3 and METTL14 form a complex that attaches to RNA molecules and adds a
methyl group ($-\mathrm{CH_3}$) to one of its bases, adenine (``methylation'').
Such chemical marks on RNA influence what happens to it in the cell ---
compare the marks on DNA in epigenetics (\cref{sec:packing}). The mutation is in
METTL3: \textbf{one arginine (R) is missing} compared with the \emph{wild type}
(the normal, non-mutated sequence):

\alnblock{%
...YGTLTDDEMRRLNIP...~~(wild type)\\
...YGTLTDDEM-RLNIP...~~(mutant)}

(Since there are two R's in a row, the gap could equally be placed under the
other one; both alignments have the same score.) The question: \emph{does this
deletion influence the binding of the complex to RNA?} The doctors had no
experience with bioinformatics tools, so they asked Regl.

\paragraph{The approach}
Exactly the insulin workflow (\cref{fig:tool-workflow}), first for METTL3 only:
get the sequences of wild type and mutant (UniProtKB); test whether the PDB
already has structures; check with BLAST that the sequences really differ in
only this one position; infer structures for both with AlphaFold; check
alternatives (Rosetta and trRosetta, Boltz, ESMFold, Modeller); compare them
with ChimeraX.

\paragraph{What came out}
\begin{itemize}
  \item There is no pre-computed structure for the mutant, and AlphaFold's
    results are ``not really convincing'': within the three domains the
    prediction is stable (dark blue), outside them it is not (yellow, orange).
  \item A prediction from trRosetta (a variant of Rosetta, a predictor that has
    been around for many years) looks quite different.
  \item AlphaFold's five predictions disagree strongly in the connecting parts
    between the domains: regions that should be aligned end up in very
    different places.
  \item \emph{Within} the domains, however, the results are good: the average
    distance between matching residues of the five predictions is mostly well
    below 1\,\AA{} (\cref{tab:mettl3}).
\end{itemize}

\begin{table}[htbp]
\centering
\small
\begin{tabular}{@{}c|cccc@{}}
 & 1 & 2 & 3 & 4 \\
\midrule
0 & 0.237 & 0.372 & 1.395 & 0.211 \\
1 & --    & 0.453 & 1.178 & 0.129 \\
2 & --    & --    & 0.366 & 0.388 \\
3 & --    & --    & --    & 0.119 \\
\end{tabular}
\caption{Average distance (in \AA) between matching residues of AlphaFold's
five predictions (0--4) for the wild type, in the regions with the best
similarity (not all residues).}
\label{tab:mettl3}
\end{table}

So the connecting sequences were cut out, and the analysis concentrated on the
\textbf{large domain}: it gives the best 3D prediction, and it contains the
functional site --- and the mutation. Aligning wild type and mutant first
showed a striking dissimilarity, but it seems to be caused by the uncertainties
in the connecting parts. The main difference found so far: \textbf{the mutation
creates a short $\alpha$-helix} at the position of the missing arginine. This
seems to be a key issue.

\paragraph{Open questions and outlook}
What happens in the METTL3/METTL14 complex? Does the change influence the
binding to RNA? The next step is to model the complex together with RNA and
test how close the contacts are. If a conformational reason is found why the
mutation causes cancer, one could think of \emph{gene editing} (re-inserting
the missing arginine) or of medication that compensates the negative effects
--- a direct implication for public health.

\begin{intbox}[title={The lesson}]
Real research uses the same handful of tools as the toy example. But predicted
structures are hypotheses, not facts: always check the confidence colours,
compare several predictions and predictors, and restrict your conclusions to
the regions where they agree.
\end{intbox}

\section{Chapter Checklist}

\begin{itemize}
  \item UniProt (sequences) vs PDB (structures) vs NCBI (BLAST); AlphaFold DB vs
    AlphaFold Server;
  \item what BLAST stands for and that it produces \emph{local} alignments;
    query and subject; read the middle line (letter, +, blank, $-$);
  \item why human and bovine insulin differ mostly outside the hormone itself;
  \item AlphaFold's confidence colours; why cores are predicted well and loops
    badly; five predictions per run;
  \item ChimeraX and matchmaker (sequence alignment first, then 3D fit);
    alternatives PyMOL and TM-align;
  \item the METTL3 case: question, workflow, the main finding and its limits.
\end{itemize}

% ==========================================================
%  CHAPTERS 10--15 — OUTLINES FOR THE REMAINING LECTURES
%  To be written when the lectures have taken place.
% ==========================================================
\chapter{Pairwise Alignment I: Global Alignment}
\label{ch:global}

\lectureref{Lecture on 16.11. --- not yet held. This chapter is an outline.}

\Cref{ch:scoring} showed that there are far too many alignments to try them all
($\approx 10^{58}$ for two sequences of length 100). Global alignment finds the
best one anyway, with \emph{dynamic programming}: the best alignment of two
sequences is built from the best alignments of their prefixes, stored in a
table and filled in row by row --- the same idea as the edit distance you may
know from algorithms courses.

\planned{the Needleman--Wunsch algorithm (score matrix, recursion with match,
gap in $x$, gap in $y$; traceback); a small example filled in by hand; linear
vs affine gap costs; time and memory $O(nm)$; end gaps and the overlap
alignment from \cref{sec:more-alignment}.}

\chapter{Pairwise Alignment II: Local Alignment}
\label{ch:local}

\lectureref{Lecture on 23.11. --- not yet held. This chapter is an outline.}

Often only a part of two sequences is related --- a shared domain or a
functional site (\cref{fig:functional-blocks}). Local alignment finds the
best-scoring pair of \emph{subsequences}; it is what BLAST does
(\cref{sec:blast}).

\planned{the Smith--Waterman algorithm (the 0 in the recursion, traceback from
the maximum); why the scoring matrix must have a negative expected score;
heuristics: the idea behind BLAST (seeds, extension) and FASTA; significance
of hits (E-values).}

\chapter{Multiple Sequence Alignment}
\label{ch:msa}

\lectureref{Lectures on 30.11. and 7.12. --- not yet held. This chapter is an
outline.}

Aligning many sequences at once is what the blocks of BLOSUM
(\cref{ch:blosum}), the family trees of PAM (\cref{ch:pam}) and AlphaFold's
input (\cref{sec:alphafold-msa}) are built on.

\planned{why exact dynamic programming does not scale to many sequences;
sum-of-pairs scores; progressive alignment along a guide tree (Clustal);
profiles and position-specific scoring (\cref{sec:more-alignment}); typical
tools.}

\chapter{Phylogenetics}
\label{ch:phylo}

\lectureref{Lecture on 14.12. --- not yet held. This chapter is an outline.}

Similarity of sequences reflects relationship (\cref{ch:alignment}); phylogenetics
turns pairwise similarities into family trees, like the hemoglobin tree in
\cref{ch:alignment} or Dayhoff's trees for PAM (\cref{fig:pam-tree}).

\planned{distances from alignments; distance-based tree building (UPGMA,
neighbour joining); parsimony; rooted vs unrooted trees; the molecular clock
revisited.}

\chapter{Hidden Markov Models}
\label{ch:hmm}

\lectureref{Lecture on 21.12. --- not yet held. This chapter is an outline.}

Markov chains already appeared with PAM (\cref{sec:markov}). Hidden Markov
models add states that are not observed directly --- for example ``inside a
gene'' vs ``between genes'', or the positions of a protein family profile.

\planned{states, transition and emission probabilities; the Viterbi algorithm
(again dynamic programming); profile HMMs for protein families.}

\chapter{AI-Based Biological Sequence Models}
\label{ch:ai}

\lectureref{Lectures on 11.01. and 18.01. --- not yet held. This chapter is an
outline.}

The course ends where it began (\cref{ch:intro}): with machine learning on
the ``languages of life''.

\planned{protein language models; how AlphaFold uses multiple sequence
alignments; structure prediction and its confidence (\cref{sec:alphafold});
applications.}

\end{document}
